% concatenated from arxiv_cont_and_comp.tex
\documentclass[11pt,letterpaper]{article}
\usepackage[letterpaper, total={6.5in, 9in}]{geometry}



\usepackage{amsmath}
\usepackage{amssymb}
\usepackage{enumitem}
\usepackage{fancyhdr}
\usepackage{tikz}
\usepackage{bbm}
\usetikzlibrary{trees}% Required for inserting images
\usepackage{hyperref}


\usepackage{url}
\usepackage{graphicx}%
\usepackage{multirow}%
\usepackage{amsmath,amsfonts}%
%\usepackage{amsthm}%
\usepackage{mathrsfs}%
\usepackage[title]{appendix}%
\usepackage{xcolor}%
\usepackage{textcomp}%
\usepackage{manyfoot}%
\usepackage{booktabs}%
\usepackage{algorithm}%
%\usepackage{algorithmic}%
\usepackage{algpseudocode}%
\usepackage{listings}%
\usepackage{chngcntr}
\usepackage{dsfont}
\usepackage{comment}
%\usepackage[
%backend=bibtex
%]{biblatex}


% \usepackage[authoryear]{natbib}
% \bibliographystyle{abbrvnat} 

% \usepackage[style=numeric]{biblatex}
% \addbibresource{refs.bib} 

%\addbibresource{refs.bib}

\usepackage[skip=10pt plus1pt]{parskip}



%thanh's macros
\newcommand{\ball}{\mathrm{B}}
\newcommand{\R}{\mathbb{R}}
\newcommand{\N}{\mathbb{N}}
\newcommand{\X}{\mathcal{X}}
\newcommand{\Y}{\mathcal{Y}}
\newcommand{\W}{\mathcal{W}}
\newcommand{\Z}{\mathcal{Z}}
\newcommand{\I}{\mathcal{I}}
\newcommand{\G}{\mathcal{G}}
\newcommand{\Ro}{\mathcal{R}}
\newcommand{\eu}{{\epsilon,w^\epsilon}}
\newcommand{\on}[1]{\operatorname{#1}}
\newcommand{\norm}[1]{\left\lVert #1 \right\rVert}
\newcommand{\abs}[1]{\left\vert #1 \right\rvert}
%\newcommand{\E}[1]{\left[ #1 \right]}
\DeclareMathOperator{\ran}{ran}
\newcommand{\mb}[1]{\mathbf{#1}}
\newcommand{\bs}[1]{\boldsymbol{#1}}
\newcommand{\BC}{\operatorname{BC}}
\newcommand{\V}[1]{\mathrm{Var}{\left( #1\right)}}
\newcommand{\cov}[1]{\mathrm{Cov}{\left( #1\right)}}


\newcommand{\local}{\mathrm{loc}}


\newcommand{\K}{\widetilde{K}}
\newcommand{\Q}{\widetilde{Q}}

\newcommand{\ceil}[1]{\lceil #1 \rceil}
\newcommand{\floor}[1]{\lfloor #1 \rfloor}

\newcommand{\opnorm}[1]{\norm{#1}_{\operatorname{op}}}

\renewcommand{\P}{\mathcal{P}}
\newcommand{\E}{\mathbb{E}}
\newcommand{\Prob}{\mathbb{P}}
\newcommand{\proj}{\operatorname{proj}}
\newcommand{\prox}{\operatorname{prox}}
\newcommand{\dist}{\operatorname{dist}}
\newcommand{\epi}{\operatorname{epi}}
\newcommand{\rin}{r_{\mathrm{in}}}
\newcommand{\Argmin}{\mathrm{Argmin}  }
\newcommand{\argmin}{\mathrm{argmin}  }
\newcommand{\Argmax}{\mathrm{Argmax}  }
\newcommand{\argmax}{\mathrm{argmax}  }
\newcommand{\dom}{\mathrm{dom}  }
\newcommand{\tx}{\tilde x}
\newcommand{\ty}{\tilde y}
\newcommand{\tz}{\tilde z}
\newcommand{\lam}{\lambda}
\newcommand{\nn}{\nonumber}

\newcommand{\D}{\mathrm{d}}
\newcommand{\Ent}{\mathsf{Ent}}
\newcommand{\cW}{\mathcal{W}}
\newcommand{\KL}{\mathsf{KL}}
\newcommand{\TV}{\mathsf{TV}}
\newcommand{\Var}{\mathsf{Var}}
\newcommand{\dvert}{\mathbin{\|}}
\newcommand{\FI}{\mathsf{FI}}
\newcommand{\vol}{\mathsf{vol}}
\newcommand{\Inner}[2]{\big \langle #1  ,#2 \big \rangle}
\newcommand{\inner}[2]{\left\langle #1  ,#2 \right\rangle}

\newcommand{\brac}[1]{\left( #1\right)}

\newcommand{\red}[1]{\textcolor{red}{#1}} 
\newcommand{\blue}[1]{\textcolor{blue}{#1}} 
\newcommand*{\QEDA}{\hfill\hbox{\vrule width1.5ex height1.5ex}}


\allowdisplaybreaks



\newtheorem{thm}{Theorem}[section]
%\newtheorem{theorem}{Theorem}[section]
\newtheorem{theorem}[thm]{Theorem}
\newtheorem{lem}[thm]{Lemma}
\newtheorem{lemma}[thm]{Lemma}
\newtheorem{prop}[thm]{Proposition}
\newtheorem{proposition}[thm]{Proposition}
\newtheorem{cor}[thm]{Corollary}
\newtheorem{corollary}[thm]{Corollary}
\newtheorem{definition}[thm]{Definition}
\newtheorem{defn}[thm]{Definition}
\newtheorem{example}[thm]{Example}
\newtheorem{rem}[thm]{Remark}
\newtheorem{remark}[thm]{Remark}
\newtheorem{problem}[thm]{Problem}
\newtheorem{question}[thm]{Question}
\newtheorem{exercise}[thm]{Exercise}
\newtheorem{application}[thm]{Application}
\newtheorem{observation}[thm]{Observation}

\setcounter{page}{1}
\def\eqnok#1{(\ref{#1})}
\def\proof{\noindent{\bf Proof}: \ignorespaces}
\def\remark{\par{\it Remark}: \ignorespaces}
\def\endproof{{\ \hfill\hbox{%
      \vrule width1.0ex height1.0ex
    }\parfillskip 0pt}\par}
\def\endremark{{\ \hfill\hbox{%
      \vrule width1.0ex height1.0ex
    }\parfillskip 0pt}\par}


    % --- subsubsubsection: numbered, JMLR-like, title on its own line (no bold) ---
\makeatletter
\newcounter{subsubsubsection}[subsubsection]
\renewcommand\thesubsubsubsection{\thesubsubsection.\arabic{subsubsubsection}}

\newcommand{\subsubsubsection}[1]{%
  \par\addvspace{0.00ex \@plus 1ex \@minus .2ex}%% space BEFORE heading
  \refstepcounter{subsubsubsection}%
  \noindent{%
    \normalfont\normalsize\rmfamily\mdseries\scshape\raggedright
    \thesubsubsubsection.\enskip #1\par
  }%
  \addvspace{0.05in \@plus .2ex}%% space after heading
}
\makeatother

\setcounter{secnumdepth}{4} % needed for 3.1.1.1-style depth


\usepackage{float}
\usepackage{caption}    % provides \captionof
\usepackage{subcaption} % provides subfigure + \subref


\begin{document}
	\title{Constrained and Composite Sampling via Proximal Sampler}
	\date{}
	\author{
    	Thanh Dang  \thanks{Department of Computer Science, University of Rochester, Rochester, NY 14620. 
        (email: {\tt ycloud777@gmail.com}).}
        \qquad
	 	Jiaming Liang \thanks{
        Goergen Institute for Data Science and Artificial Intelligence (GIDS-AI) and Department of Computer Science, University of Rochester, Rochester, NY 14620 (email: {\tt jiaming.liang@rochester.edu}). This work was partially supported by GIDS-AI seed funding.}
	}
\maketitle
	
	\begin{abstract}	
We study two log-concave sampling problems: constrained sampling and composite sampling.
First, we consider sampling from a target distribution with density proportional to $\exp(-f(x))$ supported on a convex set $K \subset \R^d$, where $f$ is convex. The main challenge is enforcing feasibility without degrading mixing. Using an epigraph transformation, we reduce this task to sampling from a nearly uniform distribution over a lifted convex set in $\R^{d+1}$. We then solve the lifted problem using a proximal sampler.
Assuming only a separation oracle for $K$ and a subgradient oracle for $f$, we develop an implementation of the proximal sampler based on the cutting-plane method and rejection sampling. Unlike existing constrained samplers that rely on projection, reflection, barrier functions, or mirror maps, our approach enforces feasibility using only minimal oracle access, resulting in a practical and unbiased sampler without knowing the geometry of the constraint set.

Second, we study composite sampling, where the target is proportional to $\exp(-f(x)-h(x))$ with closed and convex $f$ and $h$. This composite structure is standard in Bayesian inference with $f$ modeling data fidelity and $h$ encoding prior information.
We reduce composite sampling via an epigraph lifting of $h$ to constrained sampling in $\R^{d+1}$, which allows direct application of the constrained sampling algorithm developed in the first part. This reduction results in a double epigraph lifting formulation in $\R^{d+2}$, on which we apply a proximal sampler. By keeping $f$ and $h$ separate, we further demonstrate how different combinations of oracle access (such as subgradient and proximal) can be leveraged to construct separation oracles for the lifted problem.
For both sampling problems, we establish mixing time bounds measured in Rényi and $\chi^2$ divergences.
	
	
	{\bf Key words.}   Constrained sampling, composite sampling, epigraph lifting, proximal sampler, mixing times. 
		\end{abstract}


\section{Introduction}
\label{sec:intro}

% Sampling is an important research area, with classical applications such as volume computation in computer science \cite{lovasz1990mixing,dyer1991random,kannan1997random} as well as modern applications including Bayesian inference, inverse problems, and neural networks \cite{murphy2023probabilistic,dashti2015bayesian,du2019implicit,chizat2019lazy}. A long line of works have explored connections between sampling and optimization, and influential developments such as \cite{jordan1998variational} helped popularize an optimization-based viewpoint for designing and analyzing sampling methods in \cite{wibisono2018sampling,Chewi2025,liu2017stein,leestructured21,hsieh2018mirrored} among other references. 


This paper studies two log-concave sampling problems: constrained sampling and composite sampling. Using an epigraph transformation, we reduce both problems to nearly uniform sampling in lifted spaces, which can then be solved via a proximal sampler.
The proximal sampler was introduced by~\cite{leestructured21} as an unbiased method for sampling from log-concave targets, and was later substantially extended by~\cite{chen2022improved} to cover distributions satisfying general functional inequalities. The algorithm is closely related to the proximal point method in optimization, hence its name. Given a step size $\eta>0$ and a target $\pi(x)\propto \exp(-f(x))$, the proximal sampler performs Gibbs sampling on the augmented distribution $\Pi^{X,Y}(x,y)\ \propto\ \exp\!\left(-f(x)-\|x-y\|^2/(2\eta)\right)$, whose $X$-marginal is exactly $\pi$. Each iteration alternates between sampling  $y \sim \Pi^{Y|X}$, a Gaussian distribution, and sampling $x\sim \Pi^{X|Y}$, known as the restricted Gaussian oracle (RGO). In general, the RGO is unavailable or difficult to implement. In this work, we provide efficient unbiased implementations of the RGO in both sampling problems under various oracle models.


Let $f : \mathbb{R}^d \to \mathbb{R}$ be a convex function and $K \subseteq \mathbb{R}^d$ be a convex set. \emph{In the first part of this paper,} we consider sampling from a distribution $\nu$ for which
\begin{equation}
\label{silif:def:nu}
\nu(x) \propto \exp(-f(x)) \mathbf{1}_K(x) = \exp \left(-f(x) - I_K(x)\right),
\end{equation}
where $I_K$ is the extended indicator function of the set $K$, taking the value $0$ if $x \in K$ and $+\infty$ otherwise; while $ \mathbf{1}_K(x)$ equals $1$ on $K$ and $0$ otherwise. To motivate the sampling approach studied in this paper, let us consider the analogy with convex optimization. The problem $\min\{f(x): x\in K\}$ can be rewritten in epigraph form $\min \{t: (x,t) \in Q\}$,
where $Q = \{(x,t) \in K \times \mathbb{R} : f(x) \le t\}$ is the \emph{epigraph} of $f$ restricted to $K$. 
This reformulation lifts the problem from $\mathbb{R}^d$ to $\mathbb{R}^{d+1}$ and indirectly handles $f$ through the constraint $f(x) \leq t$ instead of an objective function. It has been widely used in optimization, for example, the bundle methods \cite{lemarechal1975extension,lemarechal1995new,wolfe1975method}. 

Motivated by this idea from optimization, we reformulate constrained sampling \eqref{silif:def:nu} in an epigraph form. Let us define the lifted set
\begin{equation}\label{silif:def:Q}
    Q = \{(x,t) \in \mathbb{R}^d \times \mathbb{R} : x \in K,   f(x) \le a t\},
\end{equation}
for some scaling constant $a > 0$, and consider the distribution
\begin{equation}
\label{silif:def:pi}
\pi(x,t) \propto e^{-a t} \mathbf{1}_Q(x,t) = \exp(-at-I_Q(x,t)).
\end{equation}
The crucial observation is that the $X$-marginal of $\pi$ coincides with the target distribution $\nu$.
Indeed, integrating out $t$ gives $\pi^X(x)=\int \pi(x,t) \mathrm{d}t \propto \left(\int_{f(x)/a}^{\infty} e^{-at} \mathrm{d}t\right)\mathbf{1}_K(x)=\frac{1}{a}e^{-f(x)}\mathbf{1}_K(x)\propto\nu(x)$. Thus, if one can efficiently sample $(x,t)$ from $\pi$ over the lifted set $Q$, then discarding the $t$-coordinate produces samples from the target $\nu$. There are advantages to sampling from the lifted target $\pi$ in~\eqref{silif:def:pi} rather than directly from $\nu$ in~\eqref{silif:def:nu}.


    \textit{First}, aside from the constraint $Q$, the lifted $\pi$ in~\eqref{silif:def:pi} depends only on the last coordinate $t$ through an exponential factor, and is therefore a nearly uniform distribution on the convex set $Q$. This links our setting to uniform sampling on convex bodies where a broad range of algorithmic tools are available. In particular, motivated by the recent works ~\cite{dangliang2025oracle,kook2024inandout}, we will employ the proximal sampler by~\cite{leestructured21}, to sample from the nearly uniform lifted measure $\pi$.

\textit{Second}, recall that the proximal sampler relies on the RGO, which is typically difficult to implement. For the target $\pi$ in~\eqref{silif:def:pi}, however, the RGO reduces to sampling a Gaussian distribution truncated to $Q$, which can be achieved via rejection sampling combined with the cutting-plane (CP) method by \cite{jiang2020cuttingplane}. This implementation requires only a subgradient oracle for $f$ and a separation oracle for $K$. Unlike existing constrained samplers that depend on
projections, reflections, barrier functions, or mirror maps, 
% projections \cite{brosse2017sampling,Bubeck2018}, reflections \cite{du2025non,wang2025accelerating}, barrier functions \cite{kook2022sampling,kook2024gaussian}, or mirror maps \cite{hsieh2018mirrored,zhang2020wasserstein,ahn2021efficient,srinivasan2024fast}, 
our approach enforces feasibility using only minimal separation-based oracle access.

Let $f$ and $h$ be closed convex functions on $\R^d$.    \emph{In the second part of the paper,} we study the composite sampling problem of drawing samples from
\begin{equation}
\label{dolif:def:nu}
    \tilde{\nu}(x)\ \propto\ \exp \left(-f(x)-h(x)\right), \qquad x\in\R^d.
\end{equation}
Such composite structure arises naturally in Bayesian models: $f$ typically represents the negative log-likelihood, while $h$ plays the role of a negative log-prior. Although one may instead seek a maximum a posteriori estimator by solving the corresponding composite optimization problem $\min\{f(x)+h(x):x\in \R^d\}$, sampling from $\tilde \nu$ is often more informative, as it enables uncertainty quantification beyond a single point estimate. 
Moreover, the composite sampling formulation is especially valuable when $f$ and $h$ are provided through different oracle models, which can be analyzed in a unified way by the sampling algorithm developed in this work but cannot be handled as naturally in the composite optimization setting.


We now describe a two-step epigraph lifting (\emph{double lifting}) that turns sampling $\tilde{\nu}$ on $\R^d$ into sampling a nearly uniform
distribution over a convex set in $\R^{d+2}$.
We first apply the epigraph lifting to $h$:
\begin{equation}
\label{dolif:def:gammaandK}
\gamma(x,s)\ \propto\ \exp \left(-\tilde f(x,s)-I_{\K}(x,s)\right),
\,\,
\K=\{(x,s) \in \R^d \times \R:\ h(x)\le a s\}, \,\, \tilde f(x,s)=f(x)+as.
\end{equation}
Since $\gamma$ is supported on the convex set $\K$, the first lifting reduces the composite problem to constrained sampling. We then apply epigraph lifting again to $\tilde{f}$ and obtain a nearly uniform distribution in $\R^{d+2}$:
\begin{equation}
\label{dolif:def:piandQ}
\tilde{\pi}(x,s,t)\ \propto\ \exp \left(-bt-I_{\Q}(x,s,t)\right),
\qquad
\Q=\{(x,s,t):\ (x,s)\in\K,\ \tilde f(x,s)\le bt\}.
\end{equation}
Integrating out $(s,t)$ shows that $\tilde{\pi}^X=\tilde{\nu}$, since $\int_{h(x)/a}^{\infty}\int_{\tilde f(x,s)/b}^{\infty} e^{-bt} \D t \D s=\frac{1}{ab}\exp\!\left(-f(x)-h(x)\right)$; thus sampling $(x,s,t)\sim\tilde{\pi}$ and discarding $(s,t)$ yields samples from $\tilde{\nu}$. As in the single-lifting case~\eqref{silif:def:pi}, we develop a proximal sampler for the double-lifted target~\eqref{dolif:def:piandQ}, with the RGO implemented via a separation oracle for $\Q$ and the CP method. The composite viewpoint is useful here, since different oracles for $f$ and $h$ can be combined to construct a separation oracle for $\Q$.


\paragraph*{Our contributions.}
A central goal of our work is to bring ideas from optimization into sampling. To this end, we develop efficient algorithms for both  constrained sampling \eqref{silif:def:nu} and composite sampling \eqref{dolif:def:nu} via epigraph lifting and the proximal sampler. We summarize our main theoretical results as follows.

\emph{Constrained sampling.}
We consider the lifted target $\pi$ in~\eqref{silif:def:pi} with $a=d$, assuming $K\subseteq\R^d$ is closed and convex with $\ball_d\subseteq K\subseteq R\ball_d$ ($\ball_d$ as the Euclidean unit ball in $\R^d$); $f$ is closed, convex and Lipschitz-continuous with constant $L$ on $K$; and $\nu_0$ is $M$-warm with respect to $\nu$. For the proximal sampler (Algorithm~\ref{silif:alg:ASF_pi}), we implement the RGO via Algorithm~\ref{silif:RGO} by combining the CP method with a rejection sampler. Combining Theorem~\ref{silif:theo:outer} and Theorem~\ref{silif:theo:separation} yields an \textbf{iteration complexity} of $O\!\left(d^2\log(d+1) q\left(\|\cov{\nu}\|_{\mathrm{op}}+1\right)\log\!\left(2\frac{\log M}{\epsilon}\right)\right)$ to reach $\epsilon$-accuracy in R\'enyi divergence $\mathcal{R}_q$, with an analogous bound in $\chi^2$ divergence. Regarding \textbf{oracle complexity}, each iteration of Algorithm~\ref{silif:alg:ASF_pi}  makes ${\cal O}\!\left(d\log\frac{d\gamma}{\alpha}\right)$ calls to the separation oracle for $K$ and the subgradient oracle for $f$, where $\gamma=R/\mbox{minwidth}(K)$ and $1/\alpha$ is at most polynomial in $d$ with high probability. Finally, regarding \textbf{proposal complexity}, if we further assume $L={\cal O}(\sqrt{d})$, $R={\cal O}(\sqrt{d})$, and $M={\cal O}(1)$, then the average number of proposals is $\mathcal{O}(1)$ per proximal sampler iteration.



\emph{Composite sampling.} We consider the lifted target $\tilde{\pi}$ in~\eqref{dolif:def:piandQ} under scaling $a=b=d$, and assume that $f$ and $h$ are convex and Lipschitz continuous on $\R^d$ (with respective constants $L_f$ and $L_h$), and that $\tilde{\nu}_0$ is $M$-warm with respect to $\tilde{\nu}$. For the proximal sampler (Algorithm~\ref{dolif:alg:ASF_pi}), we implement the RGO via Algorithm~\ref{dolif:RGO} by combining the CP method with a rejection sampler. Section~\ref{sec:composite:mainpaper} explains how different oracle combinations for $f$ and $h$ can be used to construct a separation oracle for $\Q$, justifying the composite structure $f+h$. Combining Theorem~\ref{dolif:theo:outer} and Theorem~\ref{dolif:theo:separation} yields an \textbf{iteration complexity} of $O\!\left(d^2\log(d+2) q\left(\|\cov{\tilde{\nu}}\|_{\mathrm{op}}+1\right)\log\!\left(2\frac{\log M}{\epsilon}\right)\right)$ to reach $\epsilon$-accuracy in R\'enyi divergence $\mathcal{R}_q$, with a similar guarantee in $\chi^2$ divergence. Regarding \textbf{oracle complexity}, each iteration of Algorithm~\ref{dolif:alg:ASF_pi} makes ${\cal O}\!\left(d\log\frac{d\gamma}{\alpha}\right)$ calls to the separation oracle for $\Q$, where $\gamma$ is bounded by a constant independent of $d$ and $1/\alpha$ is at most polynomial in $d$ with high probability. Regarding \textbf{proposal complexity}, if we further assume $L_f={\cal O}(d)$, $L_h={\cal O}(d)$, and $M={\cal O}(1)$, then the average number of proposals is $\mathcal{O}(1)$ per proximal sampler iteration. 


\paragraph*{Comparison to existing works.} 
Our work fits into a growing literature on the interplay between sampling and optimization. One line of works including~\cite{durmus2019analysis,jordan1998variational,wibisono2018sampling} and others views sampling dynamics through an optimization lens in the space of probability measures. Another line of works designs sampling algorithms by importing tools from convex optimization. In this second line, constrained samplers typically enforce feasibility using projections~\cite{brosse2017sampling,Bubeck2018}, reflections~\cite{du2025non,wang2025accelerating}, barrier functions~\cite{kook2022sampling,kook2024gaussian}, or mirror maps~\cite{ahn2021efficient,hsieh2018mirrored,srinivasan2024fast,zhang2020wasserstein}. Our work falls into the second line, but the difference is that we enforce feasibility using only minimal oracle access, namely separation oracles. We refer the reader to Appendix~\ref{appen:literaturereview} for additional related works.





We highlight a few works from Appendix~\ref{appen:literaturereview}. 
\cite{kook2025algodiffusion} also studies log-concave sampling via lifting, closely related in spirit to our approach, but assumes a ground-set condition on the potential $V$. 
In contrast, we work under (A1)-(A2) for constrained sampling (Section~\ref{sec:constrained:mainpaper}) and (B1)-(B2) for composite sampling (Section~\ref{sec:composite:mainpaper}). 
These assumptions are incomparable: the ground-set condition neither implies nor is implied by ours, and it excludes some convex Lipschitz potentials covered here (see Example~\ref{ex:groundset} in Appendix~\ref{appen:generaltechnical}). Furthermore, in the composite setting, \cite{yuan2023class} apply Gibbs sampling to an augmented target that (in its simplest formulation) has density proportional to $\exp\!\left(-f(x)-h(y)-\|x-y\|_2^2/(2\eta)\right)$, assuming $f$ and $h$ are strongly convex and smooth. Meanwhile in Section~\ref{sec:composite:mainpaper}, we work with closed, convex and Lipschitz continuous $f$ and $h$. Moreover, \cite{mou2022composite} propose a Metropolis-Hastings sampler for $\pi(x)\propto e^{-f(x)-g(x)}$ using a proximal proposal
$Y\sim p(x,y)\propto \exp\!\left(-g(y)-\|y-(x-\eta\nabla f(x))\|_2^2/(4\eta)\right)$.
It requires sampling from $p$ and its normalizer; \cite{mou2022composite} note tractable cases (separable penalties, $\ell_1$/Laplace, group Lasso), but in general the normalizer is hard to estimate.


\begin{comment}
Let us mention some specific works among those reviewed in Appendix~\ref{appen:literaturereview}. \cite{kook2025algodiffusion} also studies log-concave sampling via lifting, so their work is close in spirit to this work. They assume a ground-set condition on the potential $V$, whereas we work under (A1)-(A2) in the constrained setting (Section~\ref{sec:constrained:mainpaper}) and (B1)-(B2) in the composite setting (Section~\ref{sec:composite:mainpaper}). These assumptions are incomparable: the ground-set condition neither follows from nor implies our standing assumptions, and it excludes some convex Lipschitz potentials covered here (see Example~\ref{ex:groundset} in Appendix~\ref{appen:generaltechnical}). 
Moreover, \cite{mou2022composite} propose a Metropolis-Hastings sampler for composite log-concave targets \(\pi(x)\propto e^{-f(x)-g(x)}\). Their method uses a proximal proposal: from $x$ with step size $\eta$, draw
$Y\sim p(x,y)\propto \exp\!\left(-g(y)-\|y-(x-\eta\nabla f(x))\|_2^2/(4\eta)\right)$.
Implementing this requires sampling from the proposal and evaluating its normalizer. \cite{mou2022composite} note tractable cases such as separable penalties, $\ell_1$/Laplace regularizers, and group-Lasso structure, while in general the normalizer is hard to compute.

Moreover, \cite{mou2022composite} propose a Metropolis-Hastings sampler for composite log-concave targets \(\pi(x)\propto e^{-f(x)-g(x)}\) based on a proximal proposal: from $x$ with step size $\eta$, draw
$Y\sim p(x,\cdot)\propto \exp\!\left(-g(y)-\|y-(x-\eta\nabla f(x))_2^2/(4\eta)\right)$.
This requires sampling from the proposal and evaluating its normalizer; \cite{mou2022composite} note tractable cases (separable penalties, $\ell_1$/Laplace, group Lasso), whereas outside such structure the normalizer can be difficult to compute.

%\(\pi(x)\propto e^{-f(x)-g(x)}\)
%under global functional inequality and regularity assumptions (e.g., log-Sobolev for $\pi$, smooth dissipative $f$, and convex Lipschitz continuous $g$).
\end{comment}










%In the composite sampling setting, \cite{yuan2023class} apply the Gibbs sampler to composite targets which in the simplest formulation has density proportional to $\exp\!\left(-f(x)-h(y)-\frac{1}{2\eta}\|x-y\|_2^2\right)$. 
%They assume $f$ and $h$ are strongly convex and smooth, while we assume that $f,g$ are closed, convex and Lipschitz continuous. Meanwhile,  \cite{mou2022composite} propose a Metropolis-Hastings chain with a proximal-based proposal to  sample from a composite target with density proportional to $\exp(-(f(x)+g(x)))$.
%They assume that the target satisfies a log-Sobolev inequality, $f$ is convex, has Lipschitz gradient and  satisfies a distant dissipativity assumption, and  $h$ is convex and $M_d$-Lipschitz.

\section{Notation and Definitions}
\label{sec:prelim}
Throughout the paper, we denote $\|\cdot\|$ the Euclidean norm, $\|\cdot\|_{\mathrm{op}}$ the operator norm, and $I_n$ the $n\times n$ identity matrix. Let $B_d(c)$ denote an $\ell_2$ unit ball in $\R^d$ centered at $c$. We write $[a]_+=\max\{a,0\}$. Let ${\cal O}(\cdot)$ denote the standard big-O notation.

\textbf{Absolute continuity.} For measures $\mu,\nu$ on $(E,\mathcal{F})$, we write $\mu\ll\nu$ if there exists $f:E\to\R$ such that $\mu(A)=\int_A f d\nu$ for all $A\in\mathcal{F}$. The function $f$ is the Radon-Nikodym derivative, denoted $\frac{d\mu}{d\nu}$.

\textbf{Metric.} Let $\phi:\R_{\ge0}\to\R$ be convex with $\phi(1)=0$. For probability measures $\mu\ll\nu$ on $(E,\cal F)$, define the $\phi$-divergence
$D_\phi(\mu\|\nu)=\int_E \phi \left(\frac{d\mu}{d\nu}\right) d\nu$.
This recovers, e.g., Kullback-Leibler divergence for $\phi(x)=x\log x$ and $\chi^2$ for $\phi(x)=x^2-1$.
For $q>0$, the Rényi divergence is $\mathcal{R}_q(\mu\|\nu)=\frac{1}{q-1}\log \left(\chi^q(\mu\|\nu)+1\right)$.


\textbf{Normalizing constants.} 
For a measurable $\Theta:\R^d\to\R\cup\{+\infty\}$, define the normalizing constant $N_\Theta:=\int_{\R^d} e^{-\Theta(z)}dz$ whenever this integral is finite.
%\begin{equation}
%\label{def:normalizing}
%N_\Theta:=\int_{\R^m} e^{-\Theta(z)}dz
%\end{equation}


\textbf{Poincar\'e inequality.} A distribution $\nu$ satisfies the Poincar\'e inequality (PI) with constant $C_{\mathrm{PI}}$ if $\Var_\nu(\psi)\le C_{\mathrm{PI}}\E_\nu[\|\nabla\psi\|^2]$ for all smooth bounded $\psi$. In Lemma~\ref{silif:lem:PIconstant} of Appendix~\ref{silif:appen:supportlem} and Lemma~\ref{dolif:lem:PIconstant} of Appendix~\ref{dolif:appen:supportlem}, we bound PI constants of $\pi$ in \eqref{silif:def:pi} and of $\tilde{\pi}$ in \eqref{dolif:def:piandQ}, respectively. 

\textbf{Oracles.} Assume $S\subseteq \R^d$ and $f:\R^d\to\R$ are both closed and convex . A membership oracle for $S$ decides if $x\in S$. A separation oracle for $S$ either certifies $x\in S$ or returns $g(x)$ such that $\langle g(x),x-y\rangle >0$ for all $y\in S$ (hence it subsumes membership). A projection oracle returns $\proj_S(y)=\argmin\{\norm{x-y}^2: x\in S\}$; clearly $\proj_S(y)=y$ if $y\in S$.
An evaluation oracle returns $f(x)$, and a subgradient oracle returns some subgradient of $f$. Finally, a proximal oracle for $f$ returns $\prox_{\lambda f}(x):=\argmin_{y\in \R^d}\left\{f(y)+\frac{1}{2\lambda}\|y-x\|^2\right\}$ for a given pair $(x,\lambda)$. 

\textbf{Volumes, distance to a set, and $\mbox{minwidth}$.} Let $\vol(S)$ and $\vol_{d-1}(\partial S)$ denote the volumes of $S\subseteq \R^d$ and the boundary set $\partial S \subseteq \R^{d-1}$, respectively. Let $\dist (y,S):=\inf_{z\in S}\|y-z\|$ denote the Euclidean distance from $y$ to $S$. Let $\mbox{minwidth}(S)=\min_{\norm{a}=1}\left\{\max_{y\in S}a^\top y-\min_{y\in S}a^\top y\right\}$.

\textbf{Parallel sets} For $r\ge 0$, define the parallel set $S_r := S + r \ball_d(0) = \{y\in \R^d:\dist (y,S)\le r\}$. 

\textbf{Warmness.} We say $\mu$ is $M$-warm with respect to\ $\nu$ if $\mu\ll\nu$ and $d\mu/d\nu\le M$ a.e.


\section{Constrained Sampling}
\label{sec:constrained:mainpaper}
Recall from Section \ref{sec:intro} that, starting from the constrained target $\nu(x)\propto \exp\!\left(-f(x)-I_K(x)\right)$ as in \eqref{silif:def:nu}, we perform a \emph{single lifting} and obtain the lifted measure as in \eqref{silif:def:pi}, i.e.,
\[
\pi(x,t)\ \propto\ e^{-at} \mathbf{1}_Q(x,t)\;=\;\exp\!\left(-at-I_Q(x,t)\right),
\qquad 
Q=\{(x,t)\in\R^d\times\R:\ x\in K,\ f(x)\le at\}.
\]
If one can efficiently sample $(x,t)\sim\pi$ on the lifted set $Q$, then discarding the $t$-coordinate yields samples from the desired distribution $\nu$. In this way, the original general log-concave sampling problem is reduced to sampling from a nearly uniform distribution over a convex set. We will employ the proximal sampler for the lifted target $\pi$ in this section.




For the remainder of Section~\ref{sec:constrained:mainpaper}, we will make use of the following conditions:
\begin{itemize}
\item[(A1)] $K\subset \mathbb{R}^d$ is nonempty, closed, convex, with $\ball_d(0)\subseteq K\subseteq R \ball_d(0)$;
\item[(A2)] $f$ is convex and Lipschitz continuous on $K$ with constant $L>0$;
\item[(A3)] there exists an initial distribution $\nu_0$ that is $M$-warm w.r.t.\ $\nu$;
\item[(A4)] there exists a subgradient oracle for $f$ and a separation oracle for $K$. 
\end{itemize} 

We will use the notation $\norm{f}_\infty:=\sup_{x\in K}|f(x)|<\infty$. Furthermore, denote
\begin{equation}\label{silif:def:notation}
    w=(x,t), \quad z=(y,s-a\eta), \quad C_f=\sqrt{a^2+L^2}, \quad C_K=LR+d, \quad C_L=C_f+C_K.
\end{equation}
 For a measurable set  $S\subseteq \R^{d+1}$ and $a>0$, define
\begin{equation}
\label{def:ZandB}
Z_{S}:=\int_{S} e^{-a t} \D q,
\qquad
B_{S}:=\int_{\partial S} e^{-a t} \D S,
\end{equation}
where $\D q:=\D x \D t$ is the Lebesgue measure on $\R^{d+1}$ and $dS$ is the infinitesimal surface element on $\partial S$.

\subsection{Proximal sampler for the lifted distribution $\pi$}
\label{silif:sec:asf}
Per the idea of the proximal sampler, to sample $\pi(x,t)$ in \eqref{silif:def:pi}, we consider the augmented distribution
\begin{equation}
\label{silif:def:bigpi}
    \Pi((x,t),(y,s))\propto \exp\left(-I_Q(x,t) - at - \frac{1}{2\eta}\norm{(x,t)-(y,s)}^2 \right).
\end{equation}
For fixed $(y,s)\in \R^{d+1}$, the above potential (i.e., negative log-density) function becomes
\begin{equation}\label{silif:def:Theta}
\Theta^{\eta,Q}_{y,s}(x,t)
:= I_Q(x,t)+at+\frac{\|(x,t)-(y,s)\|^2}{2\eta}
= I_Q(x,t)+\frac{\|(x,t)-(y,s-a\eta)\|^2}{2\eta}-\frac{a^2\eta}{2}+as.
\end{equation} 
Hence, $ \Pi^{Y,S|X,T}(y,s|x,t) = \mathcal{N}\brac{(x,t),\eta I_{d+1}}$ and 
$\Pi^{X,T|Y,S}(x,t|y,s) = \mathcal{N}\brac{(y,s-a\eta),\eta I_{d+1}}|_Q$, i.e., a Gaussian distribution restricted to $Q$. 
Thus, the proximal sampler for the target $\pi(x,t)$ in \eqref{silif:def:pi} is as follows.
\begin{algorithm}[H]
	\caption{Proximal sampler for the target $\pi$ in \eqref{silif:def:pi}}
	\label{silif:alg:ASF_pi}
	\begin{algorithmic}[1]
		\State  \textbf{Gaussian}: generate $(y_k,s_k)\sim  \mathcal{N}\brac{(x_k,t_k),\eta I_{d+1}}$;
		\State \textbf{RGO}: generate $(x_{k+1},t_{k+1})\sim {\cal N}\brac{(y_k,s_k-a\eta),\eta I_{d+1}}|_Q. $
	\end{algorithmic}
\end{algorithm}

Algorithm \ref{silif:alg:ASF_pi} describes the iterative steps starting from $k = 0$.
To initialize, we generate $x_0\sim \nu_0$, an $M$-warm start for $\nu$ in~\eqref{silif:def:nu}, which guarantees $x_0\in K$ almost surely by the definition of warmness (see Section ~\ref{sec:prelim}). Next, we generate $t_0\sim e^{-at}\mathbf{1}_{\{t\ge f(x_0)/a\}}$, which can be done via the inverse transformation method with CDF $F(t)=1-e^{f(x_0)-at}$ for $t\ge f(x_0)/a$. Hence, $t_0\sim \pi^{T| X=x_0}$ and $(x_0,t_0)\sim \pi_0(x,t) = \nu_0(x) \pi^{T| X=x}(t)$. 
Starting from $(x_0,t_0) \in Q$, Algorithm \ref{silif:alg:ASF_pi} keeps generating $(x_k,t_k)\sim \pi_k$ at Step 2 and maintains feasibility $(x_k,t_k) \in Q$.
Lemma~\ref{silif:lem:warmstart} in Appendix~\ref{silif:appen:supportlem} further shows that $\pi_0$ is $M$-warm for $\pi$ in~\eqref{silif:def:pi}, and retains $M$-warmness of $\pi_k$ for every $k\ge 1$.

Assuming the RGO (i.e., Step 2) is available, we obtain the iteration complexity of Algorithm \ref{silif:alg:ASF_pi} in the next theorem. The proof is  deferred to Appendix~\ref{silif:appen:missingproofs},
as it follows directly from the one-step contraction in $\chi^2$ and R\'enyi divergences given by Proposition~\ref{prop:kook} in Appendix~\ref{appen:outer} (restated from \cite[Lemma~2.9]{kook2025algodiffusion}).
Note that we also apply Lemma~\ref{silif:lem:PIconstant} on the PI constant of $\pi$, which requires $a=d$.


\begin{theorem}\label{silif:theo:outer}
Assume (A1)-(A3) holds and $a=d$ in~\eqref{silif:def:pi}. Denote by $C_{\mathrm{PI}}(\pi)$ the Poincar\'e constant of $\pi$ in~\eqref{silif:def:pi}. By Lemma~\ref{silif:lem:PIconstant}, with $\nu$ in \eqref{silif:def:nu},
$C_{\mathrm{PI}}(\pi)=  \mathcal{O}\brac{\log (d+1)\brac{\opnorm{\cov{\nu}}+1 }}$.
Given $\epsilon>0$, the iteration complexity in $\chi^2$ divergence to reach $\epsilon$-accuracy is
$\mathcal{O} \left(\frac{C_{\mathrm{PI}}(\pi)}{\eta}\log \frac{M^2}{\epsilon}\right)$.
Moreover, for $\mathcal{R}_q$ with $q\ge2$, if $M\le e^{ 1-1/q}$, then the iteration complexity in $\mathcal{R}_q$ to reach $\epsilon$-accuracy is
${\cal O} \left(\frac{C_{\mathrm{PI}}(\pi) q}{\eta}\log \left(2\frac{\log M}{\epsilon}\right)\right)$.
\end{theorem}

\subsection{RGO implementation}\label{subsec:RGO}
The complexity bounds in Theorem~\ref{silif:theo:outer} quantify how many calls to RGO are needed to reach a target accuracy, under conditions~(A1)-(A3), the choice $a=d$ and in particular the idealization that the conditional update in Step~2 is an exact RGO. 
To turn these iteration counts into end-to-end oracle complexities, we must consider how to implement this RGO. The goal of this subsection is to implement the RGO step in Algorithm~\ref{silif:alg:ASF_pi} via rejection sampling, assuming access to a subgradient oracle for $f$ and a separation oracle for $K$ (i.e., condition (A4)). We begin by explaining why these oracles are needed.

In rejection sampling, one typically draws samples from a proposal distribution that is easier to sample from than the target, while being close to the latter to keep the rejection rate low. 
In our setting, the target distribution in RGO is $\Pi^{X,T|Y,S}(x,t|y,s) = \mathcal{N} \left((y, s - a\eta),   \eta I_{d+1}\right)|_Q$, which is concentrated around
\begin{equation}
\label{silif:optimizingThetaasprojection}
\underset{(x,t) \in \R^{d+1}}{\argmin}   \Theta^{\eta,Q}_{y,s}(x,t)=\underset{w \in Q}{\argmin}  \big\|w - z\big\|^2 
= \proj_Q(z),
\end{equation}
where $w$ and $z$ are as in \eqref{silif:def:notation}.
This suggests the ideal proposal $\mathcal{N}\left(\proj_Q(z),\eta I_{d+1}\right)$ assuming we know the exact projection $\proj_Q(z)$. In our setting, we do not assume projection onto $Q \subset \R^{d+1}$, and we therefore aim to compute an approximate solution to \eqref{silif:optimizingThetaasprojection}. Via the CP method by~\cite{jiang2020cuttingplane}, more specifically Lemma~\ref{silif:lem:equivopt}(c) and Theorem~\ref{silif:theo:cuttingplane}(a), we are able to construct an approximate solution $\tilde w$ satisfying 
\begin{equation}
\label{silif:difference:tildextildet:mainpaper}
    \norm{\tilde{w}- \proj_Q(z)}\leq \brac{1+\frac{L}{a}}\sqrt{\frac{2\eta}{d+1}}. 
\end{equation}
We set $\tilde{w}$ to be the center of the proposal in our rejection sampler for the RGO, i.e., the proposal density is proportional to $\exp\brac{-\mathcal{P}_1(w)}$ where  
\begin{equation}\label{silif:def:P1}
\mathcal{P}_1(w)
=\frac{\|w-\tilde{w}\|^2+\|\tilde{w}-z\|^2}{2\eta}
-\frac{\sqrt{2}(L+d)}{a\sqrt{\eta (d+1)}}\left(\|w-\tilde{w}\|+\|\tilde{w}-z\|\right)
-\frac{6(L+d)^2}{a^2 (d+1)}+as-\frac{a^2\eta}{2}.
\end{equation}
Notably, the CP method of~\cite{jiang2020cuttingplane} needs (i) subgradients of $\|w-z\|^2$ in~\eqref{silif:optimizingThetaasprojection} (available explicitly) and (ii) a separation oracle for $Q$, which we construct below from oracle access to $f$ and $K$.
\begin{lemma}\label{silif:lem:sep} (proof in Appendix~\ref{appen:oracles})
    Evaluation and subgradient oracles of $f$, and membership and separation oracles of $K$ give a valid separation oracle for $Q$ in \eqref{silif:def:Q}.
\end{lemma}

It is worth noting that the above claim also holds when a proximal oracle for $f$ is available (see Lemma~\ref{silif:lem:sep2}). However, for simplicity of exposition, we assume access to a subgradient oracle for $f$ throughout this section. Now we state the RGO implementation of Algorithm~\ref{silif:alg:ASF_pi}. Recall the notation $w=(x,t)$ and $z=(y,s-a\eta)$. 
\begin{algorithm}[H]
	\caption{RGO implementation  of Algorithm~\ref{silif:alg:ASF_pi}} 
	\label{silif:RGO}
	\begin{algorithmic}[1]
		\State  Compute an approximate solution $\tilde{w}$ to \eqref{silif:optimizingThetaasprojection} satisfying~\eqref{silif:difference:tildextildet:mainpaper} via the CP method by \cite{jiang2020cuttingplane}. 
        \State  Generate $U\sim {\cal U}[0,1]$ and $w\sim \exp\brac{-\mathcal{P}_1(w)}$ via Algorithm \ref{silif:alg:samplingP1} in Appendix \ref{silif:appen:specialalgo}. 
		\State  If 
		\begin{equation}\label{silif:event:separation}
			U \leq  \exp\brac{-\Theta_{y,s}^{\eta,Q}(w)+ \mathcal{P}_1(w)},
		\end{equation}
		then accept $w$; otherwise, reject $w$ and go to step 2.
	\end{algorithmic}
\end{algorithm}


\textit{Remark:}
By Lemma~\ref{silif:lem:compareThetaP1} in Appendix~\ref{silif:appen:supportlem}, $\mathcal{P}_1\le \Theta_{y,s}^{\eta,Q}$ on $\R^{d+1}$, so the acceptance test~\eqref{silif:event:separation} is well-defined. Lemma~\ref{lem:rejection} in Appendix~\ref{appen:generaltechnical} then implies Algorithm~\ref{silif:RGO} is unbiased, i.e., $w\sim \Pi^{X,T\mid Y,S}$. Any accepted sample is also feasible since $w\notin Q$ gives $I_Q(w)=\infty$ and the RHS of~\eqref{silif:event:separation} equals $0$.

The following theorem about the RGO (Algorithm~\ref{silif:RGO}) is the main result of this section; its full proof is given in Appendix~\ref{silif:appen:sepatheoproof}. Since we will combine it with the proximal sampler complexity bound in Theorem~\ref{silif:theo:outer} (see the Contributions paragraph in the introduction), we also adopt the scaling $a=d$.


\begin{theorem}
\label{silif:theo:separation} 
Assume conditions (A1)-(A4) hold, $a=d$, and $\eta =1/d^2$. Then, regarding Algorithm~\ref{silif:RGO}, 
\begin{itemize}
    \item[(a)] \textbf{(oracle complexity)} there are at most ${\cal O}\left(d \log \frac{d \gamma}{\alpha}\right)$ calls to the separation oracle of $K$ and to the subgradient oracle of $f$. Here $\gamma=R/\mbox{minwidth}(K)$. Moreover, it holds that $\Pr\brac{\alpha\leq \frac{B}{(d+1)^3}}\leq 6\exp\brac{-\frac{d^2}{8}} $ for $B:=\min\left\{\frac{1}{2R(1+3R)}, \frac{1}{12R^2 \max\{L,1\} \max\{\norm{f}_\infty,1\} }\right\}$. 
    
  \item[(b)] \textbf{(proposal complexity)} the average number of proposals is no more than 
  \begin{equation}\label{silif:proofboundsepa}
      M\exp\brac{\frac{2(L+d)^2}{d^2}+\frac{16(L+d)^2}{d^2(d+1)} + \frac{3C_L(L+d)}{d^{5/2}}} \cdot\Bigg[\frac{\sqrt{2\pi} C_L}{d}
\exp \Biggl(\frac{C_L^2}{2d^2}\Biggr)+1\Bigg].
  \end{equation}
Moreover, if we assume $L={\cal O}(\sqrt{d})$, $R={\cal O}(\sqrt{d})$, and $M={\cal O}(1)$, then \eqref{silif:proofboundsepa} simplifies to ${\cal O}(1)$.
\end{itemize}
\end{theorem}

As an example, it is easy to verify that
$K=[-1,1]^d \subseteq \sqrt d\,\ball_d(0)$ and hence $R=\sqrt d$. Also for $f(x)=\|x\|_1$, we can show $\|f'(x)\|_2\le \sqrt d$ for every $x\in \R^d$, and hence $L=\sqrt d$.

% For a general log-concave measure $\nu$, one can heuristically justify $R={\cal O}(\sqrt d)$ after an isotropic rounding step. Per the end of Section~2 in \cite{lovasz2006fast}, rounding $\nu$ means applying an affine map and working with the pushforward $\tilde\nu$ in nearly isotropic position for which $\sqrt{\mathbb E_{\tilde\nu}\|X\|_2^2}={\cal O}(\sqrt d)$. Thus most of the mass of $\nu$ lies in a ball of radius ${\cal O}(\sqrt d)$, making $R={\cal O}(\sqrt d)$ a natural assumption in the rounded coordinates.


\subsubsection{Proof sketch of Theorem~\ref{silif:theo:separation}}

\textbf{Part~a:}  This part follows from Lemma~\ref{silif:lem:equivopt} and Theorem~\ref{silif:theo:cuttingplane}, the latter of which applies the CP method by~\cite{jiang2020cuttingplane} (Theorem~\ref{theo:jiang}). 
A technical issue in solving~\eqref{silif:optimizingThetaasprojection} is that CP cannot be applied directly to $\Theta^{\eta,Q}_{y,s}$, since $Q$ in \eqref{silif:def:Q} is unbounded, whereas the CP method requires an initial bounded set. 
To address this, we invoke Lemma~\ref{silif:lem:equivopt} in Appendix~\ref{silif:appen:cp} to rewrite~\eqref{silif:optimizingThetaasprojection} as an equivalent strongly convex program over the compact set $K$, namely
\begin{equation}
\label{silif:equivopt:mainpaper}
\min_{(x,t)\in \R^{d+1}}\Theta^{\eta,Q}_{y,s}(x,t)\;=\;\min_{x\in K}\zeta^{\eta}_{y,s}(x),
\qquad 
\text{with }t(x)=\max\left\{\frac{f(x)}{a}, s-a\eta\right\},
\end{equation}
where $\zeta^{\eta}_{y,s}$ is defined in~\eqref{silif:def:zeta}. 
We then run CP on $\zeta^{\eta}_{y,s}$ for an approximate minimizer and transfer the resulting guarantee back to $\Theta^{\eta,Q}_{y,s}$ via~\eqref{silif:equivopt:mainpaper}. 
The CP method requires a separation oracle for $K$ and a subgradient oracle for $\zeta^{\eta}_{y,s}$, and the latter reduces to a subgradient oracle for $f$. 
Finally, Theorem~\ref{theo:jiang} provides a relative-accuracy parameter $\alpha\in(0,1)$ in~\eqref{ineq:opt}. We choose $\alpha$ as in~\eqref{silif:def:alpha} so that the CP output $\tilde x$ is a $(d+1)^{-1}$-solution of $\min_{x\in K}\zeta^{\eta}_{y,s}(x)$, then we compute $\tilde{t}$ using the formula for $t(x)$ in~\eqref{silif:equivopt:mainpaper}. Finally, we set $\tilde{w}=(\tilde{x},\tilde{t})$, which satisfies the important relation~\eqref{silif:difference:tildextildet:mainpaper} per Lemma~\ref{silif:lem:equivopt}(c). 
With this choice of $\alpha$, Theorem~\ref{theo:jiang} yields ${\cal O}\!\left(d\log\frac{d\gamma}{\alpha}\right)$ calls to the separation oracle for $K$ and to the subgradient oracle for $f$, where $\gamma=R/\mbox{minwidth}(K)$. 

Next, we need to show that the positive parameter $\alpha$ which is chosen in~\eqref{silif:def:alpha} and appears  in the CP oracle complexity is not too small with high probability. To do so, we derive a concentration bound for $\alpha$ using the Gaussian perturbation in Step~1 of the proximal sampler (Algorithm~\ref{silif:alg:ASF_pi})  in Theorem~\ref{silif:theo:cuttingplane}. 





\textbf{Part~b:} Fix $(y,s)$. Steps~2 and~3 of Algorithm~\ref{silif:RGO} form a rejection sampler: it generates $w$ from the proposal proportional to $e^{-\mathcal{P}_1}$ and accepts with probability
$\exp \left(-\Theta^{\eta,Q}_{y,s}(w)+\mathcal{P}_1(w)\right)\le1$. 
By a standard Bayes-rule calculation (see Lemma~\ref{lem:rejection}) and recalling the definition of normalizing constants in~Section~\ref{sec:prelim}, the number of proposals is geometric with mean
\begin{equation}
\label{silif:formulanys:mainpaper}
    n_{y,s}=\frac{N_{\mathcal{P}_1}}{N_{\Theta^{\eta,Q}_{y,s}}}, \quad \text{conditioned on $(y,s)$.}
\end{equation} 
Let $\mu_k$ be the law of $(y,s)$. Then the average number of proposals
over the random input $(Y,S)\sim\mu_k$ is $\E_{\mu_k}[n_{y,s}]$.
By Lemma~\ref{silif:lem:warmstart}, $\frac{d\mu_k}{d\Pi^{Y,S}}\le M$, so $\E_{\mu_k}[n_{y,s}]\le M  \E_{\Pi^{Y,S}}[n_{y,s}]$, 
and it suffices to bound $\E_{\Pi^{Y,S}} [n_{y,s}]$ which satisfies
\[
    \E_{\Pi^{Y,S}} [n_{y,s}]\stackrel{\eqref{silif:formulanys:mainpaper}}{=}\E_{\Pi^{Y,S}} \left[\frac{N_{\mathcal{P}_1}}{N_{\Theta^{\eta,Q}_{y,s}}}\right]\stackrel{\text{Lemma~\ref{silif:lem:compareP1P2}}}{\leq} \E_{\Pi^{Y,S}} \left[\frac{N_{\mathcal{P}_2}}{N_{\Theta^{\eta,Q}_{y,s}}}\right], 
\]
where $\mathcal{P}_2$ is the auxiliary function defined in \eqref{silif:def:P2}. The usefulness of the auxiliary  $\mathcal{P}_2$ is that $N_{\mathcal{P}_2}$ as a Gaussian-integral-like expression is easy to bound via Lemma~\ref{lem:gaussianint}(b). 

At this point, bounding $\E_{\Pi^{Y,S}}$ reduces to controlling the integration
\[
    \int_{\R^{d+1}} e^{-d z_{d+1}}  
\exp \left(-\frac{(\dist(z,Q)-\tau)^2}{2\eta}\right)  \D z
\]
with $\tau:=2\left(1+\frac{L}{d}\right)\sqrt{\frac{2\eta}{d+1}}$ and $z_{d+1}=(y,s-d\eta)_{d+1}=s-d\eta$. The \emph{crucial} tool is Lemma~\ref{silif:lem:keytechnical} (stated after this proof) with the choice $a=d$ to bound the above integral. The rest of the proof follows by applying $\E_{\mu_k}[n_{y,s}]\le M  \E_{\Pi^{Y,S}}[n_{y,s}]$, setting $\eta=1/d^2$. 

Finally, recall that $C_L=\sqrt{d^2+L^2}+LR+d$ from~\eqref{silif:def:notation}, then assuming $L={\cal O}(\sqrt{d})$, $R={\cal O}(\sqrt{d})$, and $M={\cal O}(1)$, one can verify that $C_L={\cal O}(d)$ and \eqref{silif:proofboundsepa} becomes ${\cal O}(1)$. Therefore, the proof is completed.
\QEDA

%%%%%%%%%%%%%%%%%%%%%%


\underline{\emph{Key lemmas:}} Next, we present the key lemmas used in the above proof sketch for Part~b. For readability, we again only sketch their proofs below and present the full proofs in Appendix~\ref{silif:appen:missingproofs}. 

\begin{lemma}
\label{silif:lem:keytechnical}
Assume conditions (A1) and (A2) hold. For any $\tau>0$ and $w=(x,t)\in \R^{d+1}$,  we have %Assume (A1) and (A2)
\begin{equation}\label{ineq:lem5}
    \int_{\R^{d+1}} e^{-a t}  
\exp \left(-\frac{(\dist(w,Q)-\tau)^2}{2\eta}\right)  \D w
\ \le\
Z_Q \exp(\tau C_L) \left[
 C_L\sqrt{2\pi\eta}\exp \left(\frac{\eta C_L^2}{2}\right)+1\right].
\end{equation}
\end{lemma}

\subsubsubsection{Proof sketch of Lemma~\ref{silif:lem:keytechnical}}
The integrand depends on $w$ only through $\dist(w,Q)$, so we split the domain into $Q$ and its complement $Q^c$.  Concretely, $\int_{\R^{d+1}} e^{-a t}  
\exp \left(-\frac{(\dist(w,Q)-\tau)^2}{2\eta}\right)  \D w= \int_{Q} \ldots  \D w+\int_{Q^C} \ldots  \D w. $ For  $w \in Q$, $\dist(w,Q)=0$ so the exponential term becomes a constant $\exp \left(-\tau^2/(2\eta)\right)$, giving the contribution $Z_Q  \exp \left(-\tau^2/(2\eta)\right)$ in view of \eqref{def:ZandB}. For $w \in Q^c$, we group points by their distance $r=\dist(w,Q)$: the level set $\{w: \dist(w,Q)=r\}$ is just $\partial Q_r$. So via the co-area formula and \eqref{def:ZandB}, we have 
\[
    \int_{Q^c} e^{-at}\exp \left(-\frac{(\dist(w,Q)-\tau)^2}{2\eta}\right)  \D w
\stackrel{\eqref{def:ZandB}}=\int_0^\infty \exp \left(-\frac{(r-\tau)^2}{2\eta}\right)  B_{Q_r}  \D r. 
\]
The difficulty is how to bound the integral on the right-hand side, which is  handled by Lemma~\ref{silif:lem:envelope} below. Adding the $Q$ and $Q^c$ contributions yields the final estimate.  \QEDA


\begin{lemma}
\label{silif:lem:envelope}
Assume conditions (A1) and (A2) hold. For every $\tau\geq 0$, we have %Assume (A1) and (A2)
\[
\int_{0}^{\infty} \exp \left(-\frac{(r-\tau)^2}{2\eta}\right)  
B_{Q_r}  \D r
\ \le\
 Z_Q C_L   \sqrt{2\pi\eta}\;
\exp \left(\tau C_L + \frac{\eta C_L^2}{2}\right)+Z_Q\left(e^{C_L\tau}-1\right).
\]
\end{lemma}


\textbf{Proof sketch of Lemma~\ref{silif:lem:envelope}} (full proof in Appendix~\ref{silif:appen:missingproofs}):
Bounding $\int_0^\infty e^{-(r-\tau)^2/(2\eta)} B_{Q_r} \D r$ directly is difficult since
$B_{Q_r}=\int_{\partial Q_r} e^{-at} \D S$ as a boundary integral is sensitive to the geometry of the boundary $\partial Q_r$.
In contrast, $Z_{Q_r}=\int_{Q_r} e^{-at} \D x \D t$ is a volume integral that can be bounded by enclosing $Q_r$ in an
envelope set. Hence, we will eliminate $B_{Q_r}$ in favor of $Z_{Q_r}$ via integration by parts (IBP) and
$\frac{d}{\D r}Z_{Q_r}=B_{Q_r}$ a.e.



\underline{\emph{Step 1 (envelope set).}}
We bound $Z_{Q_r}$ by replacing $Q_r$ with a slightly larger set whose
integral is easy. Recall $Q=\{(x,t):x\in K,\ f(x)\le a t\}$ and
$Q_r=Q+r\ball_{d+1}(0)$.  Define the envelope $\mathcal E_r=
\left\{(x,t): x\in K_r,\ f(x)\le a t + r C_f\right\}$. The point is that moving a distance at most $r$ from a feasible point $(x,t)\in Q$
can only (i) move $x$ into the enlarged set $K_r=K+r\ball_{d}(0)$, and (ii) violate the inequality
$f(x)\le at$ by at most a constant times $r$ (per the Lipschitz property of $f$). This gives the inclusion $Q_r\subseteq \mathcal E_r$ and $Z_{Q_r}\le Z_{\mathcal E_r}$. Then note that $\mathcal E_r$ is designed so $Z_{\mathcal E_r}$ is easy to compute exactly, and we arrive at 
$Z_{Q_r}\le Z_{\mathcal E_r}\le e^{rC_L}Z_Q$, which relates $Z_{Q_r}$ back to $Z_Q$, a quantity of our interest.


\underline{\emph{Step 2 (IBP).}}
Lemma~\ref{lem:coarea} reveals a key observation $\frac{d}{\D r}Z_{Q_r}=B_{Q_r}$ a.e. Thus, using IBP, we can rewrite $\int_0^\infty e^{-(r-\tau)^2/(2\eta)} B_{Q_r} \D r$ into an expression in terms of $Z_{Q_r}$ instead of $B_{Q_r}$.
This is advantageous because Step~1 controls $Z_{Q_r}$ by
$e^{rC_L}Z_Q$, while the Gaussian weight $w_{\eta,\tau}(r)=e^{-(r-\tau)^2/(2\eta)}$ decays fast enough that
the boundary term at $r=\infty$ vanishes. The problem then reduces to bounding a
one-dimensional integral of the form
$\int_0^\infty e^{rC_L}w_{\eta,\tau}(r) \D r$. This last part is simple and yields the desired final estimate.  \QEDA



\section{Composite Sampling}
\label{sec:composite:mainpaper}

This section follows the general structure of Section~\ref{sec:constrained:mainpaper}, so we provide only a sketch of the proximal sampler development targeting $\tilde{\pi}$ in~\eqref{dolif:def:piandQ} and of the associated RGO implementation, deferring full details to Appendix~\ref{dolif:appen:fulldetails}. We will also state our main results on iteration complexity and oracle complexity. Our goal is to highlight the key differences in our treatments of composite sampling and constrained sampling.



With respect to the quantities at~\eqref{dolif:def:nu} and~\eqref{dolif:def:piandQ}, we will use the following conditions:
\begin{itemize}
\item[(B1)] $f$ and $h$ are closed and convex functions on $\R^d$;

\item[(B2)] $f$ and $h$ are Lipschitz continuous with constants $L_f$ and $L_h$, respectively;

    \item[(B3)] there exists an initial distribution $\tilde{\nu}_0$ that is $M$-warm with respect to $\tilde{\nu}$ i.e., $\frac{d\tilde{\nu}_0}{d \tilde{\nu}}\leq M$;
    
\item[(B4)] there exists a separation oracle for $\Q$.
\end{itemize}

We will continue to use the notations $Z_S$ and $B_S$ in~\eqref{def:ZandB}. Also, denote $\delta = \sqrt{\frac{2\eta}{d+2}}$ and 
\begin{equation}
\label{dolif:def:notation}
    p=(x,s,t), \, q=(y,u,v-b\eta);\quad \widetilde{C}_f=\sqrt{b^2+a^2+L_f^2}, \, \widetilde{C}_h=\sqrt{a^2+L_h^2}, \, \widetilde{C}_L=\widetilde{C}_f+\widetilde{C}_h. 
\end{equation}
\textbf{Double lifting and oracle construction:} As discussed in the introduction, for the composite target $\tilde{\nu}(x)\propto \exp\left(-f(x)-h(x)\right)$ in $\R^d$, we apply a \emph{double lifting} that keeps $f$ and $h$ separate, lifting one function at a time. In particular, we first lift $\tilde{\nu}$ to $\gamma$ supported on $\K\subseteq \R^{d+1}$ as in~\eqref{dolif:def:gammaandK}, which we observe is a constrained log-concave measure in $\R^{d+1}$ and thus links the current composite problem to the constrained problem of Section~\ref{sec:constrained:mainpaper}. Next, we lift $\gamma$ to the nearly uniform  measure $\tilde{\pi}$ supported on $\Q\subseteq \R^{d+2}$ as in~\eqref{dolif:def:piandQ}. Then, sampling $(x,s,t)\sim\tilde{\pi}$ and discarding $(s,t)$ yields samples from $\tilde{\nu}$. Thus, this \emph{double lifting} reduces sampling from the composite target $\tilde{\nu}$ to sampling a nearly uniform target $\tilde{\pi}$, and results of Section \ref{sec:constrained:mainpaper} would apply in principle.

Similar to the discussion in Subsection \ref{subsec:RGO}, a separation oracle for $\Q$ is needed for applying the CP method to solve an optimization problem on $\Q$. Assuming different oracle access to $f$ and $h$, we provide the implementations of the separation oracle for $\Q$. First, if both subgradient oracles for $f$ and $h$ are available, then we need not apply the lifting technique and can simply treat $f+h$ as one potential. Hence, any log-concave sampling algorithm using the subgradient of $f+h$ (e.g., Langevin Monte Carlo) applies to this composite sampling problem. Next, we discuss two concrete implementations of separation oracles for $\Q$.
For simplicity, we denote $\prox_f$ and $f'$ the proximal and subgradient oracles for a function $f$, respectively. Moreover, we assume the evaluation oracles for $f$ and $h$ are available.

\textit{Case 1: $\prox_f$ and $\prox_h$ are available.} Indeed, $\prox_h$ gives a projection oracle for $\epi(h)$ by Lemma~\ref{lem:proj-epi}, and hence a separation oracle for $\epi(h)$ by Lemma~\ref{lem:proj-sep}. By a simple scaling, this yields a separation oracle for $\K$ in \eqref{dolif:def:gammaandK}. In addition, Lemma \ref{silif:lem:proxfas} shows that $\prox_f$ yields a proximal oracle for $\tilde f(x,s)$ in \eqref{dolif:def:gammaandK}.
     Now, applying Lemma \ref{silif:lem:sep2} to $\K$ and $\tilde f$, we have a separation oracle for $\Q$ in \eqref{dolif:def:piandQ}.

\textit{Case 2: $f'$ and $\prox_h$ are available.} The same argument as in the above case gives a separation oracle for $\K$. 
    In addition, it is straightforward to see that $f'$ yields a subgradient oracle for $\tilde f(x,s)$ in \eqref{dolif:def:gammaandK}.
    Now, applying Lemma \ref{silif:lem:sep} to $\K$ and $\tilde f$, we have a separation oracle for $\Q$ in \eqref{dolif:def:piandQ}.

\textbf{The proximal sampler:} (Algorithm~\ref{dolif:alg:ASF_pi} in Appendix~\ref{dolif:appen:fulldetails}) Similar to the augmented distribution $\Pi$ in \eqref{silif:def:bigpi}, we also have in the composite setting the augmented distribution
\begin{equation}
\label{dolif:def:Pi}
    \widetilde{\Pi}\brac{(x,s,t),(y,u,v)}\propto \exp\brac{ -\widetilde{\Theta}^{\eta,\Q}_{y,u,v}(x,s,t)}
\end{equation}
where
\begin{equation}
\label{dolif:def:Theta}
    \widetilde{\Theta}^{\eta,\Q}_{y,u,v}(x,s,t)=I_{\Q}(x,s,t)+\frac{1}{2\eta}\|(x,s,t)-(y,u,v-\eta b)\|^2
+bv-\frac{\eta b^2}{2}. 
\end{equation}
The proximal sampler for $\widetilde{\Pi}$ alternates between a Gaussian step $\widetilde{\Pi}^{Y,U,V|X,S,T}(\cdot|\,x,s,t)=\mathcal{N}(p,\eta I_{d+2})$ and an RGO step $ \widetilde{\Pi}^{X,S,T|Y,U,V}(\cdot|y,u,v)=\mathcal{N}(q,\eta I_{d+2})\big|_{\Q}$. Assuming an exact RGO, Proposition~\ref{prop:kook} gives the iteration complexity of Algorithm~\ref{dolif:alg:ASF_pi}. The proof parallels Theorem~\ref{silif:theo:outer} and hence is omitted. 

\begin{theorem}\label{dolif:theo:outer}
Assume (B1) and (B3) hold, and $a=b=d$ in~\eqref{dolif:def:piandQ}. Let $\epsilon>0$ and denote by $C_{\mathrm{PI}}(\tilde{\pi})$ the Poincar\'e constant of $\tilde{\pi}$ at~\eqref{dolif:def:piandQ}. By Lemma~\ref{dolif:lem:PIconstant}, with $\tilde{\nu}$ at~\eqref{dolif:def:nu}, we know that
$C_{\mathrm{PI}}(\tilde{\pi})=\mathcal{O} \left(\log(d+2)\left(\|\mathrm{Cov}(\tilde{\nu})\|_{\mathrm{op}}+1\right)\right)$.
Then the iteration complexity in $\chi^2$ divergence to reach $\epsilon$-accuracy is
$\mathcal{O} \left(\frac{C_{\mathrm{PI}}(\tilde{\pi})}{2\eta}\log \frac{M^2}{\epsilon}\right)$.
Moreover, for $\mathcal{R}_q$ with $q\ge2$, if $M\le e^{1-1/q}$, then the iteration complexity in $\mathcal{R}_q$ to reach $\epsilon$-accuracy is
$O \left(\frac{C_{\mathrm{PI}}(\tilde{\pi}) q}{\eta}\log \left(2\frac{\log M}{\epsilon}\right)\right)$.
\end{theorem}


\textbf{RGO implementation:} (Algorithm~\ref{dolif:RGO} in Appendix~\ref{dolif:appen:fulldetails}) Similar to the idea in Section \ref{sec:constrained:mainpaper}, to implement the RGO, we employ a rejection sampler for the truncated Gaussian $\mathcal{N}\brac{q,\eta I_{d+2}}|_{\Q}$. To find the center of the proposal in our rejection sampler, we apply the CP method by~\cite{jiang2020cuttingplane} with the separation oracle on $\Q$ (condition~(B4)) to approximately solve for 
\begin{equation}
\label{dolif:prob:minimizetheta}
    \argmin\left\{\widetilde{\Theta}^{\eta,\Q}_{y,u,v}(x,s,t): (x,s,t)\in\R^{d+2}\right\} =\argmin \left\{\|p-q\|^2: p\in\Q\right\} =\proj_{\Q}(q),
\end{equation} 
for $p, q$ in \eqref{dolif:def:notation}. Since the CP method begins with a bounded set initially, we will restrict the search to a local bounded region inside a ball that still contains the unique minimizer $\proj_{\Q}(q)$ in \eqref{dolif:prob:minimizetheta}. In particular, the reformulation in Lemma~\ref{silif:lem:equivopt} for a bounded convex body $K$ in Section~\ref{sec:constrained:mainpaper} does not extend to the composite setting since $f$ and $h$ are defined on the unbounded domain $\R^d$.


Let $(y,u,v)\in\R^{d+2}$ be the current input to Algorithm~\ref{dolif:RGO} (i.e., RGO implementation) and set $q=(y,u,v-b\eta)$. Let $p=(x,s,t)\in\Q$ a.s. be the previous RGO output and define $r_{\local}:=\|p-q\|$. Define the local ball and local region by
\begin{equation}
\label{dolif:def:localball}
    \ball_{\local}:=2r_{\local}\ball_{d+2}(q), \qquad  \Q_{\local}:=\Q\cap \ball_{\local}.
\end{equation}
To ensure that restricting the CP method to $\Q_{\local}$  is valid, we verify in Lemma~\ref{dolif:lem:minimizerinlocalball} in Appendix~\ref{dolif:appen:supportlem} that the minimizer $\proj_{\Q}(q)$
belongs to $\Q_{\local}$ a.s. Moreover,  the CP method requires a subgradient oracle for the objective $\|p-q\|^2$ in \eqref{dolif:prob:minimizetheta}, which is available explicitly; and a separation oracle for $\Q_{\local}$. The latter is obtained by combining the separation oracle for $\Q$ (which can be constructed from oracles for $f$ and $h$ per the earlier explanation) and a simple separation oracle for the local ball $2r_{\local}\ball_{d+2}(q)$. 



%Having verified the requirements, we are able to apply the CP method by \cite{jiang2020cuttingplane} to  obtain a $(d+2)^{-1}$-solution $\tilde{ p}$ to $\min_{ p\in \Q_\local} \|p-q\|^2$ and establish the complexity in Theorem~\ref{dolif:theo:cuttingplane}. This enables us to center the proposal at $\tilde{p}$ in our rejection sampler for the RGO, i.e., the proposal density is proportional to $\exp\brac{-\widetilde{\mathcal{P}}_1(p)}$ where 

Having verified the requirements, the CP method of \cite{jiang2020cuttingplane} applies and returns a $(d+2)^{-1}$-solution $\tilde p$ to $\min_{p\in\Q_{\local}}\|p-q\|^2$, with the complexity stated in Theorem~\ref{dolif:theo:cuttingplane}. We then center the RGO rejection proposal at $\tilde p$, i.e., the proposal density is proportional to $\exp\!\left(-\widetilde{\mathcal{P}}_1(p)\right)$, where
\begin{equation}
\label{dolif:def:tildeP1}
\widetilde{\mathcal{P}}_1(p)
= \frac{\|p-\tilde p\|^2 + \|\tilde p-q\|^2}{2\eta}  -\frac{\delta}{\eta} (\|p-\tilde p\|+\|\tilde p-q\|) + bv-\frac{b^2\eta}{2}-\frac{\delta^2}{\eta}.
\end{equation}
In summary, the RGO implementation (i.e., Algorithm~\ref{dolif:RGO}) consists of two steps: first, solving for $\tilde{p}$ via the CP method to define the potential $\widetilde{\mathcal{P}}_1$; and second building a rejection sampler with the proposal $\exp\brac{-\widetilde{\mathcal{P}}_1(p)}$, which can be implemented by Algorithm \ref{dolif:alg:samplingP1} in Appendix \ref{dolif:appen:specialalgo}.


The following is our main result for the RGO implementation; see its proof in Appendix~\ref{dolif:appen:proofsepatheorem}. Since this result will be paired with the proximal sampler iteration complexity in Theorem~\ref{dolif:theo:outer} (see the Contributions paragraph in the introduction), we specialize to the scaling $a=b=d$.

\begin{theorem}
\label{dolif:theo:separation}
Assume (B1)-(B4) holds, $a=b=d$, and $\eta =1/d^2$. Then regarding Algorithm~\ref{dolif:RGO}, 
\begin{itemize}
\item[(a)] \textbf{(oracle complexity)} there are at most ${\cal O} \left((d+2) \log \frac{(d+2) \gamma}{\alpha}\right)$ calls to the separation oracle of
$\Q$. Here $\gamma=\frac{2r_{\local}}{\mbox{minwidth}(\Q_{\local})}$
% and $\alpha\in(0,1/2]$ are from Theorem~\ref{dolif:theo:cuttingplane}. Moreover, we have  $\gamma\ \le\ \frac{2(3+\mu)}{\mu}$ for $\mu=\min\left\{\frac{d}{\sqrt{d^2+L_h^2}},\ \frac{d}{\sqrt{2d^2+L_f^2}}\right\}$; and $\Pr \left(
% \alpha\ \le\
% \min\left\{
% \frac{1}{2(d+2) (\sqrt{d+2}+d/2+1)^2},\ \frac{1}{2}
% \right\}
% \right)
% \le 2\exp \left(-\frac{d^2}{8}\right)$. 
\begin{align*}
&\gamma\ \le\ \frac{2(3+\mu)}{\mu}
\quad\text{for}\quad
\mu=\min\left\{\frac{d}{\sqrt{d^2+L_h^2}},\ \frac{d}{\sqrt{2d^2+L_f^2}}\right\};\\
&\Pr \left(
\alpha\ \le\
\min\left\{
\frac{1}{2(d+2) (\sqrt{d+2}+d/2+1)^2},\ \frac{1}{2}
\right\}
\right)
\le 2\exp \left(-\frac{d^2}{8}\right).
\end{align*}
Thus, ${\cal O} \left((d+2) \log \frac{(d+2) \gamma}{\alpha}\right)={\cal O} \left(d\log d \right)$ in high probability. 
\item[(b)] \textbf{(proposal complexity)} the average number of proposals is at most 
\begin{equation}\label{dolif:middlestepsepa}
    M \exp\!\left(\frac74+\frac{16}{d}+ \frac{3\widetilde{C}_L}{d^{3/2}}\right)
\left[
\frac{\sqrt{2\pi} \widetilde C_L}{d} \exp\!\left(\frac{ \widetilde C_L^2}{2d^2}\right)
+1 \right].
\end{equation}
Moreover, if we assume $L_f={\cal O}(d)$, $L_h={\cal O}(d)$, and $M={\cal O}(1)$, then \eqref{dolif:middlestepsepa} becomes $\mathcal{O}(1)$. 
\end{itemize}
\end{theorem}


%While the proof of the proposal complexity closely parallels the corresponding argument in Theorem~\ref{silif:theo:separation}, the proof of the oracle complexity is more involved. In particular, bounding the number of separation oracle calls requires a lower bound on $\mbox{\em minwidth}(\Q_{\local})$, which we obtain by constructing an inscribed ball $\ball_{\mathrm{in}}$ of radius $\rin$ inside $\Q_{\local}$ and using $\mbox{\em minwidth}(\Q_{\local})\ge \mbox{\em minwidth}(\ball_{\mathrm{in}})=2\rin$. The construction of the inscribed ball $\ball_{\mathrm{in}}$ in Lemma~\ref{dolif:lem:localset} is delicate and relies on a mild enlargement of the local set, namely $\Q\cap cr_{\local}\ball_{d+2}(q)$ with $c>1$ (the argument breaks for $c=1$), which is why we work with $\Q_{\local}=\Q\cap 2r_{\local}\ball_{d+2}(q)$.

The proposal complexity proof follows the argument of Theorem~\ref{silif:theo:separation}, but the oracle-complexity proof is subtler. In particular, to bound the number of separation-oracle calls we need a lower bound on $\mbox{\em minwidth}(\Q_{\local})$. We do this by inscribing a ball $\ball_{\mathrm{in}}\subseteq \Q_{\local}$ of radius $\rin$, giving $\mbox{\em minwidth}(\Q_{\local})\ge \mbox{\em minwidth}(\ball_{\mathrm{in}})=2\rin$. Constructing $\ball_{\mathrm{in}}$ in Lemma~\ref{dolif:lem:localset} is delicate and requires a slight enlargement of the local set, i.e., working with $\Q\cap c r_{\local}\ball_{d+2}(q)$ for some $c>1$ (the argument fails at $c=1$); this is why we take $\Q_{\local}=\Q\cap 2r_{\local}\ball_{d+2}(q)$.


%%%%%%%%%%%%%%%%%%%%%%%%%%%%%%%%%%%%%%%%%%%%%%%%%%%
%%%%%%%%%%%%%%%%%%%%%%%%%%%%%%%%%%%%%%%%%%%%%%%%%%%
%%%%%%%%%%%%%%%%%%%%%%%%%%%%%%%%%%%%%%%%%%%%%%%%%%%


\section{Conclusions}

This paper develops sampling algorithms for log-concave constrained and composite sampling via epigraph lifting and the proximal sampler.
Our algorithms are based on a sequence of liftings that connect these general settings to nearly uniform sampling in lifted spaces. Central to our approach is an efficient and implementable RGO used within proximal sampler.
Our RGO implementation is based on rejection sampling and the CP method. It requires only separation oracles for the lifted constraint sets $\K$ and $\Q$, and we provide concrete constructions of these oracles under various oracle models for accessing $f$, $h$, and $K$.
For both problems, we establish $\mathcal{O}\brac{d^2\log d \log \epsilon^{-1}}$ mixing time guarantees measured in Rényi and $\chi^2$-squared divergences.

A natural future direction is to relax the global Lipschitz continuity assumption in both the constrained and composite settings. In our analysis, global Lipschitz continuity provides uniform control that is crucial for the envelope arguments in Lemmas~\ref{silif:lem:envelope} and~\ref{dolif:lem:envelope}. It would be interesting to replace this requirement with weaker conditions, such as local Lipschitz continuity. For instance, \cite{kook2025algodiffusion} adopts a ground-set assumption in place of global Lipschitz continuity.
% Another direction is to improve the step size $\eta=\mathcal{O}\brac{d^{-2}}$, which currently facilitates efficient RGO implementations. Potential acceleration may be possible via biased RGOs based on approximate rejection sampling, as in~\cite{fan2023improved}. However, extending their approach to our setting is nontrivial, since our lifted RGOs involve truncated Gaussian sampling, whereas their concentration analysis relies on untruncated Gaussians.



% A future research direction is to relax the global Lipschitz continuity assumptions in both the constrained and composite settings. In our analysis, Lipschitz continuity provides uniform global control that is essential for the envelope arguments in Lemma~\ref{silif:lem:envelope} and Lemma~\ref{dolif:lem:envelope}. It would be interesting to replace this global condition with weaker alternatives such as local Lipschitz continuity. For example, \cite{kook2025algodiffusion} assumes a ground-set condition in lieu of global Lipschitzness  (see Example~\ref{ex:groundset} in Appendix~\ref{appen:generaltechnical}). In addition, this paper and~\cite{kook2025algodiffusion} work with the step size $\eta=\mathcal{O}\brac{d^{-2}}$, which enables efficient RGO implementations of the proximal sampler. A natural direction for improving the step size and potentially accelerating the proximal sampler is to use a biased RGO as in~\cite{fan2023improved}, where the RGO is realized via approximate rejection sampling. Adapting~\cite{fan2023improved} to our setting appears non-trivial in both the constrained and composite cases for the same reason: after lifting, the RGOs in our proximal samplers sample truncated Gaussian distribution, whereas the central concentration-inequality argument in~\cite{fan2023improved} is developed for untruncated Gaussian distribution. 

% \red{check citep or citet} \blue{done}


% \red{check to see if we refer to the correct appendices for all results.} \blue{correct hyperref references for the main paper. Appendices not yet checked.}


% \red{check if we refer to the correct condition A1, A2, B1, etc}







\bibliographystyle{plain}
\bibliography{refs}


%%%%%%%%%%%%%%%%%%%%%%%%%%%%%%%%%%%%%%%%%%%%%%%%%%%%%%%%%%%%%%%%%%%%%%%%%%%%%%%%%%%%%%%%%%%%%%%%%%%%%%%%%%%%%%%%%%%%%%%%%%%%%%%%%%%%%%%%%%%%%%%%%%%%%%%%%%%%%%%%%%%%%%%%%%%%%%%%%%%%%%%%%%%%%%%%%%%%%%%%%%%%%%%%%%%%%%%%%%%%%%%%%%%%%%%%%%%%%%%%%%%%


\newpage
\appendix


\setcounter{tocdepth}{2}
\tableofcontents

\section{Additional Related Works}
\label{appen:literaturereview}



In the context of constrained sampling, uniform sampling on a convex set is a fundamental case as it is closely tied to efficient volume computation. Notable works on uniform sampling over convex sets include~\cite{kook2024inandout,diaconis2012gibbs,narayanan2022mixing,laddha2023convergence,lovasz1999hit,cheneldanhitandrun,chen2018fast,kannan2009random}, among others.



Next, we focus on works about sampling general constrained targets supported on a set $K$ and, when appropriate, classify them by the operations they employ on $K$. Works such as~\cite{brosse2017sampling,lamperski2021projected,durmus2018efficient,Bubeck2018} study algorithms similar to Langevin Monte Carlo (LMC) that rely on the Euclidean projection onto $K$. Meanwhile, \cite{du2025non} (see also~\cite{wang2025accelerating}) investigates LMC-type algorithms based on a skew-projection operator built from the standard Euclidean projection; and moreover, the algorithm in~\cite{du2025non} contains that of~\cite{Bubeck2018} as a special case. \cite{ahn2021efficient,srinivasan2024fast,hsieh2018mirrored,zhang2020wasserstein,li2022mirror,jiang2021mirror} incorporate a mirror map that is self-concordant on $K$ into the LMC framework. \cite{kook2024gaussian,cousins2018gaussian,kook2022sampling} combine barrier methods with Gaussian cooling in the first two papers, and with Riemannian HMC in the last one. \cite{kook2025algodiffusion} performs a epigraph lifting and then applies the proximal sampler as in our paper. Their proximal step is implemented via rejection sampling and a membership oracle, with a failure event triggered if the number of trials exceeds a prescribed cutoff. This is followed by \cite{kook2025zeroth}, which eliminates the failure event.




Finally, we turn to works on composite sampling. From the proximal perspective, \cite{mou2022composite} proposes a Metropolis-Hastings algorithm that uses a proximal proposal, while \cite{salim2019proximal,lau2022bregman,pereyra2016proximal,ehrhardt2024proximal} combine Langevin-type updates with proximal mappings. Meanwhile, a line of work develops distributed Gibbs samplers; see \cite{yuan2023class,vono2022admm,vono2019admm,rendell2020global}. For example, \cite{yuan2023class} considers targets of the form $\exp\!\big(-\sum_{i=1}^n f_i(x)-\sum_{j=1}^m g_j(x)\big)$ and designs Gibbs schemes for an augmented density
\[
\exp\!\left(-\sum_{i=1}^n f_i(x_i)-\sum_{j=1}^m g_j(y_j)-\sum_{i=1}^n\sum_{j=1}^m\frac{\sigma_{ij}}{2\eta}\|x_i-y_j\|^2\right),
\]
noting that when $\eta$ is sufficiently small, the augmented distribution is close to the original target. Other composite sampling works include \cite{eftekhari2023forward,ghaderi2024smoothing} which study Bregman forward-backward envelope smoothing schemes  and~\cite{habring2024subgradient} which study subgradient-based Langevin schemes. 





%%%%%%%%%%%%%%%%%%%%%%%%%%%%%%%%%%%%%%%%%%%
% \section{Literature Review}
% \red{shorten to one page or one and a half.}
% Works on uniform sampling on convex set: \cite{kook2024inandout,diaconis2012gibbs,narayanan2022mixing,laddha2023convergence,lovasz1999hit, cheneldanhitandrun,  
% chen2018fast,
% kannan2009random}


% \red{in the intro, cover the kook vempala algo stab, the yuan liang paper, mou et al. }


% We group constrained sampling papers according to the optimization operations they use or draw inspiration from.
% This list is not intended to be comprehensive. \red{list of paper+key operation applied to the constrain set.}

% \begin{itemize}
%     \item \cite{kook2025algodiffusion}: lifting, no $K$ in the paper so move this from constrained sampling to its own thing. \red{there are more paper. kook restart procedure. gaussian cooling paper with lifting.}
%     \item \cite{brosse2017sampling}: projection on $K$
%     \item \cite{kook2024inandout}: membership on $K$
%     \item \cite{kook2024gaussian}: barrier for $K$
%     \item \cite{kook2022sampling}: barrier for $K$
%     \item \cite{kannan2009random}: barrier
%     \item \cite{chen2018fast}: barrier
%     \item \cite{lamperski2021projected} projection
%     \item \cite{ahn2021efficient}: \red{unsure. either mirror or barrier.}
% \item \cite{li2022mirror} mirror. \red{there are more }
% \end{itemize}

% \red{move uniform sampling to a single paragrpah. make this the opennig paragraph.}


% \red{kook and vempala in the category of lifting $K$ plus membership}
% \begin{itemize}
%     \item \textbf{Constrained sampling with a proximal map.}
%     Kook and Vempala~\cite{kook2025algodiffusion} consider sampling from a log-concave target with density proportional to $e^{-V}$, given only an evaluation oracle for a convex $V$ satisfying a ground-set regularity condition (see Example~\ref{ex:groundset} in Appendix~\ref{appen:generaltechnical}).
%     They perform an exponential lifting as in~\eqref{silif:def:pi} and then run the proximal sampler (introduced in~\cite{leestructured21} and later significantly extended in~\cite{chen2022improved}).
%     The proximal step is implemented by rejection sampling, with a Failure event if the number of trials exceeds a prescribed cutoff. Consequently, for any $\eta,\varepsilon\in(0,1)$, with probability at least $1-\eta$ the output $X$ satisfies $R_q(\mathrm{Law}(X) \| \pi^X)\le \varepsilon$ where $R_q(\mu\|\nu)$ denotes the order-$q$ R\'enyi divergence.
%     The expected number of evaluation queries is
%     $\widetilde{O}\!\left(d^2  \mathrm{polylog} \frac{1}{\eta\varepsilon}\right)$.


%     \cite{brosse2017sampling} (see also \cite{durmus2018efficient}) study sampling from constrained log-concave targets
% $\pi(x)\propto \exp(-f(x)) \mathbf{1}_{\{x\in K\}}$ where $K$ is convex and compact and $f$ is smooth.
% They handle the hard constraint by replacing the indicator $\mathbf{1}_{\{x\in K\}}$ with the Moreau-Yosida envelope of the indicator potential $\iota_K$, whose gradient is expressed using the proximal map $\mathrm{prox}_{\lambda \iota_K}=\mathrm{proj}_K$. \red{this paper will be projection}
% They then run LMC on the smoothed potential, so each iteration requires evaluating $\nabla f$ and a projection onto $K$.
% They prove convergence in Wasserstein distance; for example, with smoothing level $\epsilon$ and step size $h=\mathcal{O}(\epsilon^2/d)$, they obtain $W_2(\mathcal{L}(x_N),\pi)\le \delta$ after $N=\widetilde{O}(d/\delta^2)$ iterations.






% \item \textbf{Constrained sampling with interior-point geometry/barrier methods.} \cite{kook2024gaussian} study Gaussian cooling with Dikin walk for sampling log-concave targets
% $\pi(x)\propto e^{-f(x)}\mathbf{1}_K(x)$ on a convex body $K$.
% They assume a barrier $\phi$ inducing a local metric $g$ that is $(\nu,\bar\nu)$-Dikin-amenable with $g$ comparable to $\nabla^2\phi$, and the relative Hessian bounds
% $\alpha \nabla^2\phi(x)\preceq \nabla^2 f(x)\preceq \beta \nabla^2\phi(x)$ for all $x\in K$.
% Their algorithm implements a Gaussian cooling schedule by sampling a sequence of truncated Gaussians on $K$ and replacing the classical Ball walk update~\cite{cousins2018gaussian} with a Dikin walk step in the barrier geometry.
% They obtain an $\varepsilon$-total-variation sample from $\pi$ after $\mathcal{O}\!\left(
% d(d\vee \nu\vee \bar\nu)\log\frac{d}{\varepsilon}
% \right)$ total Dikin-walk steps across all cooling stages.

% Next, for $A^{0}\in\R^{m\times n}, b^{0}\in\R^{m}$ and $A\in\R^{p\times n}$, \cite{kook2022sampling} study sampling from log-concave targets $\pi(x)\propto e^{-f(x)}$ supported on $M=\{x\in\R^n:\ A^{0}x\ge b^{0},\ Ax=0\}$ (an affine slice of a polytope). They assume $f$ is smooth and the inequality region admits a self-concordant barrier, which induces a Riemannian metric on $M$.
% They run Riemannian HMC in the barrier-induced metric, and analyze a discretization via an implicit midpoint integrator.
% Assuming a warm start, they prove total-variation mixing bounds $T=O\!\left(m k^{7/6}\log^2(1/\varepsilon)\right)$ for the idealized method and
% $T=O\!\left(m k^3\log^3(1/\varepsilon)\right)$ for the discretized method, where $m$ is the number of facets and $k=\dim(M)$. 



% \item \textbf{Constrained sampling with a projection map.}
% \cite{lamperski2021projected} studies sampling from a constrained target
% $\pi(x)\propto \exp(-\beta  f(x)) \mathbf{1}_{\{x\in K\}}$ on a compact convex set $K$ (allowing nonconvex $f$).
% Their method is a projected Langevin scheme: at each step, take a Langevin update
% $x_k-\eta\nabla  f(x_k)+\sqrt{2\eta/\beta} w_k$ and then project the result back onto $K$.
% Under their regularity assumptions (including compact convex $K$ and Lipschitz gradient of $\bar f$), they prove $W_1(\mathcal{L}(x_k),\pi)\le c_1 e^{-\eta a k}+c_2(\eta\log k)^{1/4}$. 
% They assume $\eta\le 1/2$, and recommend $\eta=\frac{\log T}{4aT}$ to obtain
% $W_1(\mathcal{L}(x_T),\pi)=O\!\left(T^{-1/4}(\log T)^{1/2}\right)$.

% \cite{Bubeck2018} sample the constrained log-concave target $\pi(x)\propto e^{-f(x)}\mathbf{1}_K(x)$ on a convex body $K\subset\R^d$, assuming $f$ is convex, Lipschitz, and smooth, and that $K$ has bounded diameter and contains a Euclidean ball.
% Their method is projected LMC: from $X_k$, take an Euler--Langevin step using $\nabla f(X_k)$ plus Gaussian noise, then project back onto $K$.
% They prove a finite-time total-variation bound without requiring a warm start, using a vanishing step size.
% Their iteration complexity for $\varepsilon$-accuracy in total variation distance is $\widetilde{O}\!\Big(d^3 m^6\cdot \max\{d^6\varepsilon^{-16},\ \varepsilon^{-22} m^8\}\Big)$ with $m=\mathbb{E}_{X\sim\pi}\|X\|$. 



% \cite{du2025non} (see also \cite{wang2025accelerating}) study sampling from a constrained target $\pi(x)\propto e^{-f(x)}$ on a bounded convex domain $C\subset\R^d$.
% They introduce a continuous-time non-reversible Langevin model with a skew-reflection term (built from $C$ and a skew-symmetric matrix $J$) that enforces the constraint $X_t\in C$, and then discretize it to obtain an algorithm.
% Their discretization (SRNLMC) takes an Euler-Maruyama step with drift $(I+J(x_k))\nabla f(x_k)$ and Gaussian noise, and then maps the point back to $C$ using a skew-projection operator built from the standard Euclidean projection.
% When $J\equiv 0$, this reduces to the projected LMC in~\cite{bubeck2018sampling}. For  step size $\eta=\widetilde{O}(\varepsilon^8/d^{10})$, the iteration complexity for SRNLMC in 1-Wasserstein metric is $\widetilde{O}(d^{10}/\varepsilon^8)$. 

% \red{add \cite{li2022mirror}}

% \item \textbf{Constrained sampling with a mirror map} \red{reread and confirm the mirror map do not have to be a barrier (page 6)}
% \cite{ahn2021efficient} assume a self-concordant, Legendre mirror map $\phi$ and study sampling from log-concave targets $\pi(x)\propto e^{-V(x)}$ supported on $K=\mathrm{dom}(\phi^\ast)$. Motivated by the viewpoint that Langevin dynamics can be interpreted as a gradient-flow evolution for $D_{\mathrm{KL}}(\cdot\|\pi)$ in geometry induced by 2-Wasserstein metric, they modify the geometry using the mirror map $\phi$. They impose relative regularity of $V$ with respect to $\phi$, namely relative smoothness/convexity and a relative Lipschitz condition in the local norm induced by $\nabla^2\phi$.
% Their Mirror Langevin Algorithm is a mirror-descent-style discretization: each iteration first takes the mirror-descent update for $V$ with respect to the Bregman divergence $D_\phi$, and then applies a diffusion step in the dual variable. 
% In the log-concave regime  and no relative strong convexity, to get $\epsilon$-accuracy in KL divergence, they take step size in the order of $\mathcal{O}(\varepsilon/d)$ to get the iteration complexity $\mathcal{O}(d/\varepsilon^2)$.

% \cite{srinivasan2024fast} study sampling from targets $\pi(x)\propto e^{-f(x)}$ supported on a compact convex body $K$, using a self-concordant mirror map $\phi$ that induces the mirror metric on $K$.
% They assume $f$ is $\mu$-strongly convex and $\lambda$-smooth, and is also $\beta$-Lipschitz in the mirror metric.
% Their algorithm uses a mirror-Langevin proposal followed by a Metropolis--Hastings correction.
% With a step size $\eta=\mathcal{O}(1/d^3)$, they obtain the total-variation mixing bound
% $T_{\mathrm{mix}}=O\!\left(\frac{1}{\mu}\max\{d^3,d\lambda,\beta^2\}\right)$ when $\mu>0$, and a similar bound in the log-concave case $\mu=0$ with additional geometry constants.

% \cite{zhang2020wasserstein} first paper to use mirror map to do constrained sampling. 

% \item \textbf{Other constrained sampling works.} \cite{chamon2024constrained} study constrained sampling via a primal-dual gradient-flow viewpoint. Another recent work that does sampling on constrained set using Langevin-type equations and penalty terms is \cite{gurbuzbalaban2024penalized}. 
% A separate and important line of works studies the Hit-and-Run algorithm and its variants; see, for example,
% \cite{ ascolani2024entropy}. Some other notable works are \cite{diaconis2012gibbs,brubaker2012family}. 


% \end{itemize}


% \red{revise the composite sampling lit review}
% \begin{itemize}
%     \item \cite{mou2022composite} this is proximal sampler and not proximal mapping. Mou should go with Yuan which do RGO stuff. 
%     \item salim and lau and liu are correct (use proximal map).
%     \item vono papers do gibb sampling (alternate)
% \end{itemize}

% \red{summary here}


% %%%%%%%%%%%%%%%%%%%%%%%%%%
% \red{below is the long version}

% We group composite sampling papers according to the optimization techniques they use or draw inspiration from.
% This list is not intended to be comprehensive.
% \begin{itemize}
%     \item \textbf{Composite sampling via a proximal map.}
%     \cite{mou2022composite} study sampling from a composite log-concave target on $\mathbb{R}^d$ with density proportional to
% $\exp(-(f(x)+g(x)))$.
% They assume: (i) the target satisfies a log-Sobolev inequality with constant $\lambda_\ast$,
% (ii) $f$ is smooth and has $L$-Lipschitz gradient, (iii) $f$ is distant dissipative
% $\langle \nabla f(x),x-x_0\rangle \ge (\mu/2)\|x-x_0\|_2^2-\beta$, and (iv) $g$ is convex and $M_d$-Lipschitz.
% Their algorithm is a Metropolis-Hastings chain with a proximal proposal: from $x$ and step size $\eta$, propose
% \[
% Y \sim p(x,\cdot)\ \propto\ \exp\!\left(-\frac{1}{4\eta}\|y-(x-\eta\nabla f(x))\|_2^2-g(y)\right),
% \]
% which requires a proximal sampling oracle for this density together with the knowledge of its normalizing constant
% $Z(u)=\int \exp\!\left(-\frac{1}{4\eta}\|y-u\|_2^2-g(y)\right) dy$.
% With $\eta=(2L^2R^2+Ld)^{-1}$ for a radius parameter $R$ defined in the paper, they prove the total-variation mixing bound
% $T_{\mathrm{mix}}(\varepsilon)=\widetilde{O}\!\left(d\log\frac{d}{\varepsilon}\right)$. \red{change $g$ to $h$ to be consistent with our setup}


% \cite{salim2019proximal} study sampling from a log-concave target with density proportional $\exp(-U(x))$ with a composite potential
% $U(x)=F(x)+\sum_{i=1}^n G_i(x)$.  
% They assume a stochastic structure $F(x)=\mathbb{E}[f(x,\xi)]$ and $G_i(x)=\mathbb{E}[g_i(x,\xi)]$ (where where $F$ is a.s. smooth and the $G_i$ are  a.s convex), and access to a stochastic gradient $\nabla f(x,\xi)$ for the smooth part together with stochastic proximal operators for the nonsmooth parts. 
% Motivated by operator splitting in optimization, they propose a Langevin splitting scheme that combines a stochastic gradient step for the smooth part, Gaussian noise and sequential proximal updates for the non-smooth parts to produce the next iterate.
% They establish nonasymptotic rates: in the convex case, a constant step size $\gamma=\mathcal{O}(\varepsilon)$ yields $\mathrm{KL}(\mu_k\|\mu^\star)\le \varepsilon$ after $\mathcal{O}(1/\varepsilon^2)$ iterations, while under strong convexity of $F$ they obtain linear convergence (up to a step-size-dependent neighborhood), and choosing $\gamma=\mathcal{O}(\varepsilon)$ gives an $\varepsilon$-accurate Wasserstein/KL guarantee in $\mathcal{O}((1/\varepsilon)\log(1/\varepsilon))$ iterations.  


% \cite{lau2022bregman} study sampling from nonsmooth convex composite targets of the form $\pi(x)\propto \exp(-(f(x)+g(x)))$, where $f$ is differentiable and $g$ can be non-differentiable. 
% They motivate the method from optimization in two ways: (i) a Legendre mirror map $\phi$ induces a non-Euclidean geometry, leading to mirror-descent-like Langevin schemes , and (ii) the nonsmooth part $g$ is handled by a proximal-type smoothing, replacing $g$ by its Bregman-Moreau envelope, defined via a Bregman proximity operator. 
% They then run a mirror-Langevin discretization on the smoothed surrogate potential, which recovers Moreau-Yosida Unadjusted Langevin algorithm when $\phi$ is Euclidean. 
% Assuming other problem constants are dimension-free, they obtain a mixing time bound of order $\widetilde{O}(d/\varepsilon^2)$ for their discrete scheme. 

% Other works that also use a proximal maps are~\cite{pereyra2016proximal,ehrhardt2024proximal}. 


% \item \textbf{Composite sampling via ADMM-type splitting.}
% \cite{yuan2023class} apply an ADMM-type variable splitting to the composite potential $f+g$ by introducing two variables $x,y\in\R^d$ (one for $f$ and one for $g$) and enforcing $x\approx y$ through the quadratic penalty $\|x-y\|_2^2/(2\eta)$.
% They sample from the augmented target on $(x,y)\in\R^d\times\R^d$ with density proportional to $\exp\!\left(-f(x)-g(y)-\frac{1}{2\eta}\|x-y\|_2^2\right)$
% using the Gibbs sampler.
% They assume $f$ and $g$ are strongly convex and smooth.
% For exact conditional sampling via rejection, choosing $\eta=\mathcal{O}(1/d)$ makes each conditional update $\mathcal{O}(1)$ in expectation and yields a total-variation mixing bound $O\!\left(d\log \frac{1}{\varepsilon}\right)$.
% They also analyze an inexact conditional implementation (based on~\cite{fan2023improved}): if each conditional sampler is within $\varepsilon_{\mathrm{RGO}}$ in total variation uniformly over the conditioning value, then with
% $\eta = O\!\left(\frac{1}{d^{1/2}\log(1/\varepsilon_{\mathrm{RGO}})}\right)$ one obtains $K=\widetilde{O}\!\left(d^{1/2}\right)$.

% Meanwhile, \cite{vono2022admm} (see also~\cite{vono2019admm,rendell2020global})  consider composite targets
% $\pi(\theta)\propto \exp(-\sum_{i=1}^b U_i(A_i\theta))$, motivated by settings where each term $U_i:\R^{d_i}\to\R$ is such that $d_i\ll d$. They apply an ADMM-type splitting by introducing $z_i\in\R^{d_i}$ and a quadratic penalty:
% \[
% z_i| \theta \sim \mathcal \mathcal{N}(A_i\theta,\rho^2 I_{d_i}),\qquad
% \Pi_\rho(\theta,z_{1:b})\propto
% \exp\!\left(-\sum_{i=1}^b U_i(z_i)-\frac{1}{2\rho^2}\sum_{i=1}^b\|z_i-A_i\theta\|_2^2\right).
% \]
% Their split Gibbs sampler of $\Pi_\rho$ alternates updates of the auxiliary blocks $z_{1:b}:=(z_1,\ldots,z_b)$ given $\theta$ (i.e., $z_i|\theta$) and a Gaussian update of $\theta| z_{1:b}$.
% They prove convergence to $\pi_\rho$ with rates independent of $d$ (for fixed $\rho$), and bound the bias $\pi_\rho$ versus $\pi$ as $\rho\to0$
% (e.g., total-variation bias $\mathcal{O}(\rho\sum_i \sqrt{d_i}L_i)$ under Lipschitz assumptions, and $\mathcal{O}(\rho^2\sum_i d_i M_i)$ under stronger smoothness/convexity conditions).
% Their dimension-free convergence assumes the conditional draws $z_i|\theta$ are available exactly.


% \item \text{Other composite sampling works.} Other notable works include \cite{eftekhari2023forward,ghaderi2024smoothing} which study Bregman forward-backward envelope smoothing schemes  and~\cite{habring2024subgradient} which study subgradient-based Langevin schemes. 

% \end{itemize}

\section{General Results}

\subsection{Outer analysis of the proximal sampler}\label{appen:outer} The following important result by~\cite{kook2025algodiffusion} is about contractivity in R\'{e}nyi divergence and $\chi^2$-divergence of the proximal sampler which is applicable to constrained targets. It is based on the main result by \cite{kook2024inandout}, which uses a smoothing argument to adapt the proof technique by \cite{chen2022improved} to the case of uniform sampling over a convex body. 

\begin{proposition}\label{prop:kook} 
(\cite[Lemma~2.9]{kook2025algodiffusion}) Let $\pi$ be a probability measure on $Q$ that is absolutely continuous w.r.t.\ Lebesgue, and let $P$ be the Markov kernel of the proximal sampler with step size $\eta$. Denote by $C_{\mathrm{PI}}$ the Poincar\'e constant of $\pi$. Then for any $\mu\ll\pi$ and any $q\ge2$ with $\mathcal{R}_q(\mu\|\pi)\le1$, we have
\[
\chi^2(\mu P\|\pi)\le \frac{\chi^2(\mu\|\pi)}{(1+\eta/C_{\mathrm{PI}}(\pi))^2}, \quad \mathcal{R}_q(\mu P\|\pi)\le \frac{\mathcal{R}_q(\mu\|\pi)}{(1+\eta/C_{\mathrm{PI}}(\pi))^{2/q}}.
\]
\end{proposition}


\subsection{Function and set oracles}
\label{appen:oracles}

This subsection collects elementary results about oracles for functions (evaluation, subgradient, and proximal) and sets (membership, separation, and projection). 
Throughout this subsection, we assume that both $f:\R^d\to \R$ and $K\subset \R^d$ are closed and convex.

\begin{lemma}\label{silif:lem:proxfas}
The proximal oracle for $f(x)$ gives a proximal oracle for $\tilde f(x,s)$ defined in \eqref{dolif:def:gammaandK}.
\end{lemma}

\begin{proof}
    Let $(x_0,s_0) \in \R^d \times \R$ and $\lam>0$ be given. By definitions of proximal mapping and $\tilde f(x,s)$ in \eqref{dolif:def:gammaandK}, we have
    \[
    \prox_{\lam \tilde f}(x_0, s_0)=\underset{x,s}\argmin\left\{f(x) + as + \frac{1}{2\lam}\|x-x_0\|^2 + \frac{1}{2\lam}(s-s_0)^2\right\}.
    \]
    This objective splits into an $x$-term and an $s$-term, so the proximal mapping can be easily obtained as
    \[
    \prox_{\lam \tilde f}(x_0, s_0) = (\prox_{\tilde f}(x_0), s_0 - \lam a).
    \]
    Hence, the lemma is proved.
\end{proof}

\begin{lemma}\label{lem:proj-sep}
    For a closed and convex set $K \subset \R^d$, a projection oracle for $K$ implies a separation oracle.
\end{lemma}

\begin{proof}
    Given $x\in \R^d$, it suffices to discuss the case $x\notin K$. Let $p=\proj_K(x)$. Since $x\ne p$, then $g=x-p \ne 0$ satisfies $\inner{g}{x} > \inner{g}{p} \ge \inner{g}{y}$ for every $y \in K$. Hence, $g$ is a separator for $x$ and $K$.
\end{proof}

\begin{lemma}\label{lem:sep}
    Evaluation and subgradient oracles of $f$ give a valid separation oracle for $\epi(f)$.
\end{lemma}

\begin{proof}
    Consider a query point $(\hat x,\hat t) \in \R^d \times \R$. First, evaluating $f(\hat x)\le \hat t$ checks whether $(\hat x,\hat t)\in \epi(f)$. If infeasible (i.e., $f(\hat x) > \hat t$), then pick a subgradient $g \in \partial f(\hat x)$, and convexity of $f$ gives
    \[
    f(y) \ge f(\hat x)+\inner{g}{y-\hat x} \quad \forall y.
    \]
    For $(y,s) \in \epi(f)$,
    \[
    s\ge f(y) \ge f(\hat x)+\inner{g}{y-\hat x},
    \]
    and hence
    \[
    \inner{(g,-1)}{(y,s)} \le \inner{g}{\hat x} - f(\hat x) < \inner{g}{\hat x} - \hat t = \inner{(g,-1)}{(\hat x,\hat t)}.
    \]
    Therefore, $(g,-1)$ is a valid separator for $(\hat x,\hat t)$ and $\epi(f)$.
\end{proof}

\begin{lemma}\label{lem:proj-epi}
    Evaluation and proximal oracles of $f$ give a valid projection oracle for $\epi(f)$. 
\end{lemma}

\begin{proof}
    Consider a query point $(\hat x,\hat t) \in \R^d \times \R$. First, evaluating $f(\hat x)\le \hat t$ checks whether $(\hat x,\hat t)\in \epi(f)$. If infeasible (i.e., $f(\hat x) > \hat t$), we consider
    \[
    \min_{x, t} \left\{\frac{1}{2}\|x-\hat x\|^2+\frac{1}{2}(t-\hat t)^2: f(x) \le t\right\}.
    \]
    Optimality conditions give $\lambda \geq 0$ such that
    \[
        \hat x-x \in \lam \partial f(x), \quad t-\hat t - \lam =0, \quad \lam (f(x) - t) = 0.
    \]
    So, we have $f(x) = t= \hat t + \lam$ and
    \[
    \hat x-x \in \lam \partial f(x) \Longleftrightarrow x=\prox_{\lam f}(\hat x).
    \]
    Define 
    \[
    \phi(\lam) = f(\prox_{\lam f}(\hat x)) -\hat t - \lam.
    \]
    Then $\phi$ is continuous and strictly decreasing on $[0, \infty)$, with $\phi(0^{+})=f(\hat x)-\hat t>0$ and $\phi(\lam) \rightarrow-\infty$ as $\lambda \rightarrow \infty$.
    Therefore, there is a unique $\lam_*>0$ such that $\phi(\lam_*)=0$, and the exact projection is
    \[
    x_* = \prox_{\lam_* f}(\hat x), \quad t_* = \hat t + \lam_*.
    \]

    Finally, we verify the monotonicity of $\phi(\lam)$. For $0< \lam_1 < \lam_2$, let $x_1= \prox_{\lam_1 f}(\hat x)$ and $x_2=\prox_{\lam_2 f}(\hat x)$, by optimality
    \begin{align*} 
    & \frac{1}{2}\|x_2-x_1\|^2 + \frac{1}{2}\|x_2-\hat x\|^2+\lam_2 f(x_2) \le \frac{1}{2}\|x_1-\hat x\|^2+\lam_2 f(x_1), \\ 
    & \frac{1}{2}\|x_2-x_1\|^2 + \frac{1}{2}\|x_1-\hat x\|^2+\lam_1 f(x_1) \le \frac{1}{2}\|x_2-\hat x\|^2+\lam_1 f(x_2) .
    \end{align*}
    Summing the two inequalities and rearranging, we have
    \[
    (\lambda_2-\lambda_1)(f(x_2)-f(x_1)) \le -\|x_2-x_1\|^2 <0 ,
    \]
    and hence $f(x_2)<f(x_1)$, i.e., $\phi$ is strictly decreasing. Therefore, $\lam_*$ can be solved analytically or numerically using a bisection search.
\end{proof}

We are ready to prove Lemma~\ref{silif:lem:sep}.

\noindent
\textbf{Proof of Lemma~\ref{silif:lem:sep}:} 
Without loss of generality, we assume $a=1$ in \eqref{silif:def:Q}, i.e., $Q= (K \times \R) \cap \epi(f)$. Consider a query point $(\hat x,\hat t) \in \R^d \times \R$. There are two infeasible cases: i) $\hat x \notin K$ and ii) $(\hat x, \hat t)\notin \epi(f)$. 

    Case i). Suppose $p\in \R^d$ is a separator for $\hat x$ and $K$, i.e., $\inner{p}{y} < \inner{p}{\hat x}$ for every $y \in K$ 
    Then the hyperplane given by $(p,0)$ separates $(\hat x,\hat t)$ and $Q$. Indeed,
    \[
    \inner{(p,0)}{(y,s)} \le \inner{(p,0)}{(\hat x,\hat t)}, \quad \forall (y,s) \in Q. 
    \]
    Case ii) directly follows from Lemma \ref{lem:sep}.
\QEDA

\begin{lemma}\label{silif:lem:sep2}
    Evaluation and proximal oracles of $f$, and membership and separation oracles of $K$ give a valid separation oracle for $Q$ as in \eqref{silif:def:Q}.
\end{lemma}

\begin{proof}
    We consider a query point $(\hat x,\hat t) \in \R^d \times \R$ and the same two cases as in the proof of Lemma \ref{silif:lem:sep}.
    The same argument applies to case i) here. For case ii), Lemma \ref{lem:proj-epi} gives a projection oracle for $\epi(f)$. Since $\epi(f)$ is closed and convex, it further implies a separation oracle for $\epi(f)$ through Lemma \ref{lem:proj-sep}.
\end{proof}


%%%%%%%%%%%%%%%%%%%%%%%%%%%%%%%%%%%%%%%%%%%%%%%%%%%%%%
\subsection{Technical results}
\label{appen:generaltechnical}

First, we recall some elementary results about Gaussian integrals from \cite[Lemma~A.1]{dangliang2025oracle}. 


\begin{lemma}\label{lem:gaussianint}
The following statements hold for any $\eta>0$, $c\in \R^d$ and $b\in \R$.
\begin{itemize}
    \item[(a)] $ \int_{\R^d} \exp\brac{-\frac{1}{2\eta} \norm{x-c}^2} \D x=(2\pi \eta)^{d/2}$;

    \item[(b)] $ \int_{\R^d} \exp\left(-\frac{1}{2\eta}(\|x-c\|-b)^2\right) \D x
    \le \exp \left(\frac14 + \frac{d}{\eta}b^2\right)(2\pi\eta)^{d/2}.$

\end{itemize} 
\end{lemma}





%%%%%%%%%%%%%%%%%%%%%%%%%%%%%%%%%%%%%%%%%%%%%%%%%%%%%%%%%%%%%%%%%%%%%



The next lemma says that if you thicken an epigraph by radius $r$, you still stay inside an epigraph, but with the defining function shifted downward by an amount of order $r$.

\begin{lemma}  
\label{lem:parallelsetcontainment} Let $\phi:\R^{m}\to\R$ be $L$-Lipschitz continuous and $\mathrm{epi}(\phi)$ be the epigraph of $\phi$, while $(\mathrm{epi}(\phi))_r$ is the associated parallel set. Then
\[
(\mathrm{epi}(\phi))_r\ \subseteq\ \left\{(u,v):\ v\ \ge\ \phi(u)-r\sqrt{1+L^2}\right\}.
\]
\end{lemma}

\begin{proof}
If $(x,s)\in(\mathrm{epi}(\phi))_r$ then there exists $(y,t)\in\mathrm{epi}(\phi)$ with
$\|(x,s)-(y,t)\|\le r$ and $t\ge \phi(y)$. Put $\alpha:=\|x-y\|$, $\beta:=|t-s|$ so $\alpha^2+\beta^2\le r^2$.
Then
\[
\phi(x)-s \le |\phi(x)-\phi(y)| + (\phi(y)-s) \le L\alpha + (t-s) \le L\alpha + \beta
\le \sqrt{1+L^2} \sqrt{\alpha^2+\beta^2}\le r\sqrt{1+L^2},
\]
which is equivalent to $s\ge \phi(x)-r\sqrt{1+L^2}$.
\end{proof}
%%%%%%%%%%%%%%%%%%%%%%%%%%%%%%%%%%%%%%%%%%%%%%%%%%%%%%%%%%%%%%%%%%%%%

In the main paper, we repeatedly convert integrals over the parallel sets $Q_r$ into integrals over the boundary sets $\partial Q_r$. The following lemma formalizes this conversion. 

\begin{lemma}
\label{lem:coarea}    
Let $Q\subset\R^{m}$ be a measurable set, $Z_{Q_r}:=\int_{Q_r} e^{-bt} \D y$, and 
$B_{Q_r}:=\int_{\partial Q_r} e^{-bt} \D S^{m-1}(y)$ where $t=t(y)$ is the last coordinate in $y$. 
Then, $Z_{Q_r}$ is absolutely continuous in $r$ and for a.e.\ $r\ge 0$,
\begin{equation}\label{eq:IrderivativeofQr}
\frac{\D}{\D r}Z_{Q_r}=B_{Q_r}.
\end{equation}
\end{lemma}

\begin{proof}
The co-area formula (e.g.,\ \cite[Thm.~3.2.12]{federer2014geometric}) says: if $\varphi:\R^{m}\to\R$ is Lipschitz and
$h\in L^{1}(\R^{m})$, then
\[
\int_{\R^{m}} h(y) |\nabla\varphi(y)| \D y
\;=\;
\int_{\R}\left(\int_{\{y:\varphi(y)= p\}} h(y) \D S^{m-1}(y)\right)\D p.
\]
For $y=(x,t)\in\R^{m-1}\times\R$, set $\varphi(y)=\dist (y,Q)$ and
$h(y)=\frac{e^{-bt}}{\|\nabla\varphi(y)\|}\mathbf 1_{\{y:\varphi(y)\le r\}}$.
Then, noting that $\{y:\dist (y,Q)\le r\}=Q_r$ and $\|\nabla \varphi(y)\| = \|\nabla \dist (y,Q)\|=1$ for $\mathcal{L}^{m}$ a.e.\ $y$, we have
\begin{equation}\label{eq:coarea-split}
Z_{Q_r} = \int_{\{y:\varphi(y)\le r\}} e^{-bt} \D y
=\int_{0}^{r}\left(\int_{\{y:\varphi(y)= p\}}
{e^{-bt}} \D S^{m-1}(y)\right)\D p.
\end{equation}
Thus, $Z_{Q_r}$ is absolutely continuous in $r$ and a.e. with respect to $r\geq 0$, 
\[
\frac{d}{\D r}Z_{Q_r}=\int_{\{y:\varphi(y)= r\}}
{e^{-bt}} \D S^{m-1}(y). 
\]
Moreover, we know $\partial Q_r=\{y:\dist (y,Q)= r\}$, so that~\eqref{eq:IrderivativeofQr} holds. 
\end{proof}


%%%%%%%%%%%%%%%%%%%%%%%%%%%%%%%%%%%%%%%%%%%%%%%%%%%%%%%%%%%%%%%%%%%%%

We record some elementary identities in the upcoming result. 
\begin{lemma} \label{lem:basicfacts}     
The following facts hold. 
\begin{enumerate}[label=\alph*)]
    \item Given $\eta>0$ and $\tau>0$, set $w_{\eta,\tau}(r):=\exp \left(-\frac{(r-\tau)^2}{2\eta}\right)$ and
let $C\in\R$. Then,
\begin{equation}\label{eq:kindofibp}
\int_{\tau}^{\infty} \frac{r-\tau}{\eta}e^{Cr}w_{\eta,\tau}(r) \D r
= C\int_{\tau}^{\infty} e^{Cr}w_{\eta,\tau}(r) \D r + e^{C\tau}.
\end{equation}
\item Let $a>0, \gamma\in\R$, $\psi:\R^d\to\R\cup\{+\infty\}$, and $F:\R^d\to[0,\infty]$ be measurable functions. Then,
\begin{equation}
    \label{eq:intoverQ}
   \int_{\R^d} F(x)\left(\int_{t\ge \psi(x)-\gamma} e^{-a t} dt\right)\D x
=\frac{e^{a\gamma}}{a}\int_{\R^d} F(x) e^{-a\psi(x)} \D x.  
\end{equation}
\end{enumerate}

\end{lemma}

\begin{proof}
Part~a is due to integration by parts. Define $F_C(r):=e^{Cr}w_{\eta,\tau}(r)$. Then
\[
F_C'(r)=Ce^{Cr}w_{\eta,\tau}(r)+e^{Cr}w_{\eta,\tau}'(r)
=\left(C-\tfrac{r-\tau}{\eta}\right)F_C(r),
\]
since $w_{\eta,\tau}'(r)=-(r-\tau)/\eta\;w_{\eta,\tau}(r)$.
Rearranging gives $\frac{r-\tau}{\eta}F_C(r)=C F_C(r)-F_C'(r)$.
Integrating over $[\tau,\infty)$ and using $F_C(r)\to 0$ as $r\to\infty$ and
$F_C(\tau)=e^{C\tau}$, we obtain the stated result \eqref{eq:kindofibp}.

Part~b is due to 
\[
    \int_{\R^d} F(x)\left(\int_{t\ge \psi(x)-\gamma} e^{-a t} dt\right)\D x
=\int_{\R^d} F(x) \left( \frac{1}{a}e^{-a(\psi(x)-\gamma)} \right) \D x \nonumber=\frac{e^{a\gamma}}{a}\int_{\R^d} F(x) e^{-a\psi(x)} \D x. 
\]
This completes the proof.
\end{proof}


At several locations in the main paper, we use a simple rejection-sampling scheme: propose from a tractable density proportional to $\exp(-\mathcal{P})$, and accept with the standard likelihood-ratio rule in order to target a density proportional to $\exp(-\Theta)$.
The following standard lemma records the resulting target law, the acceptance probability, and the average number of proposals.

\begin{lemma}\label{lem:rejection}  
Let $\Theta,\mathcal{P}:\R^{m}\to\R\cup\{+\infty\}$ be measurable and assume
\[
0<N_\Theta:=\int_{\R^m} e^{-\Theta(z)}  \D z<\infty,
\qquad
0<N_{\mathcal{P}}:=\int_{\R^m} e^{-\mathcal{P}(z)}  \D z<\infty,
\]
and that $\Theta(z)\ge \mathcal{P}(z)$ for all $z\in\R^m$.
Consider the rejection sampler that repeats the following trial until acceptance:
\begin{enumerate}
    \item draw a proposal $Z\in\R^m$ with density $k(z)=e^{-\mathcal{P}(z)}/N_{\mathcal{P}}$;
    \item draw $U\sim{\cal U}[0,1]$ independent of $Z$ and accept if and only if
    \begin{equation}\label{eq:generic-rejection-event}
        U\le \exp\left(-\Theta(Z)+\mathcal{P}(Z)\right).
    \end{equation}
\end{enumerate}
Then the output has density
\[
\pi(z)=\frac{e^{-\Theta(z)}}{N_\Theta}.
\]
Moreover, let $E$ denote the acceptance event \eqref{eq:generic-rejection-event} for a single trial and $F$ denote the number of trials until acceptance. Then, the acceptance probability and the average number of proposals are
\begin{equation}\label{eq:generic-accept-rate}
p:=\Pr(E)=\frac{N_\Theta}{N_{\mathcal{P}}},
\qquad
\mathbb{E}[F] =\frac{1}{p}=\frac{N_{\mathcal{P}}}{N_\Theta}.
\end{equation}
\end{lemma}

\begin{proof}
For one trial, let $E$ denote the event \eqref{eq:generic-rejection-event} happens.
Since $\Theta\ge \mathcal{P}$, we have $0\le \Pr(E| Z=z)=\exp(-\Theta(z)+\mathcal{P}(z))\le 1$,
so the acceptance rule is well-defined.

We first compute the acceptance probability:
\begin{align*}
p=\Pr(E)
=\int_{\R^m}\Pr(E| Z=z) k(z)  \D z &=\int_{\R^m}\exp\left(-\Theta(z)+\mathcal{P}(z)\right)\cdot \frac{e^{-\mathcal{P}(z)}}{N_{\mathcal{P}}}  \D z \\
&=\frac{1}{N_{\mathcal{P}}}\int_{\R^m}e^{-\Theta(z)}  \D z
=\frac{N_\Theta}{N_{\mathcal{P}}}.
\end{align*}
Next, we compute the conditional density of the accepted proposal:
\begin{align*}
k(z| E)
&=\frac{\Pr(E| Z=z) k(z)}{\Pr(E)} =\frac{\exp\left(-\Theta(z)+\mathcal{P}(z)\right)\cdot e^{-\mathcal{P}(z)}/N_{\mathcal{P}}}{N_\Theta/N_{\mathcal{P}}}
=\frac{e^{-\Theta(z)}}{N_\Theta}.
\end{align*}
Thus the accepted proposal has density $\pi(z)\propto e^{-\Theta(z)}$, as claimed.

Finally, each trial is i.i.d. and succeeds with probability $p$, hence $F\sim\mathrm{Geom}(p)$ and
$\E[F]=1/p$.
\end{proof}







%%%%%%%%%%%%%%%%%%%%%%%%%%%%%%%%%%%%%%%%%%%%%%%%%%%%%%%%%%%%%%%%%%%%%





The following example illustrates the ground-set assumption in~\cite[beginning of Sec.~2.1]{kook2025algodiffusion} is not automatic under our standing assumptions (A1) and (A2) in Section~\ref{sec:constrained:mainpaper}. 
\begin{example}
\label{ex:groundset}
Fix $d\ge 4$ and take $K=2\ball_{d}(0)\subset\R^d$ and $f(x)=d^3\sum_{i=1}^d|x_i|$ on $K$. Then $V=f$ is convex and Lipschitz on $K$ and (A1)-(A2) holds. Moreover, we have  $\min_{x\in K}V(x)=0$. In the notation of \cite{kook2025algodiffusion},
\[
\mathsf L_{V,g}
=\left\{x:\ V(x)-\min_{x\in K}V(x)\le 10d \right\}
=\left\{x:\ \sum_{i=1}^d |x_i|\le \frac{10}{d^2}\right\}
\subseteq \left[-\frac{10}{d^2},\frac{10}{d^2}\right]^d,
\]
so $\mathsf L_{V,g}$ contains no unit ball. This violates the ground-set assumption in \cite{kook2025algodiffusion}. 
\end{example}
We have a similar example in the composite setting by taking $h\equiv 0$, $f(x):=d^3\sum_{i=1}^d |x_i|$ on $\R^d$, which satisfies our (B1) and (B2) but yields the same violation to the ground-set condition.

\section{Supporting Results and Proofs for Section \ref{sec:constrained:mainpaper}}


\subsection{Proof of Theorem~\ref{silif:theo:separation}}
\label{silif:appen:sepatheoproof}
 
 Regarding Part~a, the proof sketch in Section~\ref{sec:constrained:mainpaper} has provided all the details, so what remains is to show Part~b. Denote $\mu_k$ the distribution of $(y_k,s_k)$ for the first step of Algorithm \ref{silif:alg:ASF_pi} and $n_{y,s}$ the expected number of proposals conditioned on the RGO input $(y_k,s_k)$. The fact that $d\mu_k/d\Pi^{Y,S}\leq M$ from Lemma~\ref{silif:lem:warmstart} implies 
\begin{align}
\label{silif:sepawarmstartimplication}
	\E_{\mu_k}[n_{y,s}] \le M \E_{\Pi^{Y,S}}[n_{y,s}], 
\end{align}
so our focus will be on bounding $\E_{\Pi^{Y,S}}[n_{y,s}]$.
 
In Algorithm~\ref{silif:RGO}, Steps~2 and~3 are a rejection sampler where the true potential function $\Theta$ and its proposal $\mathcal{P}$ are respectively 
 \begin{equation}\label{eq:corr2}
      \Theta = \Theta^{\eta,Q}_{y,s}, \quad \mathcal{P} = \mathcal{P}_1
 \end{equation}
 in view of Lemma \ref{lem:rejection}, so that the latter result applies. It follows from $n_{y,s}=\E[F]$ in~\eqref{eq:generic-accept-rate} with the specification \eqref{eq:corr2} and inequality~\eqref{silif:ineq:P2-P3} in Lemma~\ref{silif:lem:compareP1P2} that
\begin{align}
\label{silif:for:averagerejectionnotwarmsepa}
 \E_{\Pi^{Y,S}} [n_{y,s}]
 &\stackrel{\eqref{eq:generic-accept-rate}}= \int_{\R^{d+1}} \frac{\int_{\R^{d+1}} \exp(-\mathcal{P}_1(x,t))  \D x  \D t}{\int_{\R^{d+1}} \exp(-\Theta^{\eta,Q}_{y,s}(x,t))  \D x  \D t} \Pi^{Y,S}(y,s) \D y  \D s\nn\\
 & \stackrel{\eqref{silif:ineq:P2-P3}}\leq \int_{\R^{d+1}} \frac{\int_{\R^{d+1}} \exp(-\mathcal{P}_2(x,t))  \D x  \D t}{\int_{\R^{d+1}} \exp(-\Theta^{\eta,Q}_{y,s}(x,t))  \D x  \D t} \Pi^{Y,S}(y,s) \D y  \D s .
\end{align}
Now via the definition of $\mathcal{P}_2$ at~\eqref{silif:def:P2} with $a=d$ and Part~b of Lemma~\ref{lem:gaussianint}, 
\begin{align}
\label{silif:sepa:boundintP3}
&\int_{\R^{d+1}} \exp(-\mathcal{P}_2(x,t))  \D x  \D t
\leq \exp\brac{\frac{1}{4}+\frac{2(L+d)^2}{d^2}}(2\pi\eta)^{(d+1)/2}\nn\\
&\cdot \exp\brac{-\frac{1}{2\eta}\left(\|z-\operatorname{proj}_Q(z)\| - \frac{2(L+d)}{d}\sqrt{\frac{2\eta}{d+1}}\right)^2 }
\cdot\exp\brac{\frac{d^2\eta}{2} - d s+\frac{16(L+d)^2}{d^2(d+1)} }, 
\end{align}
noting that $z=(y,s-d\eta)$. 
Moreover, applying Lemma~\ref{lem:gaussianint}(a) and the definition of $\Theta^{\eta,Q}_{y,s}$ in~\eqref{silif:def:Theta}, we have
\begin{align}
    \Pi^{Y,S}(y,s)\stackrel{\eqref{silif:def:bigpi}}=\frac{1}{(2\pi \eta)^{(d+1)/2} \int_{Q} \exp(- d t)\D x \D t} \int_Q \exp\brac{- d t-\frac{1}{2\eta}\norm{(x,t)-(y,s)}^2 }\D x \D t, \label{silif:for:piYS} \\
    \int_{\R^{d+1}} \exp(-\Theta^{\eta,Q}_{y,s}(x,t))  \D x  \D t \stackrel{\eqref{silif:def:Theta}}= \int_Q \exp\brac{- d t-\frac{1}{2\eta}\norm{(x,t)-(y,s)}^2 }\D x \D t. \label{eq:int-theta}
\end{align}
Plugging \eqref{silif:sepa:boundintP3}, \eqref{silif:for:piYS}, and \eqref{eq:int-theta} into \eqref{silif:for:averagerejectionnotwarmsepa}, we obtain
\begin{align}
\label{silif:separation:intermediate}
     &\E_{\Pi^{Y,S}} [n_{y,s}]
     \stackrel{\eqref{silif:for:averagerejectionnotwarmsepa}}\leq \frac{1}{\int_Q\exp(-dt) \D x \D t}\exp\brac{\frac{1}{4}+\frac{2(L+d)^2}{d^2}+\frac{16(L+d)^2}{d^2(d+1)}}\nonumber\\
     &\quad\cdot\int_{\R^{d+1}}\exp\brac{\frac{d^2\eta}{2} - ds }\exp\brac{-\frac{1}{2\eta}\left(\|z-\operatorname{proj}_Q(z)\| - \frac{2(L+d)}{d}\sqrt{\frac{2\eta}{d+1}}\right)^2 } \D y \D s.
\end{align}
Recall the notation $z_{d+1}=(y,s-d\eta)_{d+1}=s-d\eta$. Then
\begin{align}
  &\int_{\R^{d+1}}\exp\brac{\frac{d^2\eta}{2} - d s }\exp\brac{-\frac{1}{2\eta}\left(\|z-\operatorname{proj}_Q(z)\| - \frac{2(L+d)}{d}\sqrt{\frac{2\eta}{d+1}}\right)^2 } \D y \D s \nn\\
    =&\exp\brac{-\frac{d^2\eta}{2}}\int_{\R^{d+1}} \exp(-d z_{d+1}) 
    \exp\left(-\frac{\left(\dist(z,Q)- \tau\right)^2}{2\eta}\right)  \D z, \label{obs}
\end{align}
where
\begin{equation}\label{ineq:tau}
    \tau = \frac{2(L+d)}{d}\sqrt{\frac{2\eta}{d+1}} \leq \frac{3(L+d)}{d^{5/2}}.
\end{equation}
Applying Lemma~\ref{silif:lem:keytechnical} to \eqref{obs} with this $\tau$, plugging the resulting inequality into \eqref{silif:separation:intermediate}, and using $\eta=1/d^2$, we conclude that
\begin{align*}
 \E_{\Pi^{Y,S}} [n_{y,s}]
 &\stackrel{\eqref{ineq:lem5},\eqref{silif:separation:intermediate}}\leq \exp\brac{\frac{2(L+d)^2}{d^2}+\frac{16(L+d)^2}{d^2(d+1)} + \tau C_L}
 \cdot  \Biggl[ \frac{\sqrt{2\pi} C_L}{d}\exp \left(\frac{C_L^2}{2d^2}\right) + 1 \Biggr] \\
 &\stackrel{\eqref{ineq:tau}}\leq \exp\brac{\frac{2(L+d)^2}{d^2}+\frac{16(L+d)^2}{d^2(d+1)} + \frac{3C_L(L+d)}{d^{5/2}}} \cdot\Bigg[\frac{\sqrt{2\pi} C_L}{d}
\exp \Biggl(\frac{C_L^2}{2d^2}\Biggr)+1\Bigg].
\end{align*}
Therefore, the stated bound \eqref{silif:proofboundsepa} immediately follows from \eqref{silif:sepawarmstartimplication}.

Finally, recall that $C_L=\sqrt{d^2+L^2}+LR+d$ from~\eqref{silif:def:notation}, then assuming $L={\cal O}(\sqrt{d})$, $R={\cal O}(\sqrt{d})$, and $M={\cal O}(1)$, one can verify that $C_L={\cal O}(d)$ and \eqref{silif:proofboundsepa} becomes ${\cal O}(1)$.
\QEDA

% Finally, setting $a=d$, we have $\left(1+\frac{L}{a}\right)=1+O(1/d)$, so the prefactor on the right hand side of~\eqref{silif:proofboundsepa} satisfies
% \[
% \exp\!\left(\frac{1}{4}+2\left(1+\frac{L}{a}\right)^2+16\left(1+\frac{L}{a}\right)^2\frac{1}{d+1}\right)
% = e^{9/4}\left(1+O(1/d)\right).
% \]
% Recall $\tau=2\left(1+\frac{L}{a}\right)\sqrt{\frac{2\eta}{d+1}}$ and $C_L=\sqrt{a^2+L^2}+LR+d$, so under $a=d,\eta=1/d^2$ we have $\tau=O(d^{-3/2})$ and hence $\tau C_L=O\left(d^{-\frac{1}{2}}\right)$

% Inside the square bracket on the right hand side of~\eqref{silif:proofboundsepa}, the exponent
% $2\left(1+\frac{L}{a}\right)C_L\frac{2}{d^{3/2}}$ is $O\left(d^{-\frac{1}{2}}\right)$, hence
% \[
% C_L\sqrt{\frac{2\pi}{d^2}}
% \exp\!\left(2\left(1+\frac{L}{a}\right)C_L\frac{2}{d^{3/2}}+\frac{C_L^2}{2d^2}\right)
% =
% \left(2\sqrt{2\pi} e^{2}\right)\left(1+O\left(d^{-\frac{1}{2}}\right)\right),
% \]
% while the last exponential equals
% $\exp\!\left(2C_L\left(1+\frac{L}{a}\right)\frac{2}{d^{3/2}}\right)=1+O\left(d^{-\frac{1}{2}}\right)$.
% Therefore, the square bracket lem:volis
% \[
% 2+2\sqrt{2\pi} e^{2}+O\left(d^{-\frac{1}{2}}\right).
% \]
% Therefore, the right hand side of~\eqref{silif:proofboundsepa}  simplifies to
% \[
% M e^{9/4}\left(2+2\sqrt{2\pi} e^{2}+O\left(d^{-\frac{1}{2}}\right)\right),
% \]
% which is the second stated bound in Part~b of Theorem~\ref{silif:theo:separation}. 


\subsection{Missing proofs}
\label{silif:appen:missingproofs}
This appendix contains  proofs not available in Subsection~\ref{silif:sec:asf} and Section~\ref{sec:constrained:mainpaper} due to space limit. 


\medskip
   \textbf{Proof of Theorem~\ref{silif:theo:outer}}:
  Per Lemma~\ref{silif:lem:warmstart}, the $M$-warm start condition~(A3) implies $\chi^2\brac{\pi_0||\pi}\leq M^2-1$. Then by applying $\chi^2(\mu P\|\pi)\le \frac{\chi^2(\mu\|\pi)}{(1+\eta/C_{\mathrm{PI}}(\pi))^2}$ in Proposition~\ref{prop:kook} iteratively,  we can solve for 
\[
        \frac{M^2}{\brac{1+\eta/C_{\mathrm{PI}}(\pi)}^{2k}} \leq \epsilon
\]
to get 
\[
k \ \ge\  \frac{\log\!\left(\frac{M^2}{\epsilon}\right)}{2\log\!\left(1+\eta/C_{\mathrm{PI}}(\pi)\right)}.
\]
In particular, using $\log(1+u)\ge \frac{u}{1+u}$ for $u>0$, it suffices to take
\[
k \ \ge\ \frac{C_{\mathrm{PI}}(\pi)+\eta}{2\eta}\,\log\!\left(\frac{M^2}{\epsilon}\right)
\ =\ \mathcal{O}\!\left(\frac{C_{\mathrm{PI}}(\pi)}{\eta}\log\!\left(\frac{M^2}{\epsilon}\right)\right).
\]
Moreover, Lemma~\ref{silif:lem:PIconstant} bounds $C_{\mathrm{PI}}(\pi)$. Therefore, we get the stated iteration complexity with respect to $\chi^2$ divergence in Theorem \ref{silif:theo:outer}. 




The calculation for R\'{e}nyi divergence is along the same line with the use of the one-step contraction $\mathcal{R}_q(\mu P\|\pi)\le \frac{\mathcal{R}_q(\mu\|\pi)}{(1+\eta/C_{\mathrm{PI}}(\pi))^{2/q}}$ in Proposition~\ref{prop:kook} and Lemma~\ref{silif:lem:PIconstant}. In particular, the condition $M\leq e^{1-1/q}$ implies $\mathcal{R}_q\brac{\pi_0||\pi}\leq \frac{q}{q-1}\log M\leq 1$, so that the one-step contraction is applicable. This completes the proof. 
\QEDA


%%%%%%%%%%%%%%%%%%%%%%%%%%%%%%%%%%%%%%%%%%%%%%%%%%%%%%%%%%%%%%%%%%%%%%%%%%%%%%%%%%%%%%%%%%%%%%%%%%%%%%%%%%%%%%%%%%%%%%%%%%%%%%%%%%%%%%%%%%%%%%%%%%%%%%%%%%%%%%%%%%%%%%%%%%

\medskip
\textbf{Proof of Lemma~\ref{silif:lem:keytechnical}:}    
Split the integral over $Q$ and $  Q^c$. For $w=(x,t) \in Q$, $\dist(w,Q)=0$, hence we have
\[
\int_Q e^{-a t}\exp \left(-\frac{(\dist(w,Q)-\tau)^2}{2\eta}\right)  \D w
=\exp \left(-\frac{\tau^2}{2\eta}\right)\int_Q e^{-a t}  \D w
\stackrel{\eqref{def:ZandB}}=Z_Q\exp \left(-\frac{\tau^2}{2\eta}\right),
\]
where the second relation is due to the definition of $Z_Q$ in \eqref{def:ZandB}.
On $Q^c$, applying the co-area formula, we obtain
\begin{align*}
&\int_{  Q^c} e^{-a t}\exp \left(-\frac{(\dist(w,Q)-\tau)^2}{2\eta}\right)  \D w
=\int_{0}^{\infty}\left(\int_{\{y:\dist(y,Q)=r\}}
e^{-a t}\exp \left(-\frac{(r-\tau)^2}{2\eta}\right)  \D S(y)\right)\D r\\
&\quad=\int_0^\infty \exp \left(-\frac{(r-\tau)^2}{2\eta}\right)  
\left(\int_{\{y:\dist (y,Q)=r\}} e^{-a t}  \D S(y)\right)\D r\stackrel{\eqref{def:ZandB}}=\int_0^\infty \exp \left(-\frac{(r-\tau)^2}{2\eta}\right)  B_{Q_r}  \D r,
\end{align*}
where the last identity is due to $\partial Q_r=\{y:\dist(y,Q)= r\}$ and the definition of $B_{Q_r}$ in \eqref{def:ZandB}.
Combining the above two relations yields 
\[ \int_{\R^{d+1}} e^{-a t}  
\exp \left(-\frac{(\dist(w,Q)-\tau)^2}{2\eta}\right)  \D w
=
Z_Q\exp \left(-\frac{\tau^2}{2\eta}\right)
+ \int_0^\infty \exp \left(-\frac{(r-\tau)^2}{2\eta}\right)  B_{Q_r}  \D r.\]
Then we can bound the integral on the right hand side by Lemma~\ref{silif:lem:envelope} to get
\begin{align*}
& \int_{\R^{d+1}} e^{-a t}  
\exp \left(-\frac{(\dist(w,Q)-\tau)^2}{2\eta}\right)  \D w\\
&\leq  Z_Q\left[\exp \left(-\frac{\tau^2}{2\eta}\right)+C_L   \sqrt{2\pi\eta}
\exp \left(\tau C_L + \frac{\eta C_L^2}{2}\right)+e^{C_L\tau}-1\right]. 
\end{align*}
 The stated estimate in Lemma~\ref{silif:lem:keytechnical} follows after applying the bound $\exp \left(-\frac{\tau^2}{2\eta}\right)\le 1$. 
\QEDA





\medskip\textbf{Proof of Lemma~\ref{silif:lem:envelope}:}   The proof consists of two steps. In Step 1, we construct an envelope set which contains $Q_r$ and is much easier to integrate on compared to $Q_r$; while in Step 2, we perform integration by parts (IBP) to obtain the desired bound. The reason we need IBP to bound the quantity $\int_{0}^{\infty} \exp \left(-\frac{(r-\tau)^2}{2\eta}\right)  
B_{Q_r}  \D r$ is that $B_{Q_r}$ as an integral on the boundary $\partial Q_r$ is quite hard to deal with; so via IBP, $B_{Q_r}$ is replaced by $Z_{Q_r}$, which is easier to handle.  

\smallskip
\emph{Step 1.}
Since $f$ is only assumed on $K$, we extend it to $\R^d$ as follows.
Let $\proj_K:\R^d\to K$ be the Euclidean projection onto the closed convex set $K$,
and define $\bar f(x):=f(\proj_K(x)),x\in\R^d$. Then $\bar f=f$ on $K$ and since the projection map onto the closed convex set $K$ is $1$-Lipschitz, $\bar f$ is $L$-Lipschitz on $\R^d$: 
\begin{equation}\label{ineq:phi-Lip}
    \abs{\bar f(x)-\bar f(y)} = \abs{f(\proj_K(x))-f(\proj_K(y))}\leq  L\abs{\proj_K(x)-\proj_K(y)}\leq L\norm{x-y}. 
\end{equation}
Recall the definition of $Q$ in \eqref{silif:def:Q}, we write $Q=S_0\cap S_1 \subset \R^{d+1}$ where $S_0:=K\times \R, S_1:=\{(x,t)\in\R^{d+1}:\ t\ge \bar f(x)/a\}$. Fix $r>0$. Since $(A\cap B)+H\subseteq (A+H)\cap(B+H)$ for any sets $A,B,H$, applying this observation to $A=S_0,B=S_1$ and $H=r\ball_{d+1}(0)$, we get
\begin{align}
\label{silif:setQr}
Q_r=(S_0\cap S_1)_r \subseteq (S_0)_r\cap (S_1)_r,
\end{align}
where $(S_0)_r = K_r\times \R$ and $K_r=K+r\ball_d(0)\subset\R^d$.
Considering $\phi=\bar f/a$, we note that $\phi$ is $(L/a)$-Lipschitz continuous in view of \eqref{ineq:phi-Lip} and $S_1=\mathrm{epi}(\phi)$.
Hence, applying Lemma~\ref{lem:parallelsetcontainment} with $\phi=\bar f/a$, we obtain 
\[
(S_1)_r \subseteq \left\{(x,t):\ t\ge \frac{\bar f(x)}{a}-r\sqrt{1+\frac{L^2}{a^2}}\right\}=\left\{(x,t):\ t\ge \frac{\bar f(x)}{a}-\frac{r}{a}C_f\right\},
\]
where $C_f$ is as in \eqref{silif:def:notation}.
Therefore, via \eqref{silif:setQr}, we conclude 
\begin{align}
\label{silif:envelopeinclusion}
Q_r\subseteq \mathcal E_r:=\left\{(x,t):x\in K_r, t\ge \frac{\bar f(x)}{a}-\frac{r}{a}C_f \right\}.
\end{align}
The definition of $Z_S$ for a set $S$ in \eqref{def:ZandB}, the inclusion in~\eqref{silif:envelopeinclusion}, and Lemma~\ref{lem:basicfacts}(b) (under the specification $F(x)=\mathbf{1}_{K_r}(x)$, $\psi(x)=\bar f(x)/a$, $\gamma=rC_f/a$) imply that
\begin{equation}\label{silif:firstbound_Zer}
    Z_{Q_r}\stackrel{\eqref{def:ZandB},\eqref{silif:envelopeinclusion} }{\le} Z_{\mathcal E_r}\stackrel{\eqref{def:ZandB}}{=}\int_{K_r}\int_{t\ge (\bar f(x)-rC_f)/a} e^{-a t}  \D t  \D x
\stackrel{\eqref{eq:intoverQ}}{=}\frac{e^{rC_f}}{a}\int_{K_r} e^{-\bar f(x)}  \D x. 
\end{equation}
Since $0\in K$ and $\ball_d (0)\subseteq K$, we have $r\ball_d (0)\subseteq rK$ and hence $K_r=K+r\ball_d (0)\subseteq K+rK=(1+r)K$. Let $\lambda:=1+r$ and $y:=x/\lam$. Then, we have
\begin{align}
\label{silif:firstboundintKr}
\int_{K_r} e^{-\bar f(x)}  \D x
\le \int_{\lambda K} e^{-\bar f(x)}  \D x
=\lambda^d\int_{K} e^{-\bar f(\lambda y)}  \D y.
\end{align}
% Using the facts that $\proj_K(y)=y$ for $y\in K$ and $\proj_K$ is $1$-Lipschitz, we have
% \[
% \|\proj_K(\lambda y)-y\|
% =\|\proj_K(\lambda y)-\proj_K(y)\|
% \le \|\lambda y-y\|
% =(\lambda-1)\|y\|=r\|y\|.
% \]
It follows from \eqref{ineq:phi-Lip} with $(x,y)=(\lam y, y)$ and $\|y\|\le R$ for $y\in K$ that 
\[
\bar f(\lambda y) \ge \bar f(y)-L r\|y\|
\ge f(y)-LR  r,
\]
where we also use the fact that $\bar f(y) = f(y)$ for $y \in K$.
Plugging the above inequality into \eqref{silif:firstboundintKr} and using the fact that $\lambda^d=(1+r)^d\le e^{dr}$ for $r\ge 0$, we have
\begin{align}
\label{silif:secondboundintKr}
\int_{K_r} e^{-\bar f(x)}  \D x
\le \exp \left((LR+d)r\right)\int_{K} e^{-f(y)}  \D y
= e^{rC_K}\int_K e^{-f(y)}  \D y,
\end{align}
where $C_K$ is as in \eqref{silif:def:notation}.
Using the definition of $Q$ in \eqref{silif:def:Q} and Lemma~\ref{lem:basicfacts}(b) with specification $F(x)=\mathbf{1}_{K}(x)$, $\psi(x)=f(x)/a$, and $\gamma=0$, we have
\[
\frac{1}{a} \int_K e^{-f(y)} \D y \stackrel{\eqref{eq:intoverQ}}= \int_K \int_{t\ge f(y)/a} e^{-at} \D t  \D y \stackrel{\eqref{silif:def:Q}}= \int_{Q} e^{-a t} \D w \stackrel{\eqref{def:ZandB}}= Z_Q,
\]
where the last identity is due to \eqref{def:ZandB}.
Therefore, plugging the above relation and \eqref{silif:secondboundintKr} into \eqref{silif:firstbound_Zer}, we obtain 
\begin{equation}
\label{silif:lemmafirststep}
Z_{Q_r} \stackrel{\eqref{silif:firstbound_Zer},\eqref{silif:secondboundintKr}}\le \frac{e^{r(C_f+C_K)}}{a} \int_K e^{-f(y)}  \D y
=e^{r(C_f+C_K)}  Z_Q
\stackrel{\eqref{silif:def:notation}}=e^{rC_L}  Z_Q,
\end{equation}
where $C_L$ is defined in \eqref{silif:def:notation}.

%%%%%%%%%%%%%%%%%%%%%%%%%
\emph{Step 2.} Set $w_{\eta,\tau}(r):=\exp \left(-(r-\tau)^2/(2\eta)\right)$.
Since $w_{\eta,\tau}(r)$ decays exponentially, the conclusion~\eqref{silif:lemmafirststep} of Step 1 implies that $Z_{Q_r}w_{\eta,\tau}(r)\to0$ as $r\to\infty$.
Recall from Lemma~\ref{lem:coarea} that
$\frac{\D}{\D r}Z_{Q_r}=B_{Q_r}$ a.e.\ $r$.
Hence, integrating by parts yields
\begin{align}\label{silif:eq:IBP-master}
&\int_{0}^{\infty} B_{Q_r}   w_{\eta,\tau}(r)  \D r
\stackrel{\eqref{eq:IrderivativeofQr}}=\int_{0}^{\infty} Z_{Q_r}'   w_{\eta,\tau}(r)  \D r
= Z_{Q_r}w_{\eta,\tau}(r)\bigg|_{0}^{\infty}
-\int_{0}^{\infty} Z_{Q_r}   w_{\eta,\tau}'(r)  \D r\nonumber\\
\quad&= -Z_Q e^{-\tau^2/(2\eta)}
+ \int_{0}^{\infty} Z_{Q_r}  \frac{r-\tau}{\eta}  w_{\eta,\tau}(r)  \D r. 
\end{align}
We will bound the second term on the right hand side of~\eqref{silif:eq:IBP-master} using the fact that $Z_Q\leq Z_{Q_r}$ and \eqref{silif:lemmafirststep},  
\begin{align}
    \int_{0}^{\infty} Z_{Q_r}  \frac{r-\tau}{\eta}  w_{\eta,\tau}(r)  \D r &= \int_0^\tau Z_{Q_r}  \frac{r-\tau}{\eta}  w_{\eta,\tau}(r)  \D r + \int_\tau^\infty Z_{Q_r}  \frac{r-\tau}{\eta}  w_{\eta,\tau}(r)  \D r \nn \\
    &\stackrel{\eqref{silif:lemmafirststep}}\leq Z_Q\int_0^\tau \frac{r-\tau}{\eta}    w_{\eta,\tau}(r)  \D r+ Z_Q\int_\tau^\infty \frac{r-\tau}{\eta}  e^{rC_L}  w_{\eta,\tau}(r)  \D r.\label{ineq:ZQ}
\end{align}
In view of the definition of $w_{\eta,\tau}(r)$, it is easy to compute
\[
\int_0^\tau \frac{r-\tau}{\eta}    w_{\eta,\tau}(r) \D r = - w_{\eta,\tau}(r) \bigg|_{0}^{\tau} 
    = - 1 + e^{-\tau^2/(2\eta)}.
\]
Combining the above observation, \eqref{silif:eq:IBP-master}, and \eqref{ineq:ZQ}, and using Lemma~\ref{lem:basicfacts} with $C=C_L$, we have
\begin{align}
     \int_{0}^{\infty} B_{Q_r}   w_{\eta,\tau}(r)  \D r & \stackrel{\eqref{silif:eq:IBP-master},\eqref{ineq:ZQ}}\leq -Z_Q + Z_Q\int_\tau^\infty \frac{r-\tau}{\eta}  e^{rC_L}  w_{\eta,\tau}(r)  \D r \nn \\
     &\stackrel{\eqref{eq:kindofibp}} \le Z_Q\left(e^{C_L\tau}-1\right)+C_L Z_Q\int_{\tau}^{\infty} e^{rC_L}w_{\eta,\tau}(r)  \D r. \label{finalboundintBQr}
\end{align}
Finally, using the definition of $w_{\eta,\tau}(r)$ and noting 
\[
C_L r-\frac{(r-\tau)^2}{2\eta}
=-\frac{1}{2\eta}\left(r-(\tau+\eta C_L)\right)^2+\tau C_L+\frac{\eta C_L^2}{2},
\]
we obtain 
\[
\int_{\tau}^{\infty} e^{rC_L}w_{\eta,\tau}(r)  \D r
\le e^{\tau C_L+\frac{\eta C_L^2}{2}}
\int_{-\infty}^{\infty}\exp \left(-\frac{(r-(\tau+\eta C_L))^2}{2\eta}\right)  \D r=\sqrt{2\pi\eta}\;\exp \left(\tau C_L+\frac{\eta C_L^2}{2}\right),
\] 
where the identity is due to Lemma \ref{lem:gaussianint}(a).
Therefore, the conclusion immediately follows from \eqref{finalboundintBQr}.
\QEDA


\subsection{Supporting lemmas}
\label{silif:appen:supportlem}
We collect supporting lemmas for Section~\ref{sec:constrained:mainpaper} below. First, we show how to bound the PI constant of the lifted distribution~\eqref{silif:def:pi}. 



\begin{lemma}
\label{silif:lem:PIconstant}
    Assume the distribution $\pi$ at \eqref{silif:def:pi} has $a=d$, $f$ is a convex function and $K$ is a convex set. Then, it satisfies a PI and the PI constant is bounded as $  C_{\mathrm{PI}}(\pi)\leq C\log (d+1)\brac{\opnorm{\cov{\pi^X}}+1 }$ for a universal constant $C$. 
\end{lemma}

\begin{proof}   Note that $\pi(x,t)\propto \exp\brac{-dt-I_Q(x,t)}$ is a log-concave measure on $\R^{d+1}$. It is a well-known fact \cite{kannan1995isoperimetric,bobkov1999isoperimetric} that any log-concave measure satisfies a PI with finite PI constant. Then per \cite[Remark 7.12]{klartag2025lecturenoteisoperimetric} (see also \cite{klartag2023logarithmic}), we know that $ C_{\mathrm{PI}}(\pi)\leq C'\log (d+1) \opnorm{\cov{\pi}}$
for a universal constant $C'$. Moreover, \cite[Lemma 2.5]{kook2025algodiffusion} says $\opnorm{\cov{\pi}}\leq 2\brac{\opnorm{\cov{\pi^X}}+160 }$ which leads to the desired estimate. 
\end{proof}

\textit{Remark:}
    A long line of works \cite{cheeger2015lower,kannan1995isoperimetric,lee2024eldan,chen2021almost,klartag2022bourgain,klartag2023logarithmic} establish $C_{\mathrm{PI}}(\nu)=  \mathcal{O}\brac{\opnorm{\cov {\nu}}\log d}$ for a general log-concave distribution $\nu$. This matches, up to constants, the order of the bound on $C_{\mathrm{PI}}(\pi)$ in Lemma~\ref{silif:lem:PIconstant}.




The next lemma shows that warm starts are preserved by the lifting and remain valid throughout the iterations of the proximal sampler~\eqref{silif:alg:ASF_pi}. 

\begin{lemma}\label{silif:lem:warmstart} 
An $M$-warm start $\nu_0$ for $\nu=\pi^X$ (condition~(A3)) induces an $M$-warm start $\pi_0$ for $\pi$, i.e., $d\pi_0/d\pi\le M$. Moreover, if $\pi^{Y,S}$ denotes the law of $(Y,S)$ after Step~1 of Algorithm~\ref{silif:alg:ASF_pi}, and $\pi_k$ denotes the law of $(x_k,t_k)$ in Algorithm~\ref{silif:alg:ASF_pi}, then $d\pi^{Y,S}/d\Pi^{Y,S}\le M$ and $d\pi_k/d\Pi^{X,T}\le M$, so warmness holds for every step of Algorithm~\ref{silif:alg:ASF_pi}.
\end{lemma}

% \red{could mention \cite{kook2025algodiffusion} here}

% for $t \in [f(x)/a, \infty]$
% \[
% F(t)=\frac{\int_{-\infty}^t e^{-at}1_Qdt}{\int_{-\infty}^{\infty} e^{-at}1_Qdt}=\frac{1_{x\in K}\frac{1}{a}\brac{e^{-f(x)} - e^{-at}}}{1_{x\in K}\frac{1}{a}\brac{e^{-f(x)}}} = \frac{1_{x\in K}\brac{1- e^{f(x)-at}}}{1_{x\in K}}
% \]
% then $t=F^{-1}(u)$. 

% In particular, $x\sim X\sim \nu$ so $x\in K$ almost surely, and we can say the CDF of $\pi^{T|X=x}$ equals $1- e^{f(x)-at}$ almost surely. 




\begin{proof}    For the first part of the lemma, we recall the idea in \cite[Theorem 2.15]{kook2025algodiffusion}: generate $x\sim \nu_0$ and $t\sim \pi^{T|X=x}\propto e^{-at}\mathbf{1}_Q$. 
\begin{itemize}
    \item In particular, one way to construct an initialization $\nu_0$ that is $M$-warm with respect to $\nu$ in~\eqref{silif:def:nu} is to use the Gaussian cooling technique of~\cite[Section~3]{kook2025algodiffusion}, which extends the framework of~\cite{cousins2018gaussian} from generating warm starts for the uniform distribution to generating warm starts for general log-concave targets. We note that by condition~(A3), $\nu_0\ll \nu$. Since $\nu$ is supported on $K$, it follows that $\nu_0(\R^d\setminus K)=0$, and hence $x\in K$ almost surely under $x\sim \nu_0$. 
    \item Meanwhile, $\pi^{T|X=x}$ is a one-dimensional distribution and can be generated by sampling $u\sim U[0,1]$ and let $t=F^{-1}(u)$ where $F(t)$ is the distribution of $\pi^{T=t|X=x}$ and equals $1- e^{f(x)-at}$ almost surely with respect to $x$.
\end{itemize}
Denoting the law of $(x,t)$ by $\pi_0$, then $\pi_0$ is $M$-warm with respect to $\pi$ since 
\[
\frac{\D \pi_0}{\D\pi}(x,t)
=\frac{\left(\frac{\D \nu_0}{\D x}\right)(x)\,\pi^{T\mid X=x}(t)}
{\left(\frac{\D \nu}{\D x}\right)(x)\,\pi^{T\mid X=x}(t)}
=\frac{\D\nu_0}{\D\nu}(x)
\le M.\]
    For the second part of the lemma, assume any $U\subseteq \R^{d+1}$. For $(y,s)\in \R^{d+1}$, set $U-(y,s)=\{(x,t)\in \R^{d+1}:(x,t)+(y,s)\in U \}$. Denote $\frak{g}(\cdot)$ the density of $\mathcal{N}(0,\eta I_{d+1})$. Then, using $\D \pi_0/\D \pi=\D \pi_0/\D \Pi^{X,T}\leq M$, $\pi^{Y,S} = \pi_0 * \frak{g}$, and $\Pi^{Y,S} = \Pi^{X,T} * \frak{g}$, we have
\begin{align*}
    \pi^{Y,S}(U)= \int_{\R^{d+1}} \pi_0(U-y) \frak{g}(y,s)\D y\D s&=\int_{\R^{d+1}} \left( \int_{U-(y,s)} \frac{\D \pi_0}{\D \Pi^{X,T}}(x,t)  \Pi^{X,T}(\D x,\D t) \right) \frak{g}(y,s) \D y\D s\\
    &\leq M\int_{\R^{d+1}} \Pi^{X,T}(U-(y,s))\frak{g}(y,s) \D y\D s = M\Pi^{Y,S}(U).
\end{align*}
A similar argument will give $\frac{\D\pi_k}{\D\Pi^{X,T}}\leq M$. Thus, we can conclude the warmness holds for every step of Algorithm \ref{silif:alg:ASF_pi}.
\end{proof}




In the upcoming result, we ensure the acceptance test at \eqref{silif:event:separation} is well-defined.
\begin{lemma}\label{silif:lem:compareThetaP1}  
Recall $z=(y,s-a\eta)$ defined in~\eqref{silif:def:notation}, $\tilde{w}$ defined in Algorithm~\ref{silif:RGO}, and the function $\mathcal{P}_1$ in~\eqref{silif:def:P1}. Assuming~\eqref{silif:difference:tildextildet:mainpaper} holds, then we have for every $w\in \R^{d+1}$,
\[
    \mathcal{P}_1(w)\leq \Theta_{y,s}^{\eta,Q}(w),
\]
and hence the acceptance test \eqref{silif:event:separation} is well-defined.

\end{lemma}
\begin{proof}
Set 
\begin{equation}
\label{silif:def:P0}
    \mathcal{P}_0(w)=\frac{1}{2\eta}\brac{\norm{\proj_Q(z)-z}^2+\norm{w - \proj_Q(z)}^2}  - \frac{a^2\eta}{2} + a s. 
\end{equation}
Then it follows from \eqref{silif:def:Theta} and \eqref{silif:optimizingThetaasprojection} that
    \begin{align}
    \label{silif:ineq:compareThetaP1}
    \Theta^{\eta,Q}_{y,s}(w) &\stackrel{\eqref{silif:def:Theta}}=I_Q(w) +\frac{1}{2\eta}\norm{w-z}^2- \frac{a\eta^2}{2} + a s\nn\\
        & \stackrel{\eqref{silif:optimizingThetaasprojection}}\geq \frac{1}{2\eta}\norm{\proj_Q(z)-z}^2 +\frac{1}{2\eta}\norm{w - \proj_Q(z)}^2- \frac{a\eta^2}{2} + a s\stackrel{\eqref{silif:def:P0}}=\mathcal{P}_0(w).
    \end{align}
It follows from the triangle inequality and \eqref{silif:difference:tildextildet:mainpaper} that for every $w\in \R^{d+1}$, 
\begin{align}
\label{silif:estimatewihoutsquare}
    &\norm{w-\tilde{w}}\leq \norm{w-\proj_Q(z)}+\norm{\proj_Q(z)-\tilde{w}} \stackrel{\eqref{silif:difference:tildextildet:mainpaper}}\leq \norm{w-\proj_Q(z)}+\brac{1+\frac{L}{a}}\sqrt{\frac{2\eta}{d+1}}
\end{align}
    Taking the square, and applying the triangle inequality, we have
    \begin{align*}
        &\norm{w-\tilde{w}}^2\\
        &\leq \norm{w-\proj_Q(z)}^2+2\norm{w-\proj_Q(z)}\brac{1+\frac{L}{a}}\sqrt{\frac{2\eta}{d+1}}+\brac{1+\frac{L}{a}}^2\frac{2\eta}{d+1}\\
        &\leq \norm{w-\proj_Q(z)}^2+2\brac{\norm{w-\tilde{w}} +\norm{\tilde{w}- \proj_Q(z)}}\brac{1+\frac{L}{a}}\sqrt{\frac{2\eta}{d+1}}+\brac{1+\frac{L}{a}}^2\frac{2\eta}{d+1}\\
        &\stackrel{\eqref{silif:difference:tildextildet:mainpaper}}{\leq} \norm{w-\proj_Q(z)}^2+2\norm{w-\tilde{w}}\brac{1+\frac{L}{a}}\sqrt{\frac{2\eta}{d+1}}+\brac{1+\frac{L}{a}}^2\frac{6\eta}{d+1}. 
    \end{align*}
This can be rearranged as 
\begin{align}
\label{silif:firstestimate}
    \norm{w-\tilde{w}}^2-2\norm{w-\tilde{w}}\brac{1+\frac{L}{a}}\sqrt{\frac{2\eta}{d+1}}-\brac{1+\frac{L}{a}}^2\frac{6\eta}{d+1}\leq \norm{w-\proj_Q(z)}^2. 
\end{align}
Taking $w=z$ in \eqref{silif:firstestimate}, we obtain
\begin{align}
    \label{silif:secondestimate}
      &\norm{z-\tilde{w}}^2-2\norm{z-\tilde{w}}\brac{1+\frac{L}{a}}\sqrt{\frac{2\eta}{d+1}}-\brac{1+\frac{L}{a}}^2\frac{6\eta}{d+1}\leq \norm{z-\proj_Q(z)}^2. 
\end{align}
In view of \eqref{silif:def:P1} and \eqref{silif:def:P0}, combining \eqref{silif:firstestimate} and \eqref{silif:secondestimate} leads to $\mathcal{P}_0(w)\geq \mathcal{P}_1(w),\forall w$. 
It further follows from \eqref{silif:ineq:compareThetaP1} that
\[
    \mathcal{P}_1(w)\leq \mathcal{P}_0(w)\leq \Theta_{y,s}^{\eta,Q}(w), \forall w. 
\]
This completes the proof. 
\end{proof}


In the upcoming lemma, we construct an auxiliary function  $\mathcal{P}_2$ that will be useful in the rejection analysis of Algorithm~\ref{silif:RGO}. 
\begin{lemma}
\label{silif:lem:compareP1P2}  
 Recall $z=(y,s-a\eta)$ defined in~\eqref{silif:def:notation} and $\tilde{w}$ defined in Algorithm~\ref{silif:RGO}. For $w\in \R^{d+1}$, define the real-valued function 
\begin{align}
\label{silif:def:P2}
\mathcal{P}_2(w)
&= \frac{1}{2\eta}\Bigg[
\left(\|w-\tilde{w}\| - \left(1+\frac{L}{a}\right)\sqrt{\frac{2\eta}{d+1}}\right)^2 \nonumber\\
&
+ \left(\|z-\operatorname{proj}_Q(z)\| - 2\left(1+\frac{L}{a}\right)\sqrt{\frac{2\eta}{d+1}}\right)^2
- 32\left(1+\frac{L}{a}\right)^2\frac{\eta}{d+1}
\Bigg] -\frac{a^2\eta}{2} + a s.
\end{align}



Then, recall $\mathcal{P}_1$ at \eqref{silif:def:P1} and assume ~\eqref{silif:difference:tildextildet:mainpaper} holds, then
\begin{equation}\label{silif:ineq:P2-P3}
    \mathcal{P}_1(w)\geq \mathcal{P}_2(w), \quad \forall w\in \R^{d+1}.
\end{equation}
\end{lemma}

\begin{proof}
Applying~\eqref{silif:estimatewihoutsquare} with $w=z$ to get
\begin{equation}
\label{silif:estimatewihoutsquareP3}
    -\norm{z-\proj_Q\brac{z}}-\brac{1+\frac{L}{a}}\sqrt{\frac{2\eta}{d+1}}\leq -\norm{z-\tilde{w}}
\end{equation}
It follows from the triangle inequality and \eqref{silif:difference:tildextildet:mainpaper} that
\[
\norm{z-\proj_Q(z)}\leq \norm{\tilde{w} -z} +\norm{\tilde{w}-\proj_Q(z)}\stackrel{\eqref{silif:difference:tildextildet:mainpaper}}{\leq} \norm{\tilde{w} -z} +\brac{1+\frac{L}{a}}\sqrt{\frac{2\eta}{d+1}},
\]
and hence that
\begin{align}
\label{silif:firstestimateP3}
    &\norm{z-\proj_Q(z)}^2-2\norm{z-\proj_Q(z)}\brac{1+\frac{L}{a}}\sqrt{\frac{2\eta}{d+1}}-\brac{1+\frac{L}{a}}^2\frac{6\eta}{d+1}\leq \norm{\tilde{w}-z}^2. 
\end{align}
Using \eqref{silif:estimatewihoutsquareP3}, \eqref{silif:firstestimateP3}, and the definition of $\mathcal{P}_1$ in \eqref{silif:def:P1}, we have
\begin{align*}
    &\mathcal{P}_1(w)\\
    &\geq \frac{1}{2\eta}\Bigg[\norm{w - \tilde{w}}^2 + \norm{z-\proj_Q(z)}^2-2\norm{z-\proj_Q(z)}\brac{1+\frac{L}{a}}\sqrt{\frac{2\eta}{d+1}} -\brac{1+\frac{L}{a}}^2\frac{6\eta}{d+1} \\
    &-2\sqrt{\frac{2\eta}{d+1}}\brac{1+\frac{L}{a}} \bigg(\norm{w-\tilde{w}} + \norm{z-\proj_Q\brac{z}}+\brac{1+\frac{L}{a}}\sqrt{\frac{2\eta}{d+1}}\bigg) -\brac{1+\frac{L}{a}}^2\frac{12\eta}{d+1}\Bigg] \\
    & - \frac{a^2\eta}{2} + a s,
\end{align*}
and the right hand side simplifies to the stated formula for $\mathcal{P}_2$. 
\end{proof}


\subsection{Results about the cutting-plane method by \cite{jiang2020cuttingplane}}
\label{silif:appen:cp}

We first restate \cite[Theorem~C.1]{jiang2020cuttingplane}, which is about the iteration complexity and running time of the CP method by \cite{jiang2020cuttingplane}. 

\begin{theorem}(\cite[Theorem~C.1]{jiang2020cuttingplane})\label{theo:jiang} Let $f$ be a convex function on $\R^d$. $K$ is a convex set that contains a minimizer
of $f$ and $K\subseteq B_\infty (0,R)$, where $B_\infty (0,R)$ denotes a ball of radius $R$ in $\ell_\infty$ norm, i.e., $\norm{x}_\infty=\sup_{1\leq i\leq d}\abs{x_i}$. 

Suppose we have a subgradient oracle for $f$ with cost $T$ and a separation oracle for $K$ with cost $S$. Using $B_\infty (0,R)$ as the initial polytope for our CP Method, for any $0< \alpha<1$, we can compute $\tilde x\in K$ such that 
\begin{equation}\label{ineq:opt}
f(\tilde x)-\min_{x \in K} f(x) \le \alpha \left(\max_{x \in K} f(x)-\min_{x \in K} f(x)\right).    
\end{equation}
with a running time of $ {\cal O}\left(T\cdot d \log \frac{d \gamma}{\alpha}+S\cdot  d \log \frac{d \gamma}{\alpha}+d^3\log \frac{d \gamma}{\alpha}\right)$. In particular, the number of subgradient oracle calls and the number of separation oracle calls are of the order 
\begin{align*}
    \mathcal{O}\brac{d \log \frac{d \gamma}{\alpha}}, 
\end{align*}
where
\begin{align*}
    \gamma=R/\mbox{minwidth}(K), \quad \mbox{minwidth}(K)=\min_{\norm{a}=1}\left\{\max_{y\in K}a^Ty-\min_{y\in K}a^Ty\right\}. 
\end{align*}

\end{theorem}


In the upcoming result, we rewrite~\eqref{silif:optimizingThetaasprojection}, a strongly convex optimization problem over the unbounded set $Q$, as an equivalent strongly convex program over the compact set $K$.


\begin{lemma} 
\label{silif:lem:equivopt}
Assume conditions (A1)-(A2) and define  
\begin{equation}
\label{silif:def:zeta}
     \zeta^{\eta}_{y,s}(x)= \frac{1}{2\eta}\|x-y\|^2+ \frac{1}{2\eta}\left[ \frac{f(x)}{a} - (s - a\eta)\right]_{+}^2. 
\end{equation}
The following statements hold: 
\begin{itemize}
    \item[(a)] the solution to $\underset{(x,t)\in\R^{d+1}}{\min} \Theta^{\eta,Q}_{y,s}(x,t)$ in \eqref{silif:optimizingThetaasprojection} is 
\[
x_* =\underset{x\in\R^d}{\argmin} \zeta^{\eta}_{y,s}(x), \quad
t_*=t_*(x_*)= \max\left\{ \frac{f(x_*)}{a},~ s - a\eta \right\};
\]
    \item[(b)]  $\zeta^{\eta}_{y,s}(x)$ is $\eta^{-1}$-strongly convex;
    \item[(c)] if $\tilde{x}$ is a $(d+1)^{-1}$-solution to $\min_{x\in K} \zeta^{\eta}_{y,s}(x)$ and $\tilde{t}=\max\left\{ f(\tilde{x})/a, s - a\eta \right\}$, then $\tilde{w}=(\tilde{x},\tilde{t})$ satisfies \eqref{silif:difference:tildextildet:mainpaper}.
\end{itemize}
\end{lemma}

\begin{proof}  
~\

\textbf{Part~a:} Per the definitions of $Q$ and $\Theta^{\eta,Q}_{y,s}(x,t)$ at respectively~\eqref{silif:def:Q} and~\eqref{silif:def:Theta}, 
 \begin{align}
 \label{silif:Thetaexpand}
 \Theta^{\eta,Q}_{y,s}(x,t)
= I_{\{x\in K,f(x)\leq at\}}(x,t)+ \frac{1}{2\eta}\norm{x-y}^2+\frac{1}{2\eta}(t-(s-a\eta))^2+as-\frac{a^2\eta}{2}. 
 \end{align}
 Then regarding  $\underset{(x,t)\in\R^{d+1}}{\min} \Theta^{\eta,Q}_{y,s}(x,t)$, if we fix $x\in K$, the optimization over $t$ is 
\[
t^*= \underset{t \ge f(x)/a}{\argmin}\left\{
\frac{1}{2\eta}\left(t-(s - a\eta)\right)^2 \right\}
=\max\left\{ \frac{f(x)}{a}, s - a\eta \right\}.
\]
Then 
\begin{align}
\label{quadratictermt}
\frac{1}{2\eta}\left(t^*-(s - a\eta)\right)^2&=\begin{cases}
    \frac{1}{2\eta}\left[\frac{f(x)}{a}-(s - a\eta)\right]^2 &\text{if} \quad \frac{f(x)}{a}\geq s-a\eta,\\
    0 &\text{otherwise} 
\end{cases}\nn\\
&=\frac{1}{2\eta}\left[ \frac{f(x)}{a} - (s - a\eta)\right]_{+}^2. 
\end{align}
Plugging~\eqref{quadratictermt} into~\eqref{silif:Thetaexpand} yields the stated formula for $x_*$. Having $x_*$ also leads to the formula for $t_*$. 

\textbf{Part~b:} Since $f$ is convex, $a>0$, and the maps $t \mapsto \max\{t,0\}$ and $t \mapsto t^2, t\ge0$ are convex and non-decreasing, $\left[ \frac{f(x)}{a} - (s - a\eta)\right]_{+}^2$ is convex. The strong convexity of $\zeta^{\eta}_{y,s}(x)$ is due to the quadratic term $\frac{1}{2\eta}\|x-y\|^2$. 


\textbf{Part c:}  As we have shown that $\zeta^{\eta}_{y,s}$ is $\eta^{-1}$-strongly convex per Part~b, so
    \begin{align}
    \label{silif:difference:tildex}
        \norm{\tilde{x}-x_*}\leq \sqrt{2\eta\brac{\zeta^{\eta}_{y,s}(\tilde{x})-\zeta^{\eta}_{y,s}(x_*) }}\leq \sqrt{\frac{2\eta}{d+1}}. 
    \end{align}
Next, the fact that $f(x)$ is $L$-Lipschitz implies that $h(x)=\max \left\{f(x)/a,s-a\eta\right\}$ is $L/a$-Lipschitz. Combining this with \eqref{silif:difference:tildex} and $\tilde{t}=h(\tilde{x})$, we get 
\begin{align}
\label{silif:difference:tildet}
    \abs{\tilde{t}-t_*}=\abs{h(\tilde{x})-h(x_*) }\stackrel{\eqref{silif:difference:tildex}}{\leq} \frac{L}{a}\norm{\tilde{x}-x_*}\leq \frac{L}{a}\sqrt{\frac{2\eta}{d+1}}.  
\end{align}
We now set $\tilde{w}=(\tilde{x},\tilde{t})$. Using \eqref{silif:difference:tildex},\eqref{silif:difference:tildet} and the fact that $(x_*,t_*)=\proj_Q(z)$ per \eqref{silif:optimizingThetaasprojection}, we arrive at the desired bound on $\norm{\tilde{w}- \proj_Q(z)}$. 
\end{proof}

Next, we apply Theorem~\ref{theo:jiang} by \cite{jiang2020cuttingplane} to the function $\zeta^{\eta}_{y,s}$ defined in Lemma~\ref{silif:lem:equivopt}, with the goal of producing $\tilde{w}$ satisfying \eqref{silif:difference:tildextildet:mainpaper}. 

\begin{theorem}
\label{silif:theo:cuttingplane} 
Assume conditions (A1), (A2) and (A4). Recall $x_*=\argmin \zeta^{\eta}_{y,s}(x)$ in Lemma~\ref{silif:lem:equivopt}. Set
    \begin{align}
    \label{silif:def:alpha}
        \alpha=\min\left\{\frac{\eta}{(d+1)\brac{2R^2+2R\norm{x_*-y}+ \left[\frac{\max_{x\in K}f(x)}{a}-s+a\eta\right]_{+}\frac{RL}{a} }},\frac{1}{2}\right\}. 
    \end{align}
The following statements hold. 
\begin{itemize}
    \item[(a)] The CP method by \cite{jiang2020cuttingplane} makes ${\cal O}\left(d \log \frac{d \gamma}{\alpha}\right)$ calls to the separation oracle of $K$ and ${\cal O}\left(d \log \frac{d \gamma}{\alpha}\right)$ calls to the subgradient oracle of $f$ in order to generate a $(d+1)^{-1}$-solution $\tilde{x}\in K$ to the optimization problem $\min_{x\in K} \zeta^{\eta}_{y,s}(x)$. Here $\gamma=R/\mbox{minwidth}(K)$ and $\mbox{minwidth}(K)$ are defined in~Section~\ref{sec:prelim}. 
    \item[(b)] Under the additional assumption that $a=d$ and $\eta=1/d^2$,  we have  the following concentration inequality for $\alpha$ in Part~a:
 \begin{align*}
     \Pr\brac{\alpha\leq \frac{\min\left\{\frac{1}{4R(1+3R)}, \frac{1}{12R^2 \max\{L,1\}\max\{\norm{f}_\infty,1\} }\right\} }{(d+1)^3}}\leq 6\exp\brac{-\frac{d^2}{8}}. 
 \end{align*}
\end{itemize}
\end{theorem}

\begin{proof} 
~\

\textbf{Part~a:}   Let $x^*$ and $x_*$ be respectively the maximizer and minimizer over $K$ of  the objective $\zeta^{\eta}_{y,s}$ in Lemma~\ref{silif:lem:equivopt}. Then we have  
    \begin{align}
    \label{silif:eq:firstboundmaxandmin}
        &\zeta^{\eta}_{y,s}(x^*)- \zeta^{\eta}_{y,s}(x_*) \stackrel{\eqref{silif:def:zeta}} = \frac{1}{2\eta}\brac{\norm{x^*-y}^2-\norm{x_*-y}^2 }\nn\\
        &+\frac{1}{2\eta}\brac{\left[ \frac{f(x^*)}{a} - (s - a\eta)\right]_{+}^2-\left[ \frac{f(x_*)}{a} - (s - a\eta)\right]_{+}^2}. 
    \end{align}
By the triangle inequality and the fact that $x^*,x_*\in K\subseteq R\ball_d(0)$ (condition~(A1)), 
\begin{align}
\label{silif:zetafirstpart}
   \frac{1}{2\eta}\brac{\norm{x^*-y}^2-\norm{x_*-y}^2 }&\leq \frac{1}{2\eta}\brac{\brac{\norm{x^*-x_*}+\norm{x_*-y}}^2-\norm{x_*-y}^2 }\nonumber\\
    &\leq \frac{1}{2\eta}\brac{4R^2+4R\norm{x_*-y} }. 
\end{align}
Moreover, since $f(x)$ is $L$-Lipschitz, the function $\left[\frac{f(x)}{a}-s+a\eta\right]_{+}$ is $L/a$-Lipschitz, and hence 
\begin{align}
\label{silif:zetasecondpart}
    &\frac{1}{2\eta}\brac{\left[ \frac{f(x^*)}{a} - (s - a\eta)\right]_{+}^2-\left[ \frac{f(x_*)}{a} - (s - a\eta)\right]_{+}^2}\nonumber\\
    &\leq \frac{2}{2\eta}\max_{x\in K} \left\{\left[ \frac{f(x)}{a} - (s - a\eta)\right]_{+}\right\}\norm{x^*-x_*}\frac{L}{a}\nonumber\\
    &\leq\frac{1}{\eta}  \left[\frac{\max_{x\in K}f(x)}{a}-s+a\eta\right]_{+}\frac{RL}{a}. 
\end{align}
Combining \eqref{silif:eq:firstboundmaxandmin},\eqref{silif:zetafirstpart} and \eqref{silif:zetasecondpart} to get
\begin{align*}
     \zeta^{\eta}_{y,s}(x^*)- \zeta^{\eta}_{y,s}(x_*)&\leq \frac{2R}{\eta}\brac{R+\norm{x_*-y} }+\frac{1}{\eta}  \left[\frac{\max_{x\in K}f(x)}{a}-s+a\eta\right]_{+}\frac{RL}{a}. 
\end{align*}
In view of the above equation, \eqref{ineq:opt} from Theorem~\ref{theo:jiang} and the choice of $\alpha\in (0,1)$ at \eqref{silif:def:alpha}, Theorem~\ref{theo:jiang} guarantees the CP method by~\cite{jiang2020cuttingplane} produces a $(d+1)^{-1}$-solution $\tilde{x}$  to  $\underset{x\in\R^d}{\argmin} \zeta^{\eta}_{y,s}(x)$, i.e.,
\[
\zeta^{\eta}_{y,s}(\tilde x)-\min_{x \in K} \zeta^{\eta}_{y,s}(x)\stackrel{\eqref{ineq:opt}}{\leq}    \alpha\brac{ \zeta^{\eta}_{y,s}(x^*)- \zeta^{\eta}_{y,s}(x_*)}\stackrel{\eqref{silif:def:alpha}\eqref{silif:zetafirstpart}, \eqref{silif:zetasecondpart}}{\leq} \frac{1}{d+1}. 
\]
Therefore, Theorem~\ref{theo:jiang} shows that the CP method by~\cite{jiang2020cuttingplane} requires  ${\cal O}\left(d \log \frac{d \gamma}{\alpha}\right)$ separation oracle calls and subgradient oracle calls. This completes the first part of the proof. 

\textbf{Part~b:} In this part, we set $a=d$, $\eta=1/d^2$ and derive the concentration inequality for $\alpha$. Recall from Section~\ref{sec:constrained:mainpaper}, there is a positive constant $\norm{f}_\infty$ such that $\abs{f(x)}\leq \norm{f}_\infty,\forall x\in K$. Let us define
  \begin{align}
  \label{silif:def:n*}
 n_*=n_1\wedge n_2, \quad\text{where}\quad      n_1=\frac{1}{(d+1)^3}\frac{1}{4R(1+3R)}, n_2=\frac{1}{(d+1)^3}\frac{1}{12R^2 \max\{L,1\}\max\{\norm{f}_\infty,1\} }. 
  \end{align}
Now let $n>0$ be a generic positive constant to be determined later. Then
\begin{align}
\label{silif:splitT1T2}
    \Pr\brac{\alpha\leq n}
    &\leq \Pr\brac{\frac{2\eta}{(d+1)\brac{4R^2+4R\norm{x_*-y}+2 \left[\frac{\max_{x\in K}f(x)}{a}-s+a\eta\right]_{+}\frac{RL}{a} }}\leq n}\nonumber\\
    &\leq \Pr\brac{4R\norm{x_*-y}+2\left[\frac{\max_{x\in K}f(x)}{a}-s+a\eta\right]_{+}\frac{RL}{a}\geq \frac{2}{n d^2(d+1)}-4R^2 }\nonumber\\
    &\leq \Pr\brac{4R\norm{x_*-y}\geq  \frac{1}{n d^2(d+1)}-2R^2}\nonumber\\
    &+\Pr\brac{ 2\left[\frac{\max_{x\in K}f(x)}{a}-s+a\eta\right]_{+}\frac{RL}{a}\geq \frac{1}{n d^2(d+1)}-2R^2}:=T_1(n)+T_2(n). 
\end{align}
We know that $(y,s)$ is the output of step 1 in the proximal sampler (Algorithm \ref{silif:alg:ASF_pi}), which means
\begin{align}
    \label{silif:outputstep1}
    y=x_{k-1}+\frac{1}{d}Z_1,\quad  s=t_{k-1}+\frac{1}{d}Z_2
\end{align}
where $Z_1\sim \mathcal{N}(0,I_d)$ and $Z_2\sim \mathcal{N}(0,1)$. 

Then regarding $T_1$ on the right hand side of \eqref{silif:splitT1T2}, we have $\norm{x_*-y}\leq \norm{x_*-x_{k-1}-\frac{1}{d}Z_1}\leq 2R+\frac{1}{d}\norm{Z_1}$, and thus
\begin{align*}
    T_1(n)&=\Pr\brac{4R\norm{x_*-y}\geq  \frac{1}{n d^2(d+1)}-2R^2}\leq \Pr\brac{\norm{Z_1}\geq d\brac{\frac{1}{nd^2(d+1)4R}-R-2R }}. 
\end{align*}
  Plugging in $n=n_*$ at \eqref{silif:def:n*} and using $n_1\leq n_*$ to get 
  \begin{align*}
      T_1(n_*)&\leq \Pr\brac{\norm{Z_1}\geq d\brac{\frac{1}{n_*d^2(d+1)4R}-3R }}\\
      &\leq \Pr\brac{\norm{Z_1}\geq d\brac{\frac{1}{n_1 d^2(d+1)4R}-3R }}. 
  \end{align*}
Further note that $ \frac{1}{n_1 d^2(d+1)4R}-3R\geq \frac{1}{n_1 (d+1)^3 4R}-3R \geq 1$,
which implies 
\begin{align}
\label{silif:boundT1}
     T_1(n_*)&\leq \Pr\brac{\norm{Z_1}\geq d}\leq 4\exp\brac{-\frac{d^2}{8}} 
\end{align}
per the Gaussian concentration inequality from \cite[Equation (3.5)]{ledoux2013probability}. 

Regarding $T_2(n)$ on the right hand side of \eqref{silif:splitT1T2}, recall $\max_{x\in K}f(x)\leq \norm{f}_\infty$. We also know $(x_{k-1},t_{k-1})\in Q$ almost surely (see the remark following Algorithm \ref{silif:RGO}), so that $-t_{k-1}\leq -\frac{f(x_{k-1})}{d}\leq \frac{\norm{f}_\infty}{d}$ almost surely. Then
\begin{align}
\label{silif:T2middle}
    T_2(n)&\leq \Pr\brac{ 2\brac{\frac{\max_{x\in K}f(x)}{a}-s+a\eta}\frac{RL}{a}\geq \frac{1}{n d^2(d+1)}-2R^2}\nonumber\\
    &\stackrel{\eqref{silif:outputstep1}}{\leq} \Pr\brac{Z_2\leq \frac{\norm{f}_\infty}{d}-t_{k-1}+\frac{1}{d}-\frac{d}{2RL}\brac{\frac{1}{nd^2(d+1)}-2R^2} }\nonumber\\
    &\leq \Pr\brac{Z_2\leq \frac{2\norm{f}_\infty}{d}+1-\frac{d}{2RL}\brac{\frac{1}{nd^2(d+1)}-2R^2} }.
\end{align}
Set 
\begin{align*}
    n_3=\brac{\brac{\frac{2RL}{d}\brac{\frac{2\norm{f}_\infty}{d}+1+d}+2R^2 }d^2(d+1)}^{-1}
\end{align*}
then \begin{align*}
    n_3&=\brac{4RL(d+1)+2RL(1+d)^2d+2R^2 d^2(d+1) }^{-1}\\
    &\geq \brac{3\max\{4,2,2\}\max\{R,R^2\}\max\{L,1\}\max\{\norm{f}_\infty,1\}\max\{d+1,(1+d)^2d,d^2(d+1)\} }^{-1}\\
    &\geq \frac{1}{(d+1)^3}\frac{1}{12R^2 \max\{L,1\}\max\{\norm{f}_\infty,1\} }=n_2.
\end{align*}
Therefore, we have $n_*\leq n_2\leq n_3$. Now plugging $n_*$ into \eqref{silif:T2middle} to get
\begin{align}
\label{silif:boundT2}
    T_2(n_*)&\leq \Pr\brac{Z_2\leq \frac{2\norm{f}_\infty}{d}+1-\frac{d}{2RL}\brac{\frac{1}{n_*d^2(d+1)}-2R^2} }\nonumber\\
   & \leq \Pr\brac{Z_2\leq \frac{2\norm{f}_\infty}{d}+1-\frac{d}{2RL}\brac{\frac{1}{n_3d^2(d+1)}-2R^2} }\nonumber\\
   &=\Pr\brac{ Z_2\leq -d}\leq 2\exp\brac{-\frac{d^2}{8}}.
\end{align}
The last line is again due to \cite[Equation (3.5)]{ledoux2013probability}.  

Combining \eqref{silif:splitT1T2}, \eqref{silif:boundT1} and \eqref{silif:boundT2} to get $ \Pr\brac{\alpha\leq n_*}\leq 6\exp\brac{-\frac{d^2}{8}}$ which is the desired bound. 
\end{proof}


\subsection{Sampling $w\sim \exp\brac{-\mathcal{P}_1(w)}$ in Algorithm~\ref{silif:RGO}}
\label{silif:appen:specialalgo}   

Recall the definition of $\mathcal{P}_1$ at~\eqref{silif:def:P1} and the point $\tilde{w}$ at~\eqref{silif:difference:tildextildet:mainpaper}. By completing the square, one can easily see that sampling $w\sim \exp\brac{-\mathcal{P}_1(w)}$ is equivalent to
\begin{align*}
 w\sim \Lambda(w)\propto \exp\brac{-\frac{1}{2\eta}\brac{\norm{w-\tilde{w} }-\sqrt{\frac{2\eta}{d+1} }\brac{1+\frac{L}{a}} }^2}.
\end{align*}



While $\Lambda(w)$ is not a Gaussian density, generating $w\sim \Lambda(w)$ is straightforward since it can be turned into a one-dimensional sampling problem. We state here a generic procedure for this sampling problem. An explanation is given in Lemma~\ref{silif:lem:samplingP2} below. 
\begin{algorithm}[H]
	\caption{Sample $w \sim \Lambda(w)$} 
	\label{silif:alg:samplingP1}
	\begin{algorithmic}
		\State 1. Generate $W\sim {\cal N}(0,I_{d+1})$ and set $\theta = W/\|W\|$;
        \State 2. Generate $r \propto r^{d} \exp\left(-\frac{1}{2\eta} \brac{r-\sqrt{\frac{2\eta}{d+1} }\brac{1+\frac{L}{a}} }^2\right)$ by Adaptive Rejection Sampling for one-dimensional log-concave distribution by \cite{gilks1992adaptive}.
        \State 3. Output $w = \tilde{w} + r \theta$.
	\end{algorithmic}
\end{algorithm}

\begin{lemma}
\label{silif:lem:samplingP2}
    Algorithm \ref{silif:alg:samplingP1} generates $w \sim \Lambda(w)$.
\end{lemma}

\begin{proof}  
We rewrite $\Lambda(w)$ in polar coordinates.
Let $r = \| w - \widetilde{w} \|$. 
Since $\mathrm{d}x \mathrm{d}t = r^{d} \mathrm{d}r \mathrm{d}S(\theta)$,
where $\mathrm{d}S$ is surface measure on $\mathbb S^{d}$, for $r \ge 0$ and $\theta \in \mathbb{S}^{d}$ we have
\[
 q(w) = p(r,\theta) \propto r^{d} \exp \left(-\frac{1}{2\eta} \left(r - \sqrt{\frac{2\eta}{d+1}}  \left(1 + \frac{L}{a}\right)\right)^2\right).
\]
Its first marginal is 
\[
p_r(r) \propto r^{d} \exp \left(-\frac{\left(r - \sqrt{\frac{2\eta}{d+1}}  \left(1 + \frac{L}{a}\right)\right)^2} { 2\eta}\right).
\]
Since 
\[
\log p_r(r) = d \log r - \frac{\left(r - \sqrt{\frac{2\eta}{d+1}}  \left(1 + \frac{L}{a}\right)\right)^2 }{2\eta} + \text{const}, 
\quad 
\frac{\mathrm{d}^2}{\mathrm{d}r^2}\log p_r(r) = -\frac{d}{r^2} - \frac{1}{\eta} < 0,
\]
$p_r(r)$ is a one-dimensional log-concave distribution.
Thus, we can sample $r \sim p_r$ using any standard one-dimensional log-concave sampler (e.g., \cite{gilks1992adaptive}).
Then sample $W \sim \mathcal{N}(0, I_{d+1})$, set $\theta = W / \|W\|$ \cite{muller1959note}, and output 
\[
w = \tilde{w} + r \theta
\]
as the sample from $p(r,\theta) = \Lambda(w)$.
\end{proof}


\section{Supporting Results and Proofs for Section~\ref{sec:composite:mainpaper}}


\subsection{Full details of Section~\ref{sec:composite:mainpaper}}
\label{dolif:appen:fulldetails}

\subsubsection{Proximal sampler for the composite sampling problem}
\label{dolif:sec:ASF}

We follow the idea of the proximal sampler for the target $\tilde{\pi}$ at \eqref{dolif:def:piandQ} and start with the augmented distribution
\begin{equation}
\label{dolif:def:Pi:appen}
    \widetilde{\Pi}\brac{(x,s,t),(y,u,v)}\propto \exp\brac{-bt-I_{\Q}(x,s,t)-\frac{1}{2\eta}\norm{(x,s,t)-(y,u,v) }^2}.
\end{equation}
For a fixed point $(y,u,v)\in \R^{d+2}$,  set
\begin{align}
\label{dolif:def:Theta:appen}
    \widetilde{\Theta}^{\eta,\Q}_{y,u,v}(x,s,t)&=I_{\Q}(x,s,t)+bt+\frac{1}{2\eta}\norm{(x,s,t)-(y,u,v) }^2\nn\\
    &=I_{\Q}(x,s,t)+\frac{1}{2\eta}\|(x,s,t)-(y,u,v-\eta b)\|^2
+bv-\frac{\eta b^2}{2}. 
\end{align}




Then $\widetilde{\Pi}^{Y,U,V|X,S,T}(y,u,v|x,s,t)=\mathcal{N}\brac{p,\eta I_{d+2}}$ and $ \widetilde{\Pi}^{X,S,T|Y,U,V}(x,s,t|y,u,v)=\mathcal{N}\brac{q,\eta I_{d+2}}|_{\Q}$. The proximal sampler for $\tilde{\pi}$ is given below.


\begin{algorithm}[H]
	\caption{Proximal sampler for the target $\tilde{\pi}$ at \eqref{dolif:def:piandQ}}
	\label{dolif:alg:ASF_pi}
	\begin{algorithmic}[1]
		\State \textbf{Gaussian}: Generate $(y_k,u_k,v_k)\sim  \mathcal{N}\brac{(x_k,s_k,t_k),\eta I_{d+2}}$; 
		\State \textbf{RGO}: Generate $(x_{k+1},s_{k+1},t_{k+1})\sim {\cal N}\brac{(y_k,u_k,v_k-b\eta),\eta I_{d+2}}|_{\Q}.$ 
	\end{algorithmic}
\end{algorithm}




 We initialize Algorithm~\ref{dolif:alg:ASF_pi} by drawing $x_0\sim \tilde{\nu}_0$, an $M$-warm start for $\tilde{\nu}$ in~\eqref{dolif:def:nu}. 
Given $x_0$, we draw $s_0$ from the one-dimensional law $e^{-as}\mathbf{1}_{\{s\ge h(x_0)/a\}}$, which can be sampled by inverse transform using the CDF $F_S(s)=1-e^{h(x_0)-as}$ for $s\ge h(x_0)/a$, so that $s_0\sim \tilde{\pi}^{S |  X=x_0}$. 
Given $(x_0,s_0)$, we then draw $t_0$ from $e^{-bt}\mathbf{1}_{\{t\ge (f(x_0)+as_0)/b\}}$, again via inverse transform with CDF $F_T(t)=1-e^{f(x_0)+as_0-bt}$ for $t\ge (f(x_0)+as_0)/b$, hence $t_0\sim \tilde{\pi}^{T |  X=x_0,S=s_0}$. 
In particular, $(x_0,s_0,t_0)\sim \tilde{\pi}_0(x,s,t)=\tilde{\nu}_0(x) \tilde{\pi}^{S |  X=x}(s) \tilde{\pi}^{T |  X=x,S=s}(t)$ and $(x_0,s_0,t_0)\in \Q$. 
Thereafter, Algorithm~\ref{dolif:alg:ASF_pi} generates $(x_k,s_k,t_k)\sim \tilde{\pi}^k$ and maintains feasibility $(x_k,s_k,t_k)\in \Q$ at each iteration. Finally, 
Lemma~\ref{dolif:lem:warmstart} shows that $\tilde{\pi}_0$ is $M$-warm with respect to $\tilde{\pi}$ in~\eqref{dolif:def:piandQ}, and that $M$-warmness is retained for $\tilde{\pi}^k$ for every $k\ge 1$.

At this point, we refer to Theorem~\ref{dolif:theo:outer} in Section~\ref{sec:composite:mainpaper} for the iteration complexity of Algorithm~\ref{dolif:alg:ASF_pi} with exact RGO.  



%%%%%%%%%%%%%%%%%%%%%%%%%%%%%%%%%%%%%%%%%%%%%%%%%%%%%%%%%%%%%%%%%%%%%

\subsubsection{RGO implementation}

In Theorem~\ref{dolif:theo:outer}, the analysis of Algorithm~\ref{dolif:alg:ASF_pi} treats Step~2 as an exact draw from a Gaussian distribution truncated to $\Q$ at~\eqref{dolif:def:piandQ}. The goal of this section is to provide a realization of this truncated-Gaussian sampler via rejection sampling and to quantify the resulting acceptance probabilities and oracle usage. Similar to the explanation given in the beginning of Section~\ref{sec:constrained:mainpaper}, we need to approximately solve for $\argmin_{(x,s,t)\in\R^{d+2}}\widetilde{\Theta}^{\eta,\Q}_{y,u,v}(x,s,t)
=\proj_{\Q}( q)$, where $\widetilde{\Theta}^{\eta,\Q}_{y,u,v}$ is defined at~\eqref{dolif:def:Theta:appen}, in order to center our proposal in the rejection sampler. To do so, we use the CP method by~\cite{jiang2020cuttingplane}. Since this method requires an initial enclosing polytope, we restrict the search to a local bounded region inside a ball that still contains the unique minimizer of $\widetilde{\Theta}^{\eta,\Q}_{y,u,v}$. In particular, the reformulation in Lemma~\ref{silif:lem:equivopt} for a bounded convex body $K$ does not extend to the composite setting since $f$ and $h$ are defined on the unbounded domain $\R^d$.


Let the current input of Algorithm~\ref{dolif:alg:ASF_pi} be $(y,u,v)$ with $q=(y,u,v-b\eta)$. Let the previous RGO output be $ p_{k-1}=(x_{k-1},s_{k-1},t_{k-1})\in\Q$ a.s., set $r_{\local}:=\| p_{k-1}- q\|$ and define the local ball and local region as
\begin{equation}
\label{dolif:def:localball:appen}
    \ball_{\local}:=2r_{\local}\ball_{d+2}(q), \qquad  \Q_{\local}:=\Q\cap \ball_{\local}.
\end{equation}
First, to ensure that restricting the CP method to $\Q_{\local}$ is valid, we verify in Lemma~\ref{dolif:lem:minimizerinlocalball} in Appendix~\ref{dolif:appen:supportlem} that the minimizer $ p_*=\argmin_{ p\in \R^{d+2}}\widetilde{\Theta}^{\eta,\Q}_{y,u,v}( p)$
belongs to $\Q_{\local}$ a.s. Second, we have already identified in Section~\ref{sec:composite:mainpaper} oracle assumptions on $f$ and $h$ that yield a separation oracle for $\Q$, and hence the cutting planes required by the CP method on $\Q_{\mathrm{loc}}$. 



At this point, we can solve for a $(d+2)^{-1}$-solution $\tilde{ p}$ of $\argmin_{ p\in \Q_\local} \widetilde{\Theta}^{\eta,\Q}_{y,u,v}( p)$ via the CP method by \cite{jiang2020cuttingplane} and center at $\tilde{p}$ the proposal in our rejection sampler for the RGO, i.e., the proposal density is proportional to $\exp\brac{-\widetilde{\mathcal{P}}_1(p)}$ where 
\begin{equation}
\label{dolif:def:tildeP1:appen}
    \widetilde{\mathcal{P}}_1( p)
  = \frac{\| p-\tilde p\|^2 + \|\tilde p- q\|^2}{2\eta}
     - \frac{\delta}{\eta} \left(\| p-\tilde p\|+\|\tilde p- q\|\right) - \frac{\delta^2}{\eta} -\frac{b^2\eta}{2} + b v, 
\end{equation}
and $\delta = \sqrt{\frac{2\eta}{d+2}}$ as defined at the beginning of Section~\ref{sec:composite:mainpaper}. 

The following is the RGO implementation of Algorithm~\ref{dolif:alg:ASF_pi} with a separation oracle on $\Q$.
\begin{algorithm}[H]
	\caption{RGO implementation of Algorithm~\ref{dolif:alg:ASF_pi}}
	\label{dolif:RGO}
	\begin{algorithmic}[1]
		\State  Compute a $(d+2)^{-1}$-solution $\tilde{ p}$ of
        $\argmin_{ p\in \Q_\local} \widetilde{\Theta}^{\eta,\Q}_{y,u,v}( p)$ via the CP method by \cite{jiang2020cuttingplane}.
		\State Generate $U\sim {\cal U}[0,1]$ and $ p\sim \exp\brac{-\widetilde{\mathcal{P}}_1( p)}$ via Algorithm \ref{dolif:alg:samplingP1} in Appendix \ref{dolif:appen:specialalgo}. 
		\State  If
        \begin{equation}\label{dolif:event:separation}
			U \leq \exp\brac{-\widetilde{\Theta}_{y,u,v}^{\eta,\Q}( p)+\widetilde{\mathcal{P}}_1( p) }
		\end{equation}
		then accept $ p$; otherwise, reject $ p$ and go to step 2.
	\end{algorithmic}
\end{algorithm}
 Lemma~\ref{dolif:lem:compareP1P2} in Appendix~\ref{dolif:appen:supportlem} shows $\widetilde{\mathcal{P}}_1( p)\le \widetilde{\Theta}_{y,u,v}^{\eta,\Q}( p),\forall p$, so the acceptance test~\eqref{dolif:event:separation} is proper.


Finally, we refer to Theorem~\ref{dolif:theo:separation} in Section~\ref{sec:composite:mainpaper} for the proposal complexity and oracle complexity of Algorithm~\ref{dolif:RGO}. 

\textit{Remark:}
    Here we motivate the specific choice of radius $2r_{\local}$ in \eqref{dolif:def:localball:appen}. 
Per Theorem~\ref{theo:jiang}, bounding the oracle complexity of the CP method requires a lower bound on $\mbox{minwidth}(\Q_{\local})$. We obtain this bound by explicitly constructing an inscribed ball $\ball_{\mathrm{in}}$ of radius $\rin$ inside $\Q_{\local}$, which allows us to use $\mbox{minwidth}(\Q_{\local})\ge \mbox{minwidth}(\ball_{\mathrm{in}})=2\rin$. 
However, the construction of this inscribed ball (see Lemma~\ref{dolif:lem:localset}) is delicate and relies on a strict enlargement of the local set. Specifically, we require the local set to be $\Q\cap c r_{\local}\ball_{d+2}(q)$ with $c>1$ to ensure sufficient interior volume (the argument for the existence of $\ball_{\mathrm{in}}$ breaks down when $c=1$). Therefore, we work with the choice of $c=2$.


%%%%%%%%%%%%%%%%%%%%%%%%%%%%%%%%%%%%%%%%%%%%%%%%%%%%%%%%%%%%%%%%%%%%%%%%%%%%%%%%%%%%%%%%%%%%%%%%%%%%%%%%%%%%%%%%%%%%%%%%%%%%%%%%%%%%%%%%%%%%%%%%%%%%%%%%%%%%%%%%%%%%%%%%%%%%%%%%%%%%%%%%%%%%%%%%%%%%%%%%%%%%%%%%%%%%%%%%%%%%%%



\subsection{Proof of Theorem~\ref{dolif:theo:separation}}
\label{dolif:appen:proofsepatheorem}

 

Part~a is a consequence of Theorem~\ref{dolif:theo:cuttingplane} after taking $R=2r_{\local}$ and letting $a=b=d$, so what remains is to verify the claim in Part~b.  Denote $\mu_k$ the distribution of $(y_k,u_k,v_k)$ for the first step of Algorithm~\ref{dolif:alg:ASF_pi} and $n_{y,u,v}$ the expected number of proposals conditioned on the RGO input $(y_k,u_k,v_k)$. The fact that $d\mu_k/d\widetilde{\Pi}^{Y,U,V}\leq M$ from Lemma~\ref{dolif:lem:warmstart} implies 
\begin{align}
\label{dolif:warmstartimplicationsepa}
	\E_{\mu_k}[n_{y,u,v}] \le M \E_{\widetilde{\Pi}^{Y,U,V}}[n_{y,u,v}], 
\end{align}
and hence we will focus on bounding $\E_{\widetilde{\Pi}^{Y,U,V}}[n_{y,u,v}]$.



In Algorithm~\ref{dolif:RGO}, Steps~2 and~3 are a rejection sampler where the true potential function $\Theta$ and its proposal $\mathcal{P}$ are respectively 
 \begin{equation}\label{eq:corr4}
      \Theta = \widetilde{\Theta}^{\eta,\Q}_{y,u,v}, \quad \mathcal{P} = \widetilde{\mathcal{P}}_1
 \end{equation}
 in view of Lemma \ref{lem:rejection}. Thus, the latter result applies: it follows from $n_{y,u,v}=\E[F]$ in~\eqref{eq:generic-accept-rate} with the specification \eqref{eq:corr4} and inequality~\eqref{dolif:ineq:P2-P3} in Lemma~\ref{dolif:lem:compareP2P3} that
\begin{align}
\label{dolif:for:averagerejectionnotwarmsepa}
 \E_{\widetilde{\Pi}^{Y,U,V}} [n_{y,u,v}]&\stackrel{\eqref{eq:generic-accept-rate}}{=}\E_{\widetilde{\Pi}^{Y,U,V}}\left[\frac{N_{\widetilde{\mathcal{P}}_1}}{N_{\widetilde{\Theta}^{\eta,\Q}_{y,u,v}}} \right]\nn\\
 &=\int_{\R^{d+2}} \frac{\int_{\R^{d+2}} \exp(-\widetilde{\mathcal{P}}_1( x,s,t))  \D  x\D s\D t}{\int_{\R^{d+2}} \exp(-\widetilde{\Theta}^{\eta,\Q}_{y,u,v}( x,s,t))  \D  x\D s\D t} \widetilde{\Pi}^{Y,U,V}(y,u,v) \D y  \D u \D v\nn\\
 &\stackrel{\eqref{dolif:ineq:P2-P3}}{\leq} \int_{\R^{d+2}} \frac{\int_{\R^{d+2}} \exp(-\widetilde{\mathcal{P}}_2( x,s,t))  \D  x\D s\D t}{\int_{\R^{d+2}} \exp(-\widetilde{\Theta}^{\eta,\Q}_{y,u,v}( x,s,t))  \D  x\D s\D t} \widetilde{\Pi}^{Y,U,V}(y,u,v) \D y  \D u \D v .
\end{align}
Via the definition of $\widetilde{\mathcal{P}}_2$ at~\eqref{dolif:def:P2} with $b=d$ and Part~b of Lemma~\ref{lem:gaussianint}, 
\begin{align}
\label{dolif:inttildeP3}
&\int_{\R^{d+2}} \exp(-\widetilde{\mathcal{P}}_2( x,s,t))  \D  x\D s\D t
\leq \exp\brac{\frac{1}{4}+2}(2\pi\eta)^{(d+2)/2}\nn\\
&\qquad\cdot \exp\brac{-\frac{1}{2\eta}\left(\| (y,u,v-d\eta)-\operatorname{proj}_{\Q}(y,u,v-d\eta)\| - 2\delta\right)^2 }\cdot\exp\brac{\frac{d^2\eta}{2} - d v+\frac{16}{d+2} },
\end{align}
where $\delta = \sqrt{\frac{2\eta}{d+2}}$.
In addition, 
\begin{align}
\label{eq:intThetasimple}
    \int_{\R^{d+2}} \exp(-\widetilde{\Theta}^{\eta,\Q}_{y,u,v}( x,s,t))  \D  x\D s\D t\stackrel{\eqref{dolif:def:Theta}}{=}\int_{\R^{d+2}} \exp\brac{-dt-\frac{1}{2\eta}\norm{(x,s,t)-(y,u,v) }^2}   \D  x\D s\D t
\end{align}






Moreover, via part a of Lemma~\ref{lem:gaussianint},  we have %the marginal $ \widetilde{\Pi}^{Y,U,V}$ of $\widetilde{\Pi}$ is
\begin{align}
\label{dolif:for:tildepimarginal}
   \widetilde{\Pi}^{Y,U,V}(y,u,v)=\frac{1}{(2\pi \eta)^{\frac{d+2}{2}} \int_{\Q} \exp(- d t)\D x \D s \D t}\int_{\Q} \exp\brac{- d t-\frac{1}{2\eta}\norm{(x,s,t)-(y,u,v)}^2 }\D x \D s \D t.
   \end{align}
Plugging in $Z_{\Q}=\int_{\Q} e^{-dt}\D x\D s\D t$ and ~\eqref{dolif:inttildeP3},~\eqref{eq:intThetasimple},~\eqref{dolif:for:tildepimarginal} into~\eqref{dolif:for:averagerejectionnotwarmsepa} to get
\begin{align}
\label{dolif:separation:intermediate}
    & \E_{\widetilde{\Pi}^{Y,U,V}} [n_{y,u,v}]
     \leq \frac{1}{Z_{\Q}}\exp\brac{\frac{1}{4}+2+\frac{16}{d+2}}\nn\\&\cdot
     \int_{\R^{d+2}}\exp\brac{\frac{d^2\eta}{2} - d v }
     \exp\brac{-\frac{1}{2\eta}\left(\| (y,u,v-d\eta)-\operatorname{proj}_{\Q}(y,u,v-d\eta)\| - 2\delta\right)^2 } \D y \D u \D v .
\end{align}
Recall $ q=(y,u,v-d\eta)$ and $ q_{d+2}=(y,u,v-d\eta)_{d+2}=v-d\eta$. Then
\begin{align*}
  A:=&
  \int_{\R^{d+2}}\exp\brac{\frac{d^2\eta}{2} - d v }
  \exp\brac{-\frac{1}{2\eta}\left(\| (y,u,v-d\eta)-\operatorname{proj}_{\Q}(y,u,v-d\eta)\| - 2\delta\right)^2 } \D y \D u \D v\\
  =&\exp\brac{-\frac{d^2\eta}{2}}
  \int_{\R^{d+2}} \exp(-d q_{d+2})
  \exp\left(-\frac{\left(\dist( q,\Q)- 2\delta\right)^2}{2\eta}\right)  \D  q.
\end{align*}
Applying Lemma~\ref{dolif:lem:keytechnical} with $\tau=2\delta$ leads to 
\begin{align*}
A
&\leq Z_{\Q} \exp\brac{-\frac{d^2\eta}{2}+2\delta \widetilde{C}_L}
\Biggl[
\widetilde{C}_L\sqrt{2\pi\eta}\exp \left(\frac{\eta \widetilde{C}_L^2}{2}\right)
+1\Biggr].
\end{align*}
Substituting the bound on $A$ into \eqref{dolif:separation:intermediate}, we get
\begin{align*}
 \E_{\widetilde{\Pi}^{Y,U,V}} [n_{y,u,v}]
 &\leq
 \exp\brac{\frac{9}{4}+\frac{16}{d+2}} \cdot \exp\brac{-\frac{d^2\eta}{2}+2\delta \widetilde{C}_L}
\Biggl[
\widetilde{C}_L\sqrt{2\pi\eta}\exp \left(\frac{\eta \widetilde{C}_L^2}{2}\right)
+1\Biggr].
\end{align*}
Using $\eta=1/d^2$, noting that
\[
2\delta=2\sqrt{ \frac{2\eta}{d+2}}=\frac{2\sqrt{2}}{d\sqrt{d+2}} \le \frac{3}{d^{3/2}},
\]
and applying \eqref{dolif:warmstartimplicationsepa}, we obtain
\[
\E_{\mu_k}[n_{y,u,v}]
\le
M \exp\!\left(\frac74+\frac{16}{d}+ \frac{3\widetilde{C}_L}{d^{3/2}}\right)
\left[
\frac{\sqrt{2\pi} \widetilde C_L}{d} \exp\!\left(\frac{ \widetilde C_L^2}{2d^2}\right)
+1 \right],
\]
which is the declared bound \eqref{dolif:middlestepsepa}.

Finally, recall $\widetilde C_L=\sqrt{b^2+a^2+L_f^2}+\sqrt{a^2+L_h^2}$  (with $a=b=d$) from \eqref{dolif:def:notation}, then assuming $L_f={\cal O}(d)$, $L_h={\cal O}(d)$, and $M={\cal O}(1)$, one can verify that \eqref{dolif:middlestepsepa} becomes ${\cal O}(1)$. This completes the proof. \QEDA








%%%%%%%%%%%%%%%%%%%%%%%%%%%%%%%%%%%%%%%%%%%%%%%%%%%%%%%%%%%%%%%%%%%%%%%%%%%%%%%%%%%%%%%%%%%%%%%%%%%%%%%%%%%%%%%%%%%%%%%%%%%%%%%%%%%%%%%%%%%%%%%%%%%%%%%%%%%%%%%%%%%%%%%%%%

\subsection{Supporting lemmas}
\label{dolif:appen:supportlem}

To ensure that restricting the CP method to $\Q_{\local}$ is valid, we need the following result.

\begin{lemma}
\label{dolif:lem:minimizerinlocalball}  
Recall $\widetilde{\Theta}^{\eta,\Q}_{y,u,v}$ defined at~\eqref{dolif:def:Theta} and $\Q_{\local}=\Q\cap \ball_{\local}$ defined at~\eqref{dolif:def:localball}. The minimizer $ p_*=\argmin_{ p\in \R^{d+2}}\widetilde{\Theta}^{\eta,\Q}_{y,u,v}( p)=\proj_{\Q}( q)$
belongs to $\Q_{\local}$ a.s. In particular,
\[
 p_*=\underset{p\in \Q_\local}\argmin\widetilde{\Theta}^{\eta,\Q}_{y,u,v}( p).
\]
\end{lemma}


\begin{proof}
Recall in Section~\ref{sec:composite:mainpaper}, we denote $p$ as the previous RGO output (Algorithm~\ref{dolif:RGO}), so that $ p\in \Q$ a.s. This implies
$\norm{ p_*- q}\leq \norm{ p- q}=r_\local$. This implies
$ p_*\in r_\local \ball_{d+2}(q)\subseteq 2r_\local \ball_{d+2}( q)=\ball_\local$ a.s.
Since $ p_*\in \Q$ by definition ($\Q$ is closed), we deduce $ p_*\in \Q\cap \ball_\local=\Q_\local$ a.s.
Since $ p_*$ minimizes $\widetilde{\Theta}^{\eta,\Q}_{y,u,v}$ over $\Q$, it also minimizes $\widetilde{\Theta}^{\eta,\Q}_{y,u,v}$ over the smaller domain $\Q_\local=\Q\cap \ball_\local$.
\end{proof}



 



%%%%%%%%%%%%%%%%%%%%%%%%%%%%%%%%%%%%%%%%%%%%%%%%%%%%%%%%%%%%%%%%%%%%%%%%%%%%%%%%%%%%%%%%%%%%%%%%%%%%%%%%%%%%%%%%%%%%%%%%%%%%%%%%%%%%%%%%%%%%%%%%%%%%%%%%%%%%%%%%%%%%%%%%%%%%%%%%%%%%%%%%%%%%%%%%%%%%%%%%

Here we provide a bound on the PI constant of $\tilde{\pi}$ at \eqref{dolif:def:piandQ}. 
This result is analogous to Lemma~\ref{silif:lem:PIconstant}.

\begin{lemma}
\label{dolif:lem:PIconstant}
    Assume the distribution $\tilde{\pi}$ at \eqref{dolif:def:piandQ} has $a=b=d$. Then, it satisfies a PI and $  \widetilde{C}_{\mathrm{PI}}(\tilde{\pi})\leq C\log (d+2)\brac{\opnorm{\cov{\tilde{\nu}}}+1 }$ for a universal $C$, where $\tilde{\nu}$ is defined at \eqref{dolif:def:nu}. 
\end{lemma}


\begin{proof}   
Note that $\tilde{\pi}\propto \exp\brac{-dt-I_{\Q}}$ is a log-concave measure on $\R^{d+2}$. It is a well-known fact (\cite{kannan1995isoperimetric,bobkov1999isoperimetric}) that any log-concave measure satisfies a PI with finite PI constant. Then per \cite[Remark 7.12]{klartag2025lecturenoteisoperimetric} (see also \cite{klartag2023logarithmic}), the PI constant of $\tilde{\pi}$ is bounded as $ \widetilde{C}_{\mathrm{PI}}(\tilde{\pi})\leq C'\log (d+2) \opnorm{\cov{\tilde{\pi}}}$
for a universal constant $C'$. 

With respect to the distributions $\tilde{\nu},\gamma$ and $\tilde{\pi}$ at \eqref{dolif:def:nu}, \eqref{dolif:def:gammaandK} and \eqref{dolif:def:piandQ}, we have the relations $\tilde{\pi}^{X,S}=\gamma,\gamma^X=\tilde{\nu}$. Then $\tilde{\pi}^{X,S}=\gamma$ and \cite[Lemma 2.5]{kook2025algodiffusion} imply 
\begin{align*}
    \opnorm{\cov{\tilde{\pi}}}\leq 2\brac{\opnorm{\cov{\tilde{\pi}^{X,S}}}+160 }= 2\brac{\opnorm{\cov{\gamma}}+160 }.
\end{align*}
Similarly, $\gamma^X=\tilde{\nu}$ and \cite[Lemma 2.5]{kook2025algodiffusion} give 
\begin{align*}
    \opnorm{\cov{\gamma}}\leq 2\brac{\opnorm{\cov{\gamma^{X}}}+160 }= 2\brac{\opnorm{\cov{\tilde{\nu}}}+160 }.
\end{align*}
Combining the above calculations yield the desired conclusion. 
\end{proof}


The following warm-start result is analogous to Lemma~\ref{silif:lem:warmstart} in the constrained sampling case.
\begin{lemma}
\label{dolif:lem:warmstart}
    An $M$-warm start $\tilde{\nu}_0$ for $\tilde{\nu}=\tilde{\pi}^X$ (condition~(B.3)) naturally induces an $M$-warm start $\tilde{\pi}_0$ for $\tilde{\pi}$, i.e., $\frac{d\tilde{\pi}_0}{d\tilde{\pi}}\leq M$. Moreover, denote $\tilde{\pi}_\eta$ and $\tilde{\pi}^k$ respectively the output of the first step of Algorithm~\ref{dolif:alg:ASF_pi} and the output of the $k$-th iterate of Algorithm~\ref{dolif:alg:ASF_pi}, then it holds that $    \frac{d\tilde{\pi}_\eta}{d\widetilde{\Pi}^{Y,U,V}}\leq M$ and $\frac{d\tilde{\pi}^k}{d\widetilde{\Pi}^{X,S,T}}\leq M$. In other words, the warmness condition holds for every step of Algorithm~\ref{dolif:alg:ASF_pi}. 
\end{lemma}


\begin{proof}  
For the first part of the lemma, we follow the same construction as in the proof of Lemma~\ref{silif:lem:warmstart}.
Generate $x\sim \tilde{\nu}_0$ (by, for instance, the Gaussian cooling procedure in~\cite[Section 3]{kook2025algodiffusion}) and then sample $(s,t)\sim \tilde{\pi}^{S,T |  X=x}$.
Concretely, under $\tilde{\pi}$ the conditional $S |  X=x$ is one-dimensional with density proportional to $e^{-as}\mathbf{1}\{s\ge h(x)/a\}$, and can be generated by sampling $u_1\sim U[0,1]$ and setting $s=F_S^{-1}(u_1)$, where
\[
F_S(s)=\Pr(S\le s |  X=x)=1-e^{h(x)-as},\qquad s\ge \frac{h(x)}{a}.
\]
Likewise, $T |  (X,S)=(x,s)$ has density proportional to $e^{-bt}\mathbf{1}\{t\ge (f(x)+as)/b\}$, and can be generated by sampling $u_2\sim U[0,1]$ and setting $t=F_T^{-1}(u_2)$, where
\[
F_T(t)=\Pr(T\le t |  X=x,S=s)=1-e^{f(x)+as-bt},\qquad t\ge \frac{f(x)+as}{b}.
\]
Denote by $\tilde{\pi}_0$ the law of $(x,s,t)$ produced in this way. Then $\tilde{\pi}_0$ is $M$-warm with respect to $\tilde{\pi}$ since the conditional factors cancel:

\[
\frac{\D\tilde{\pi}_0}{\D\tilde{\pi}}(x,s,t)=
\frac{\left(\frac{\D\tilde{\nu}_0}{\D x}\right)(x)\,\tilde{\pi}^{S,T \mid X=x}(s,t)}
{\left(\frac{\D \tilde{\nu}}{\D x}\right)(x)\,\tilde{\pi}^{S,T \mid X=x}(s,t)}
=
\frac{\D \tilde{\nu}_0}{\D \tilde{\nu}}(x)
\le M .\]
For the second part of the lemma, let $U\subseteq \R^{d+2}$ be measurable. For $(y,u,v)\in \R^{d+2}$, set
$U-(y,u,v):=\{(x,s,t)\in \R^{d+2}:(x,s,t)+(y,u,v)\in U \}$, and denote by $\frak{g}(\cdot)$ the density of
$\mathcal{N}(0,\eta I_{d+2})$.
By $\frac{d\tilde{\pi}_0}{d\tilde{\pi}}\le M$, $\tilde{\pi}_\eta=\tilde{\pi}_0*\frak{g}$, and
$\widetilde{\Pi}^{Y,U,V}=\widetilde{\Pi}^{X,S,T}*\frak{g}$ with $\widetilde{\Pi}^{X,S,T}=\tilde{\pi}$, we have
\begin{align*}
\tilde{\pi}_\eta(U)
&=\int_{\R^{d+2}} \tilde{\pi}_0\left(U-(y,u,v)\right) \frak{g}(y,u,v)  \D y \D u \D v\\
&=\int_{\R^{d+2}}\brac{\int_{U-(y,u,v)} \frac{\D\tilde{\pi}_0}{\D\widetilde{\Pi}^{X,S,T}}(x,s,t) \widetilde{\Pi}^{X,S,T}(\D x,\D s,\D t) }
\frak{g}(y,u,v)  \D y \D u \D v\\
&\le M\int_{\R^{d+2}} \widetilde{\Pi}^{X,S,T}\left(U-(y,u,v)\right) \frak{g}(y,u,v)  \D y \D u \D v
= M \widetilde{\Pi}^{Y,U,V}(U),
\end{align*}
and hence $\frac{d\tilde{\pi}_\eta}{d\widetilde{\Pi}^{Y,U,V}}\le M$.
A similar argument (using the fact that Step~2 of Algorithm~\ref{dolif:alg:ASF_pi} is the exact conditional update with stationary
law $\widetilde{\Pi}^{X,S,T |  Y,U,V}$) gives $\frac{d\tilde{\pi}^k}{d\widetilde{\Pi}^{X,S,T}}\le M$ for every $k$.
Thus, warmness holds for every step of Algorithm~\ref{dolif:alg:ASF_pi}.
\end{proof}





 %%%%%%%%%%%%%%%% %%%%%%%%%%%%%%%% %%%%%%%%%%%%%%%% %%%%%%%%%%%%%%%% %%%%%%%%%%%%%%%% %%%%%%%%%%%%%%%% %%%%%%%%%%%%%%%% %%%%%%%%%%%%%%%% %%%%%%%%%%%%%%%%

The next result ensures that the acceptance test~\eqref{dolif:event:separation} is well-defined.
Although the upcoming Lemma~\ref{dolif:lem:compareP1P2} is the composite analogue of
Lemma~\ref{silif:lem:compareThetaP1} from the constrained sampling case, there is one important difference:
Lemma~\ref{silif:lem:compareThetaP1} relies on the Lipschitz assumption, whereas
Lemma~\ref{dolif:lem:compareP1P2} does not.

\begin{lemma}\label{dolif:lem:compareP1P2}
Recall the notations $q=(y,u,v-b\eta)$, $\delta = \sqrt{\frac{2\eta}{d+2}}$ and the function $\widetilde{\mathcal{P}}_1$ at~\eqref{dolif:def:tildeP1}. Then under condition~(B1), we have $\widetilde{\mathcal{P}}_1( p)\le \widetilde{\Theta}_{y,u,v}^{\eta,\Q}( p)$ for all $ p\in\mathbb{R}^{d+2}$, and thus the acceptance test at~\eqref{dolif:event:separation} is well-defined.
\end{lemma}

\begin{proof}  
We assume $a,b> 0$ in the introduction, which combines with  condition~(B1) implies the set $\Q$ at~\eqref{dolif:def:piandQ} is convex. Now set
\begin{align*}
    \widetilde{\mathcal{P}}_0( p) = \frac{1}{2\eta}\norm{ p - \proj_{\Q}( q)}^2 + \frac{1}{2\eta}\norm{\proj_{\Q}( q) -  q}^2 + bv - \frac{b^2\eta}{2}. 
\end{align*}
Then
    \begin{align}
    \label{dolif:eq:Thetarewrite}
        \widetilde{\Theta}^{\eta,\Q}_{y,u,v}( p) &= I_{\Q}( p) + bt + \frac{1}{2\eta}\norm{ p - (y,u,v)}^2 = I_{\Q}( p) + \frac{1}{2\eta}\norm{ p -  q}^2 + bv - \frac{b^2\eta}{2} \nn\\
        &\ge \frac{1}{2\eta}\norm{\proj_{\Q}( q) -  q}^2 + \frac{1}{2\eta}\norm{ p - \proj_{\Q}( q)}^2 + bv - \frac{b^2\eta}{2} = \widetilde{\mathcal{P}}_0( p).
    \end{align}
Specifically, the inequality in the second-to-last line holds trivially if $ p \notin \Q$, and by the Pythagorean property of projections onto the convex set $\Q$ if $ p \in \Q$.


The rest of the proof is devoted to showing  $\widetilde{\mathcal{P}}_1( p)\le \widetilde{\mathcal{P}}_0( p),\forall p$.  Let $  p_* = \argmin_{ p\in \Q_\local} \widetilde{\Theta}^{\eta,\Q}_{y,u,v}( p)$. Since $\Q_\local=\Q\cap \ball_\local$ is closed, $\tilde{ p}$ as the $(d+2)^{-1}$-solution to the problem
$\argmin_{ p\in \Q_\local}\widetilde{\Theta}^{\eta,\Q}_{y,u,v}( p)$ belongs to $\Q_\local$. Moreover, the map $ p\mapsto \frac{1}{2\eta}\| p- q\|^2$ is $\eta^{-1}$-strongly convex, hence
\begin{align}
\label{dolif:difference:tildexi}
\|\tilde{ p}- p_*\|
\le \sqrt{2\eta\left(\widetilde{\Theta}^{\eta,\Q}_{y,u,v}(\tilde{ p})-\widetilde{\Theta}^{\eta,\Q}_{y,u,v}( p_*)\right)}
\le \sqrt{\frac{2\eta}{d+2}}
=: \delta.
\end{align}

At this point, \eqref{dolif:difference:tildexi} and the triangle inequality imply that, for every $ p\in\mathbb{R}^{d+2}$,
\begin{align}
\label{dolif:estimatewihoutsquare-xi}
\| p- p_*\|
\ge \| p-\tilde{ p}\|-\|\tilde{ p}- p_*\|
\ge \| p-\tilde{ p}\|-\delta.
\end{align}
It follows that
\begin{align*}
\| p- p_*\|^2
\ge \left(\| p-\tilde{ p}\|-\delta\right)^2
= \| p-\tilde{ p}\|^2-2\delta\| p-\tilde{ p}\|+\delta^2
\ge \| p-\tilde{ p}\|^2-2\delta\| p-\tilde{ p}\|-\delta^2,
\end{align*}
and therefore, after rearrangement,
\begin{align}
\label{dolif:firstestimate-composite}
\| p-\tilde{ p}\|^2-2\delta\| p-\tilde{ p}\|-\delta^2
\le \| p- p_*\|^2.
\end{align}

Similarly, \eqref{dolif:difference:tildexi} and the triangle inequality give
\begin{align}
\label{dolif:estimatewihoutsquare-rho}
\| p_*- q\|
\ge \|\tilde{ p}- q\|-\|\tilde{ p}- p_*\|
\ge \|\tilde{ p}- q\|-\delta,
\end{align}
so
\begin{align*}
\| p_*- q\|^2
\ge \left(\|\tilde{ p}- q\|-\delta\right)^2
= \|\tilde{ p}- q\|^2-2\delta\|\tilde{ p}- q\|+\delta^2
\ge \|\tilde{ p}- q\|^2-2\delta\|\tilde{ p}- q\|-\delta^2,
\end{align*}
which can be rearranged as
\begin{align}
\label{dolif:secondestimate-composite}
\|\tilde{ p}- q\|^2-2\delta\|\tilde{ p}- q\|-\delta^2
\le \| p_*- q\|^2.
\end{align}
Combining \eqref{dolif:firstestimate-composite} and \eqref{dolif:secondestimate-composite} yields
\begin{align*}
&\| p-\tilde{ p}\|^2+\|\tilde{ p}- q\|^2
-2\delta\left(\| p-\tilde{ p}\|+\|\tilde{ p}- q\|\right)
-2\delta^2\le \| p- p_*\|^2+\| p_*- q\|^2.
\end{align*}
Multiplying by $1/(2\eta)$ and adding the common term $-\frac{b^2\eta}{2}+bv$ on both sides, we obtain
$\widetilde{\mathcal{P}}_1( p)\le \widetilde{\mathcal{P}}_0( p)$ for all $ p\in\mathbb{R}^{d+2}$. Finally, we can combine the previous result with $\widetilde{\mathcal{P}}_0( p)\le \widetilde{\Theta}_{y,u,v}^{\eta,\Q}( p)$ for all $ p\in\mathbb{R}^{d+2}$ at~\eqref{dolif:eq:Thetarewrite} to complete the proof. 
\end{proof}



 %%%%%%%%%%%%%%%% %%%%%%%%%%%%%%%% %%%%%%%%%%%%%%%% %%%%%%%%%%%%%%%% %%%%%%%%%%%%%%%% %%%%%%%%%%%%%%%% %%%%%%%%%%%%%%%% %%%%%%%%%%%%%%%% %%%%%%%%%%%%%%%%


We now introduce a function which will be used
only in the rejection analysis. Recall $q=(y,u,v-b\eta)$ defined at~\eqref{dolif:def:notation} and $\tilde p$ defined in Algorithm~\ref{dolif:RGO}. For $p=(x,s,t)\in \R^{d+2}$, define the real-valued function 
\begin{equation}\label{dolif:def:P2}
  \widetilde{\mathcal{P}}_2( p)
  = -\frac{b^2\eta}{2} + b v
     + \frac{1}{2\eta}\left(
        \left(\| p-\tilde p\| - \delta\right)^2
        + \left(\| q- p_*\| - 2\delta\right)^2
        - 16\delta^2
     \right). 
\end{equation}
The following result is analogous to Lemma~\ref{silif:lem:compareP1P2}.


\begin{lemma}
\label{dolif:lem:compareP2P3}
Assume the setting of Lemma~\ref{dolif:lem:compareP1P2}, and let
\(\widetilde{\mathcal{P}}_1\) and \(\widetilde{\mathcal{P}}_2\) be given by respectively
\eqref{dolif:def:tildeP1} and \eqref{dolif:def:P2}.
Then for every \( p=(x,s,t)\in\mathbb{R}^{d+2}\),
\begin{equation}\label{dolif:ineq:P2-P3}
  \widetilde{\mathcal{P}}_2( p)\;\le\;\widetilde{\mathcal{P}}_1( p).
\end{equation}
\end{lemma}

\begin{proof}  
For brevity, set $ r = \| p-\tilde p\|, n = \| q-\tilde p\|, D = \| q- p_*\|.$ Recall from \eqref{dolif:difference:tildexi} that
\begin{equation}\label{dolif:eq:delta-approx-norm-again}
  \|\tilde p- p_*\|\le \delta.
\end{equation}
By the triangle inequality, 
\begin{equation}\label{dolif:eq:estimatewihoutsquareP3-composite}
  n=\| q-\tilde p\|
  \le \| q- p_*\|+\| p_*-\tilde p\|
  \stackrel{\eqref{dolif:eq:delta-approx-norm-again}}{\le} D+\delta.
\end{equation}
Also, the reverse triangle inequality gives \[|n-D|
  =\left|\| q-\tilde p\|-\| q- p_*\|\right|
  \le \|\tilde p- p_*\|
  \stackrel{\eqref{dolif:eq:delta-approx-norm-again}}{\le} \delta.\] 
In particular, this leads to \(n\ge D-\delta\), hence $ n^2 \ge (D-\delta)^2 = D^2 - 2D\delta + \delta^2 \ge D^2 - 2D\delta - 7\delta^2$. Thus, 
\begin{equation}\label{dolif:eq:firstestimateP3-composite}
  D^2 - 2D\delta - 7\delta^2 \;\le\; n^2.
\end{equation}
Via \eqref{dolif:eq:estimatewihoutsquareP3-composite}, \eqref{dolif:eq:firstestimateP3-composite},
and the formula of \(\widetilde{\mathcal{P}}_1\) at \eqref{dolif:def:tildeP1}, we get
\begin{align*}
  \widetilde{\mathcal{P}}_1( p)
  &= -\frac{b^2\eta}{2} + b v
     + \frac{1}{2\eta}\left(
        r^2 + n^2 - 2\delta(r+n) - 2\delta^2
     \right) \\
  &\stackrel{\eqref{dolif:eq:firstestimateP3-composite}}{\ge}
     -\frac{b^2\eta}{2} + b v
     + \frac{1}{2\eta}\left(
        r^2 + (D^2-2D\delta-7\delta^2) - 2\delta r - 2\delta n - 2\delta^2
     \right) \\
  &\stackrel{\eqref{dolif:eq:estimatewihoutsquareP3-composite}}{\ge}
     -\frac{b^2\eta}{2} + b v
     + \frac{1}{2\eta}\left(
        r^2 - 2\delta r + D^2 - 4D\delta - 11\delta^2
     \right).
\end{align*}
Finally, note that $r^2 - 2\delta r = (r-\delta)^2 - \delta^2$ and  $D^2 - 4D\delta = (D-2\delta)^2 - 4\delta^2$, so that
\[
  r^2 - 2\delta r + D^2 - 4D\delta - 11\delta^2
  = (r-\delta)^2 + (D-2\delta)^2 - 16\delta^2.
\]
Substituting this identity into the previous display yields
\[
  \widetilde{\mathcal{P}}_1( p)
  \;\ge\;
  -\frac{b^2\eta}{2} + b v
  + \frac{1}{2\eta}\left(
       (r-\delta)^2 + (D-2\delta)^2 - 16\delta^2
    \right)
  = \widetilde{\mathcal{P}}_2( p),
\]
which completes the proof.
\end{proof}


%%%%%%%%%%%%%%%%%%%%%%%%%%%%%%%%%%%%%%%%%%%%%%%%%%%%%%%%%%%%
%%%%%%%%%%%%%%%%%%%%%%%%%%%%%%%%%%%%%%%%%%%%%%%%%%%%%%%%%%%%


The next result is analogous to Lemma~\ref{silif:lem:envelope}. The proof follows the same overall strategy, but the construction of the envelope set is different.
\begin{lemma}
\label{dolif:lem:envelope}
Assume condition~(B2) holds. For all $\tau>0$, we have 
\[\int_{0}^{\infty} \exp \left(-\frac{(r-\tau)^2}{2\eta}\right)
B_{\Q_r}  \D r
\ \le\
\widetilde{C}_L  Z_{\Q}  \sqrt{2\pi\eta}\;
\exp \left(\tau \widetilde{C}_L + \frac{\eta \widetilde{C}_L^2}{2}\right)+Z_{\Q}\brac{\exp(\tau \widetilde{C}_L)-1}. \]
\end{lemma}

\begin{proof}  
The proof consists of two steps. In Step 1, we construct an envelope set which contains $\Q_r$ and is much easier to integrate on compared to $\Q_r$; while in Step 2, we perform integration by parts (IBP) to obtain the desired bound. The reason we need IBP to bound the quantity $\int_{0}^{\infty} \exp \left(-\frac{(r-\tau)^2}{2\eta}\right)
B_{\Q_r}  \D r$ is that $B_{\Q_r}$ as an integral on the boundary $\partial \Q_r$ is quite hard to bound; so via IBP, $B_{\Q_r}$ is replaced by $Z_{\Q_r}$ which is easier to handle.

Note that relative to the proof of Lemma~\ref{silif:lem:envelope}, Step~1 below differs from the corresponding Step~1 there, whereas Step~2 is essentially the same.

\emph{Step 1.}
Write $\Q=\widetilde{C}_1\cap \widetilde{C}_2$ where \begin{align*}
    \widetilde{C}_1:=\left\{(x,s,t)\in\R^{d+2}:\ s\ge \frac{h(x)}a\right\},
\qquad
\widetilde{C}_2:=\left\{(x,s,t)\in\R^{d+2}:\ t\ge \frac{f(x)+as}b\right\}. 
\end{align*}
Fix $r>0$. Since $(A\cap B)+H\subseteq (A+H)\cap(B+H)$ for any sets $A,B,H$, let $A=\widetilde{C}_1,B=\widetilde{C}_2$ and $H=r\ball_{d+2}(0)$ to get
\begin{align}
\label{dolif:setQr}
\Q_r=(\widetilde{C}_1\cap \widetilde{C}_2)_r \subseteq (\widetilde{C}_1)_r\cap (\widetilde{C}_2)_r.
\end{align}
For $\widetilde{C}_1$, apply Lemma~\ref{lem:parallelsetcontainment} with $\phi=\frac{h}{a}$.
Since $\phi$ is $\frac{L_h}{a}$-Lipschitz, we obtain
\[
(\widetilde{C}_1)_r \subseteq \left\{(x,s,t):\ s\ge \frac{h(x)}{a}-r\sqrt{1+\frac{L_h^2}{a^2}}\right\}
=\left\{(x,s,t):\ s\ge \frac{h(x)}{a}-\frac r a \widetilde{C}_h\right\}.
\]
For $\widetilde{C}_2$, apply Lemma~\ref{lem:parallelsetcontainment} with $\phi(x,s):=\frac{f(x)+as}{b}$ for $(x,s)\in \R^{d+1}$. Since $\phi$ is $\frac{\sqrt{L_f^2+a^2}}{b}$-Lipschitz in $(x,s)$, we obtain \[(\widetilde{C}_2)_r \subseteq \left\{(x,s,t):\ t\ge \frac{f(x)+as}{b}-r\sqrt{1+\tfrac{L_f^2+a^2}{b^2}}\right\}
=\left\{(x,s,t):\ t\ge \frac{f(x)+as}{b}-\frac{r\widetilde{C}_f}{b} \right\}.\] 
Therefore, via \eqref{dolif:setQr}, we conclude $\Q_r\subseteq \widetilde{\mathcal{E}}_r$ where \[\widetilde{\mathcal{E}}_r
:=\left\{(x,s,t):\ s\ge \frac{h(x)}{a}-\frac r a \widetilde{C}_h,\quad
t\ge \frac{f(x)+as}{b}-\frac r b \widetilde{C}_f\right\}.\] This further implies
\begin{align}
\label{dolif:firstbound_Zer}
Z_{\Q_r}\le Z_{\widetilde{\mathcal{E}}_r},
\end{align}
for which 
\begin{align*}
Z_{\widetilde{\mathcal{E}}_r}
&=\int_{\R^d}\int_{s\ge \frac{h(x)}{a}-\frac r a \widetilde{C}_h}
\int_{t\ge \frac{f(x)+as}{b}-\frac r b \widetilde{C}_f} e^{-bt}  \D t  \D s  \D x \nonumber\\
&=\frac{e^{r\widetilde{C}_f}}{b}\int_{\R^d} e^{-f(x)}
\left(\int_{s\ge \frac{h(x)}{a}-\frac r a \widetilde{C}_h} e^{-as}  \D s\right)  \D x \nonumber\\
&\stackrel{\eqref{eq:intoverQ}}{=}\frac{e^{r\widetilde{C}_f}}{ab}\int_{\R^d} e^{-f(x)}e^{-h(x)} e^{r\widetilde{C}_h}  \D x
= e^{r(\widetilde{C}_f+\widetilde{C}_h)}  Z_{\Q}
= e^{r\widetilde{C}_L}  Z_{\Q}.
\end{align*}
In particular, we apply Lemma \ref{lem:basicfacts}(b) above with $F(x)=\frac{e^{r\widetilde C_f}}{b} e^{-f(x)},
\psi(x)=\frac{h(x)}{a}$ and $\gamma=\frac{r}{a}\widetilde C_h$. Therefore, we get $Z_{\Q_r}\le Z_{\widetilde{\mathcal{E}}_r}=e^{r\widetilde{C}_L}Z_{\Q}$. 


\emph{Step 2.}
Set $w_{\eta,\tau}(r):=\exp \left(-\frac{(r-\tau)^2}{2\eta}\right)$.
By \eqref{eq:IrderivativeofQr} in Lemma~\ref{lem:coarea},
$\frac{d}{\D r}Z_{\Q_r}=B_{\Q_r}$ a.e.\ $r$.
Moreover, since $Z_{\Q_r}\le e^{r\widetilde{C}_L}Z_{\Q},\forall r\ge0$ from the conclusion of Step~1 and $w_{\eta,\tau}(r)$ decays exponentially,
we have $Z_{\Q_r}w_{\eta,\tau}(r)\to0$ as $r\to\infty$.
Also $Z_{\Q_0}=Z_{\Q}$ and $w_{\eta,\tau}(0)=e^{-\tau^2/(2\eta)}$.
Integration by parts and the fact that  $w_{\eta,\tau}'(r)=-(r-\tau)/\eta\; w_{\eta,\tau}(r)$ imply
\begin{align}
\label{dolif:eq:IBP-master}
&\int_{0}^{\infty} B_{\Q_r}   w_{\eta,\tau}(r)  \D r
=\int_{0}^{\infty} Z_{\Q_r}'   w_{\eta,\tau}(r)  \D r
= Z_{\Q_r}w_{\eta,\tau}(r)\bigg|_{0}^{\infty}
-\int_{0}^{\infty} Z_{\Q_r}   w_{\eta,\tau}'(r)  \D r\nonumber\\
\quad&= -Z_{\Q} e^{-\tau^2/(2\eta)}
+ \int_{0}^{\infty} Z_{\Q_r}  \frac{r-\tau}{\eta}  w_{\eta,\tau}(r)  \D r.
\end{align}
We will bound the second term on the right hand side of~\eqref{dolif:eq:IBP-master} using $Z_{\Q}\leq Z_{\Q_r}\leq e^{r\widetilde{C}_L}Z_{\Q}$ for $r\geq 0$.  
\begin{align*}
    \int_{0}^{\infty} Z_{\Q_r}  \frac{r-\tau}{\eta}  w_{\eta,\tau}(r)  \D r &= \int_0^\tau Z_{\Q_r}  \frac{r-\tau}{\eta}  w_{\eta,\tau}(r)  \D r + \int_\tau^\infty Z_{\Q_r}  \frac{r-\tau}{\eta}  w_{\eta,\tau}(r)  \D r\\
    &\leq Z_{\Q}\int_0^\tau \frac{r-\tau}{\eta}    w_{\eta,\tau}(r)  \D r+ Z_{\Q}\int_\tau^\infty \frac{r-\tau}{\eta}  e^{r\widetilde{C}_L}  w_{\eta,\tau}(r)  \D r \\
    &= - Z_{\Q} w_{\eta,\tau}(r) \bigg|_{0}^{\tau} + Z_{\Q}\int_\tau^\infty \frac{r-\tau}{\eta}  e^{r\widetilde{C}_L}  w_{\eta,\tau}(r)  \D r \\
    &= - Z_{\Q} (1 - e^{-\tau^2/(2\eta)}) + Z_{\Q}\int_\tau^\infty \frac{r-\tau}{\eta}  e^{r\widetilde{C}_L}  w_{\eta,\tau}(r)  \D r.
\end{align*}
Now via formula~\eqref{eq:kindofibp} in Lemma~\ref{lem:basicfacts}, 
\begin{align*}
     \int_{0}^{\infty} Z_{\Q_r}  \frac{r-\tau}{\eta}  w_{\eta,\tau}(r)  \D r &\leq - Z_{\Q} (1 - e^{-\tau^2/(2\eta)}) +Z_{\Q}\brac{ \widetilde{C}_L\int_\tau^{\infty} e^{\widetilde{C}_Lr}w_{\eta,\tau}(r) \D r + e^{\widetilde{C}_L \tau} }. 
\end{align*}
Then \eqref{dolif:eq:IBP-master} becomes
\begin{align}
\label{finalboundintBQr_comp}
\int_{0}^{\infty} B_{\Q_r}   w_{\eta,\tau}(r)  \D r
\le Z_{\Q}\left(e^{\widetilde{C}_L\tau}-1\right)+\widetilde{C}_L Z_{\Q}\int_{\tau}^{\infty} e^{r\widetilde{C}_L}w_{\eta,\tau}(r)  \D r.
\end{align}
Finally, we estimate the Gaussian integral. Via $\widetilde{C}_L r-\frac{(r-\tau)^2}{2\eta}
=-\frac{1}{2\eta}\left(r-(\tau+\eta \widetilde{C}_L)\right)^2+\tau \widetilde{C}_L+\frac{\eta \widetilde{C}_L^2}{2}$, we can further write 
\[\int_{0}^{\infty} e^{r\widetilde{C}_L}w_{\eta,\tau}(r)  \D r
\le e^{\tau \widetilde{C}_L+\frac{\eta \widetilde{C}_L^2}{2}}
\int_{-\infty}^{\infty}\exp \left(-\frac{(r-(\tau+\eta \widetilde{C}_L))^2}{2\eta}\right)  \D r=\sqrt{2\pi\eta}\;\exp \left(\tau \widetilde{C}_L+\frac{\eta \widetilde{C}_L^2}{2}\right).\] Combining with~\eqref{finalboundintBQr_comp} to reach the stated bound. 
\end{proof}


The next result is a consequence of Lemma~\ref{dolif:lem:envelope}. 
It is also analogous to Lemma \ref{silif:lem:keytechnical}.

\begin{lemma}
\label{dolif:lem:keytechnical}
Assume condition~(B2) holds. Let $Z_{\Q}:=Z_{\Q_0}$. Then for $\tau\geq 0$ and $ q_{d+2}=v-b\eta$, we have \begin{align*}&\int_{\R^{d+2}} e^{-bq_{d+2}} \exp\brac{-\frac{\brac{\dist ( q,\Q)-\tau}^2}{2\eta}}d q &\le Z_{\Q} \exp(\tau \widetilde{C}_L) \left[ 
     \widetilde{C}_L \sqrt{2\pi\eta} \exp\brac{ \frac{\eta \widetilde{C}_L^2}{2}} +1\right].
    \end{align*}
\end{lemma}


\begin{proof}   We split the proof into two cases. 
~\\
\underline{Case I:} When $q\in \Q$, we have $\dist(q,\Q)=0$. Then
\begin{align*}
    \int_{\Q} \exp\brac{-bq_{d+2}} \exp\brac{-\frac{\brac{\dist (q,\Q)-\tau}^2}{2\eta}}\D q
    &=\int_{\Q} \exp\brac{-bq_{d+2}} \exp\brac{-\frac{\tau^2}{2\eta}}   \D q \stackrel{\eqref{def:ZandB}}{=} e^{-\frac{\tau^2}{2\eta}} Z_{\Q}.
\end{align*}

~\\
\underline{Case II:} When $q \notin \Q$, we denote $r(q) = \dist(q,\Q)$ and $w_{\eta,\tau}(r) = \exp\brac{-\frac{(r-\tau)^2}{2\eta}}$. Using the co-area formula, we write
\begin{align*}
    &\int_{\Q^C} \exp\brac{-bq_{d+2}} \exp\brac{-\frac{\brac{\dist (q,\Q)-\tau}^2}{2\eta}}\D q\\
    &=\int_0^\infty \exp\brac{-\frac{(r-\tau)^2}{2\eta}}\left(\int_{\partial \Q_r}\exp\brac{-bq_{d+2}}\D S^{d+1}(q)\right)\D r = \int_0^\infty B_{\Q_r} w_{\eta,\tau}(r) \D r.
\end{align*}


We apply the crucial bound in Lemma~\ref{dolif:lem:envelope} to get
\begin{align*}
    \int_0^\infty B_{\Q_r} w_{\eta,\tau}(r) \D r\le  Z_{\Q} \brac{\widetilde{C}_L\sqrt{2\pi\eta}\;\exp\brac{\tau \widetilde{C}_L + \frac{\eta \widetilde{C}_L^2}{2}}+\exp(\tau \widetilde{C}_L)-1}.
\end{align*}

Combining Case I, Case II and the fact that $e^{-\frac{\tau^2}{2\eta}} \leq 1$  yields the desired bound. 
\end{proof}


%%%%%%%%%%%%%%%%%%%%%%%%%%%%%%%%%%%%%%%%%%%%%%%%%%%%%%%%%%%%%%%%%%%%%%%%%%%%%%%%%%%%%%%%%%%%%%%%%%%%%%%%%%%%%%%%%%%%%%%%%%%%%%%%%%%%%%%%%%%%%%%%%%%%%%%%%%%%%%%%%%%%%%%%%%%%%%%%%%%%%%%%%%%%%%%%%%%%%%%%%%%%%%%%%%%%%%%%%%%%%%%%%%%%%%%%%%%%%%%%%%%








\subsection{Results about the cutting-plane method by \cite{jiang2020cuttingplane}}
\label{dolif:appen:cuttingplane}




Per Lemma~\ref{dolif:lem:minimizerinlocalball}, in order to find an approximate solution to $ p_*=\argmin_{ p\in \R^{d+2}}\widetilde{\Theta}^{\eta,\Q}_{y,u,v}( p)$, we will apply the CP method (Theorem~\ref{theo:jiang}) to optimize $\widetilde{\Theta}^{\eta,\Q}_{y,u,v}( p)$ over $\Q_{\local}=\Q\cap 2r_{\local}\ball_{d+2}( q)$. In view of Theorem~\ref{theo:jiang}, we need to ensure $\gamma$ is not too big, which is the content of the next result. 

This next result is also one of the subtle points of this appendix. It turns out via Lipschitz property of $f$ and $h$, we are able to  upper bound $\gamma$ if the CP method is applied to $\Q_{\local}(c):=\Q\cap cr_{\local}\ball_{d+2}( q)$ for any $c>1$. In fact, the upcoming proof argument (see specifically \eqref{dolif:whyc>1}) would not go through for $c=1$, which is why we choose to work with the local set $\Q_{\local}:=\Q_{\local}(2)=\Q\cap 2r_\local \ball_{d+2}( q)$ in Theorem~\ref{dolif:theo:separation}. 







\begin{lemma}
\label{dolif:lem:localset}
Assume condition (B2) on $f$ and $h$. Fix any $c>1$ and define
\[
\ball_{c,\local}=c r_\local \ball_{d+2}( q),
\qquad
\Q_{\local}(c)=\Q\cap \ball_{c,\local},
\]
where $ q$ and $r_\local$ are defined at \eqref{dolif:def:localball}. Recall the notations $\widetilde{C}_h=\sqrt{L_h^2+a^2}, \widetilde{C}_f=\sqrt{L_f^2+a^2+b^2}$ at~\eqref{dolif:def:notation} and set
\[
\mu=\min\left\{\frac{a}{\widetilde{C}_h},\ \frac{b}{\widetilde{C}_f}\right\},\quad \lambda=2+\frac{a}{b},\quad \kappa=\frac{c-1}{\lambda+\mu}. 
\]
Moreover, define
\[
\bar p
=
 p_*+\left(0,\ \Delta s,\ \frac{a}{b}\Delta s+\Delta t\right),
\qquad
\Delta s=\Delta t=\kappa r_\local.
\]
Then we have the inclusion $B\left(\bar p,\ \kappa\mu r_\local\right)\subseteq \Q_{\local}(c)$. In particular,
\[
\mbox{\em minwidth}(\Q_{\local}(c))\ge 2\kappa\mu r_\local
\quad\text{and}\quad
\gamma(c):=\frac{c r_\local}{\mbox{\em minwidth}(\Q_{\local}(c))}
\le \frac{c(\lambda+\mu)}{(c-1)\mu}. 
\]
\end{lemma}


\begin{proof}  
Let us provide a sketch of the three steps of the proof.
To control $\gamma(c)=\frac{c r_{\local}}{\mbox{\em minwidth}(\Q_{\local}(c))}$, it suffices to lower bound
$\mbox{\em minwidth}(\Q_{\local}(c))$, which we do by constructing an inscribed ball inside
$\Q_{\local}(c)$.

Step~1 constructs a nontrivial ball $\ball_{\mathrm{in}}$ contained in $\Q$.
A subtlety is that Lemma~\ref{dolif:lem:minimizerinlocalball} only guarantees $ p_*\in \Q\cap r_{\local} \ball_{d+2}( q)$
and gives no quantitative slack for the constraints at $ p_*$; in particular, $ p_*$ may lie on or arbitrarily close to
$\partial\Q$, so a ball centered at $ p_*$ and inside $\Q$ need not exist.
We therefore shift $ p_*$ to
$\bar p= p_*+\left(0,\Delta s,\frac{a}{b}\Delta s+\Delta t\right)$ with $\Delta s=\Delta t=\kappa r_{\local}$,
which creates slack for both constraints: increasing $s$ relaxes $h(x)\le as$, while increasing $t$ by
$\frac{a}{b}\Delta s+\Delta t$ restores and adds slack for $f(x)+as\le bt$.
The Lipschitz bounds on $f$ and $h$ then imply that all points sufficiently close to $\bar p$ remain feasible,
yielding $\kappa\mu r_{\local} \ball_{d+2}(\bar p)\subseteq \Q$. Thus, $\kappa\mu r_{\local} \ball(\bar p)$ is the inscribed ball $\ball_{\mathrm{in}}$ in $\Q$ that we need. Note that this construction is nontrivial only when $c>1$, since then $\kappa=\frac{c-1}{\lambda+\mu}>0$.

Step~2 shows that $\ball_{\mathrm{in}}=\kappa\mu r_{\local}\ball_{d+2}(\bar p)$ also lies in $c r_{\local} \ball_{d+2}( q)$ and hence in $\Q_{\local}(c)$ (by combining with Step~1).
Using $\| p_*- q\|\le r_{\local}$ and
$\|\bar p- p_*\|=\Delta s+\left|\frac{a}{b}\Delta s+\Delta t\right|=\kappa\lambda r_{\local}$,
the triangle inequality gives
$\|\bar p- q\|+\kappa\mu r_{\local}\le r_{\local}+\kappa(\lambda+\mu)r_{\local}=c r_{\local}$. The last estimate implies $\kappa\mu r_{\local}\ball_{d+2}(\bar p)\subseteq c r_{\local}\ball_{d+2}( q)$. 

Step~3 uses the inclusion $\ball_{\mathrm{in}}=\kappa\mu r_{\local} \ball_{d+2}(\bar p)\subseteq \Q_{\local}(c)$ to lower bound $\mbox{\em minwidth}(\Q_{\local}(c))$ and finally upper bound $\gamma(c)$. 





\smallskip
\underline{Step 1:} Showing the nontrivial ball $\ball_{\mathrm{in}}:=\kappa\mu r_\local \ball(\bar p)\subseteq \Q$ for $c>1$.

\smallskip
Set
\[
g_1( p)=g_1(x,s,t)=h(x)-a s,
\qquad
g_2( p)=g_2(x,s,t)=f(x)+a s-b t,
\]
so that 
\begin{align}
\label{dolif:Qrewrite}
    \Q=\{ p:\ g_1( p)\le 0,\ g_2( p)\le 0\}. 
\end{align}
Because $h$ is $L_h$-Lipschitz and $f$ is $L_f$-Lipschitz, for all $ p=(x,s,t)$ and $ p'=(x',s',t')$,
\begin{align}
\label{dolif:consequencelipschitz}
|g_1( p)-g_1( p')|
\le L_h\|x-x'\|+a|s-s'|
\le \sqrt{L_h^2+a^2} \| p- p'\|
= \widetilde{C}_h\| p- p'\|,
\end{align}
and similarly $|g_2( p)-g_2( p')|\le \widetilde{C}_f\| p- p'\|$.

Since $ p_*\in \Q$, we have $g_1( p_*)\le 0$ and $g_2( p_*)\le 0$.
Define $\bar p$ as in the statement of the lemma then
\begin{align}
\label{dolif:compareg1g2}
g_1(\bar p)=g_1( p_*)-a\Delta s\le -a\Delta s,
\qquad
g_2(\bar p)=g_2( p_*)-b\Delta t\le -b\Delta t.
\end{align}
Moreover, we have 
\begin{align}
\label{dolif:compareradiusincribedball}
\kappa\mu r_\local=\min\left\{\frac{a\Delta s}{\widetilde{C}_h},\ \frac{b\Delta t}{\widetilde{C}_f}\right\}.
\end{align}
Then for any $\zeta$ with $\|\zeta\|\le \kappa\mu r_\local$, by \eqref{dolif:consequencelipschitz}, \eqref{dolif:compareg1g2}, \eqref{dolif:compareradiusincribedball}, we can write
\begin{align*}
g_1(\bar p+\zeta)&\stackrel{\eqref{dolif:consequencelipschitz}}{\le} g_1(\bar p)+\widetilde{C}_h\|\zeta\|\stackrel{\eqref{dolif:compareg1g2}}{\le} (-a\Delta s)+\widetilde{C}_h \kappa\mu r_\local\stackrel{\eqref{dolif:compareradiusincribedball}}{\le} 0,\\
g_2(\bar p+\zeta)&\le (-b\Delta t)+\widetilde{C}_f \kappa\mu r_\local\le 0. 
\end{align*}
Thus, in view of \eqref{dolif:Qrewrite}, we conclude $\kappa\mu r_\local \ball_{d+2}(\bar p)\subseteq \Q$. Notice in particular that the inscribed ball $\kappa\mu r_\local \ball_{d+2}(\bar p)$ is non-trivial only if $\kappa=\frac{c-1}{\lambda+\mu}>0$, and for that we require
\begin{align}
\label{dolif:whyc>1}
    c>1. 
    \end{align}

\smallskip
\underline{Step 2:} Showing $\ball_{\mathrm{in}}=\kappa\mu r_\local \ball_{d+2}(\bar p)\subseteq \Q_\local(c)$.

\smallskip

Next, since $\| p_*- q\|\le \| p_{k-1}- q\|=r_\local$, we have
\begin{align*}
\|\bar p- q\|
\le \| p_*- q\|+\|\bar p- p_*\|
&\le r_\local+\left(\Delta s+\left|\frac{a}{b}\Delta s+\Delta t\right|\right)\\
&= r_\local+\left(\left(1+\frac{a}{b}\right)\Delta s+\Delta t\right)
= r_\local+\kappa \lambda r_\local,
\end{align*}
which combined with $\kappa=\frac{c-1}{\lambda+\mu}$ leads to
\begin{align}
\label{dolif:inequalityconnectcenters}
\|\bar p- q\|+\kappa\mu r_\local
\le r_\local+\kappa(\lambda+\mu)r_\local
= r_\local+(c-1)r_\local
= c r_\local. 
\end{align}
Take any point $p\in \kappa\mu r_\local \ball_{d+2}(\bar p)$, then we get $\norm{p- q}\leq \norm{p-\bar p}+\norm{ q-\bar p}\leq  \kappa\mu r_\local+\|q-\bar p\|\stackrel{\eqref{dolif:inequalityconnectcenters}}{\leq} c r_\local$.
This implies $\kappa\mu r_\local \ball_{d+2}(\bar p)\subseteq c r_\local \ball_{d+2}( q)=\ball_{c,\local}$. Combining with Step~1 to get
$\kappa\mu r_\local \ball_{d+2}(\bar p)\subseteq \Q\cap \ball_{c,\local}=\Q_\local(c)$. 

\smallskip
\underline{Step 3:} Finally, since $\ball_{\mathrm{in}}=\kappa\mu r_\local \ball_{d+2}(\bar p)\subseteq \Q_\local(c)$ per Step 2, we have $\mbox{minwidth}(\Q_\local(c))\ge 2\kappa\mu r_\local$,
so that
\[
\gamma(c)=\frac{c r_\local}{\mbox{\em minwidth}(\Q_\local(c))}
\le \frac{c r_\local}{2\kappa\mu r_\local}
= \frac{c(\lambda+\mu)}{2(c-1)\mu}
\le \frac{c(\lambda+\mu)}{(c-1)\mu}.
\]
\end{proof}



%%%%%%%%%%%%%%%%%%%%%%%%%%%%%%%%%%%%%%%%%%%%%%%%%%%%%%%%%%%%

In the upcoming result, via Theorem~\ref{theo:jiang}, we describe the oracle complexities of solving the minimization problem $\argmin_{p\in \Q_{\local}} \widetilde{\Theta}^{\eta,\Q}_{y,u,v}(p)$. We then invoke Lemma~\ref{dolif:lem:localset} to bound the parameters appearing in these complexity estimates.


\begin{theorem}
\label{dolif:theo:cuttingplane}
Assume conditions~(B1),~(B2) and recall
\[
   p_*=\underset{p\in \Q_\local}\argmin  \,\widetilde{\Theta}^{\eta,\Q}_{y,u,v}( p).
\]
and the local radius $r_\local$ at \eqref{dolif:def:localball}. Let 
\begin{align}
\label{dolif:def:deltaandalpha}
   \delta = \sqrt{\frac{2\eta}{d+2}}, \qquad \alpha=\min\left\{\frac{\delta^2}{4r_{\local}^2},\frac{1}{2}\right\}.
\end{align}

Then the CP method by \cite{jiang2020cuttingplane} makes
${\cal O} \left((d+2) \log \frac{(d+2) \gamma}{\alpha}\right)$ calls to the separation
oracle of $\Q$ in order to generate a $1/(d+2)$-solution
$\tilde p\in \Q_{\local}$ to the optimization problem
$\min_{ p\in \Q_{\local}} \widetilde{\Theta}^{\eta,\Q}_{y,u,v}( p)$. Here $\gamma=\frac{R}{\mbox{minwidth}(\Q_{\local})}$
with $R$ any radius such that $\Q_{\local}\subseteq R \ball_{d+2}( q)$.

Moreover, under the additional assumption $a=b=d$, $\eta=1/d^2$ and $R=2r_\local$,  we have 
\begin{enumerate}[label=(\alph*)]
    \item the upper bound
    \begin{align}
    \label{dolif:conc:gamma}
    \gamma\ \le\ \frac{2(3+\mu)}{\mu} \quad\text{for}\quad \mu=\min\left\{\frac{d}{\sqrt{d^2+L_h^2}},\ \frac{d}{\sqrt{2d^2+L_f^2}}\right\};
\end{align}
\item the concentration inequality
\begin{align}
\label{dolif:conc:alpha}
    \Pr \left(
        \alpha\ \le\
        \min\left\{
        \frac{1}{2(d+2) (\sqrt{d+2}+d/2+1)^2},\ \frac{1}{2}
        \right\}
    \right)
    \le 2\exp \left(-\frac{d^2}{8}\right).
\end{align}
\end{enumerate}
\end{theorem}

\begin{proof}  
As noted at the start of the proof of Lemma~\ref{dolif:lem:compareP1P2}, Condition~(B1) together with $a,b>0$ implies that the set $\Q$ in~\eqref{dolif:def:piandQ} is convex. Now let $p^*$ and $p_*$ be respectively the maximizer and minimizer over $\Q_\local$ of  the objective $\widetilde{\Theta}^{\eta,\Q}_{y,u,v}$ at~\eqref{dolif:def:Theta}. We can write
\begin{align*}
    \widetilde{\Theta}^{\eta,\Q}_{y,u,v}( p^*)-\widetilde{\Theta}^{\eta,\Q}_{y,u,v}( p_*)
    &=
    \frac{1}{2\eta}\left(\| p^*- q\|^2-\| p_*- q\|^2\right)
    \le \frac{1}{2\eta}\| p^*- q\|^2.
\end{align*}



Combining with the fact that 
$\| p- q\|\le 2r_{\local}$ for all $ p\in \Q_{\local}$ to get
\begin{align}
\label{dolif:range:Theta}
    \widetilde{\Theta}^{\eta,\Q}_{y,u,v}( p^*)-\widetilde{\Theta}^{\eta,\Q}_{y,u,v}( p_*)
    \le \frac{(2r_{\local})^2}{2\eta}
    =\frac{2r_{\local}^2}{\eta}.
\end{align}
In view of convexity of $\Q$, \eqref{dolif:range:Theta}, the inequality \eqref{ineq:opt} from
Theorem~\ref{theo:jiang}, and the choice of $\alpha\in(0,1)$ at
\eqref{dolif:def:deltaandalpha}, Theorem~\ref{theo:jiang} guarantees the CP method by~\cite{jiang2020cuttingplane} produces a $1/(d+2)$ solution $\tilde{p}$ to the optimization problem
$\min_{ p\in \Q_{\local}} \widetilde{\Theta}^{\eta,\Q}_{y,u,v}( p)$, i.e.,
\begin{align*}
    \widetilde{\Theta}^{\eta,\Q}_{y,u,v}(\tilde p)-\min_{ p\in \Q_{\local}}\widetilde{\Theta}^{\eta,\Q}_{y,u,v}( p)
    &\stackrel{\eqref{ineq:opt}}{\le}
    \alpha\left(\widetilde{\Theta}^{\eta,\Q}_{y,u,v}( p^*)-\widetilde{\Theta}^{\eta,\Q}_{y,u,v}( p_*)\right)
    \stackrel{\eqref{dolif:def:deltaandalpha},\eqref{dolif:range:Theta}}{\le}
    \frac{\delta^2}{2\eta}\stackrel{\eqref{dolif:def:deltaandalpha}}{=}\frac{1}{d+2}. 
\end{align*}
In particular, the CP method by~\cite{jiang2020cuttingplane} requires a subgradient oracle for the objective $\widetilde{\Theta}^{\eta,\Q}_{y,u,v}$, which is available explicitly; and a separation oracle for $\Q_{\local}$. The latter is obtained by combining a free separation oracle for $2r_{\local}\ball_{d+2}(q)$ with a separation oracle for $\Q$. Hence, per Theorem~\ref{theo:jiang}, we need ${\cal O} \left((d+2) \log \frac{(d+2) \gamma}{\alpha}\right)$ calls to the separation
oracle of $\Q_{\local}$, which translate to ${\cal O} \left((d+2) \log \frac{(d+2) \gamma}{\alpha}\right)$ calls to the separation
oracle of $\Q$. 

Next, we will set $a=b=d$, $\eta=1/d^2$. The upper bound on $\gamma$ is derived in Lemma~\ref{dolif:lem:localset} with $c=2$ and this result requires condition~(B1). What remains is to derive the
concentration inequalities for $r_{\local}$ and $\alpha$.
Write $Z\sim \mathcal{N}(0,I_{d+2})$. Since $ q=(y,u,v-b\eta)= p_{k-1}+\sqrt{\eta} Z-b\eta e_t$ for $e_t=\brac{0,\ldots,0,1}$, we have $ r_{\local}=\| p_{k-1}- q\|=\|\sqrt{\eta} Z-b\eta e_t\|
  \le \sqrt{\eta} \|Z\|+b\eta$
and hence
\begin{align}
\label{relation_rlocZ}
  \frac{r_{\local}}{\sqrt{\eta}}-b\sqrt{\eta}\le \|Z\|.
\end{align}
The Gaussian concentration inequality from \cite[Equation (3.5)]{ledoux2013probability} says $ \Pr\{\|Z\|\ge \sqrt{d+2}+\ell\}\le 2\exp \left(-\frac{\ell^2}{2}\right)$, so combined with~\eqref{relation_rlocZ}, we get $ \Pr \left\{
    r_{\local}\ge (\sqrt{d+2}+\ell+b\sqrt{\eta})\sqrt{\eta}
  \right\}
  \le 2\exp \left(-\frac{\ell^2}{2}\right)$.  Plugging in $\eta=1/d^2$, $a=b=d$, and $\ell=d/2$ gives
\begin{align}
\label{dolif:conc:rloc}
    \Pr \left(
        r_{\local}\ \ge\ \frac{\sqrt{d+2}+d/2+1}{d}
    \right)
    \le 2\exp \left(-\frac{d^2}{8}\right).
\end{align}

Finally, on the event $ r_{\local}\le \frac{\sqrt{d+2}+d/2+1}{d}$, we have
\begin{align*}
  \alpha=\min\left\{\frac{\delta^2}{4r_{\local}^2},\frac{1}{2}\right\}
  &\ge
  \min\left\{
  \frac{\delta^2}{4\left(\frac{\sqrt{d+2}+d/2+1}{d}\right)^2},\ \frac{1}{2}
  \right\}\\
  &=\min\left\{\frac{1}{2(d+2)(\sqrt{d+2}+d/2+1)^2},\frac12\right\}, 
\end{align*}
and therefore
\[
  \left\{\alpha\le \min\left\{
  \frac{1}{2(d+2)(\sqrt{d+2}+d/2+1)^2},\frac12\right\}\right\}
  \subseteq
  \left\{r_{\local}> \frac{\sqrt{d+2}+d/2+1}{d}\right\}.
\]
Combining this inclusion with \eqref{dolif:conc:rloc} yields \eqref{dolif:conc:alpha}.
\end{proof}


%%%%%%%%%%%%%%%%%%%%%%%%%%%%%%%%%%%%%%%%%%%%%%%%%%%%%%%%%%%%%%%%%%%%%%%%%%%%%%%%%%%%%%%%%%%%%%%%%%%%%%%%%%%%%%%%%%%%%%%%%%%%%%%%%%%%%%%%%%%%%%%%%%%%%%%%%%%%%%%%%%%%%%%%%%

\subsection{Sampling $ p\sim \exp\brac{-\widetilde{\mathcal{P}}_1( p)}$ in Algorithm~\ref{dolif:RGO}}
\label{dolif:appen:specialalgo}

This part is similar to Appendix~\ref{silif:appen:specialalgo}. Recall $\widetilde{\mathcal{P}}_1$ is defined at~\eqref{dolif:def:tildeP1} and $\tilde p=(\tilde x,\tilde s,\tilde t)$ is defined in Algorithm~\ref{dolif:RGO}.
By completing the square, one can easily see that sampling $w\sim \exp\brac{-\widetilde{\mathcal{P}}_1(p)}$ is equivalent to
\begin{align*}
  p\sim \Lambda(p) \propto
 \exp\brac{-\frac{1}{2\eta}\brac{\norm{ p-\tilde p}-\sqrt{\frac{2\eta}{d+2}} }^2}.
\end{align*}

While $\Lambda( p)$ is not a Gaussian density, generating $p\sim \Lambda(p)$ is straightforward since it can be turned into a one-dimensional sampling problem. We state here a generic procedure for this sampling problem. 
\begin{algorithm}[H]
	\caption{Sample $ p\sim  \Lambda( p)$ } 
	\label{dolif:alg:samplingP1}
	\begin{algorithmic}
		\State 1. Generate $W\sim {\cal N}(0,I_{d+2})$  and set $\theta = W/\|W\|$; 
		\State 2. Generate $r \propto r^{d+1} \exp\left(-\frac{1}{2\eta} \left(r - \sqrt{\frac{2\eta}{d+2}}\right)^2\right)$ by Adaptive Rejection Sampling for one-dimensional log-concave distribution by \cite{gilks1992adaptive}. 
		\State 3. Output $ p = \tilde p + r \theta$.
	\end{algorithmic}
\end{algorithm}


    % \printbibliography



\end{document}
